% Lesson 52 — Speaking: Exchanges information on using ICTs to enhance work or learning — Anglais Tle A — Cours signé OnBuch+
% Programme officiel MINESEC. Compiler avec : tectonic EN52-speaking-exchanges-information-on-using-icts-to-enhance.tex
\newcommand{\DOCMATIERE}{Anglais}
\newcommand{\DOCNIVEAU}{Tle A}
\newcommand{\DOCMODULE}{Chapter 5 : Media and Communication}
\newcommand{\DOCLECON}{EN52}
\newcommand{\DOCTITRE}{Lesson 52 — Speaking: Exchanges information on using ICTs to enhance work or learning}
% OnBuch+ — préambule « cours signés » ENRICHI (compilé avec tectonic / XeLaTeX)
% Système visuel « Pop » d'OnBuch+ : fond crème (jamais blanc), bordures encre
% épaisses, ombres « dures » décalées, titres Archivo Black en capitales.
% Version MATHÉMATIQUES : graphes (pgfplots/tikz), boîtes pédagogiques
% (définition, propriété, méthode, attention, le savais-tu, exercice résolu,
% exercice, TP, bilan), page de garde et signature OnBuch+ ancrée en bas.
%
% Macros attendues dans main.tex AVANT \input{preamble} :
%   \DOCMATIERE  \DOCNIVEAU  \DOCMODULE  \DOCLECON (numéro)  \DOCTITRE
\documentclass[11pt]{article}

\usepackage{fontspec}
\usepackage[english,french]{babel} % français = langue principale ; l'anglais via \eng{..} / textebox
\usepackage[a4paper,margin=2cm,top=3.4cm,bottom=2.8cm,headheight=1.4cm,footskip=1.3cm]{geometry}
\usepackage{amsmath,amssymb,amsfonts}
\usepackage{mathtools} % \xmapsto, \coloneqq... souvent utilisés par les modèles
\usepackage{cancel} % simplifications barrées (bilans de forces)
% \abs{z} (module, valeur absolue) et \norm{u} (norme), utilisés par les modèles
\providecommand{\abs}{}\let\abs\relax\DeclarePairedDelimiter{\abs}{\lvert}{\rvert}
\providecommand{\norm}{}\let\norm\relax\DeclarePairedDelimiter{\norm}{\lVert}{\rVert}
\usepackage[table]{xcolor} % [table] : fournit aussi \rowcolors (alternance de lignes)
\usepackage{pagecolor}
\usepackage{tikz}
\usetikzlibrary{arrows.meta,positioning,calc,shapes.geometric,shapes.misc,shapes.arrows,decorations.pathmorphing,decorations.pathreplacing,decorations.markings,patterns,shadows,mindmap,trees,fit,backgrounds,matrix,intersections,angles,quotes}
\usepackage{pgfplots}
\usepackage[version=4]{mhchem} % équations nucléaires \ce{^{235}_{92}U}
\usepackage[european,siunitx]{circuitikz} % schémas électriques
\pgfplotsset{compat=1.18}
\usepgfplotslibrary{fillbetween,groupplots} % aires entre courbes (calcul intégral)
\usepackage{enumitem}
\usepackage{fancyhdr}
% [most] charge toutes les bibliothèques tcolorbox (listings, minted, raster,
% documentation...) et gonfle inutilement la mémoire TeX sur un document long
% (~300 boîtes) : on ne charge que ce qui est réellement utilisé.
\usepackage[skins,breakable]{tcolorbox}
\usepackage{graphicx}
\usepackage{colortbl}
\usepackage{booktabs}
\usepackage{tabularx}
\usepackage{diagbox} % case d'en-tête diagonale des tableaux à double entrée
% Colonnes extensibles centrée / gauche / droite (C, L, R), souvent utilisées par les modèles
\newcolumntype{C}{>{\centering\arraybackslash}X}
\newcolumntype{L}{>{\raggedright\arraybackslash}X}
\newcolumntype{R}{>{\raggedleft\arraybackslash}X}
\usepackage{array}
\usepackage{multicol}
\usepackage{multirow}
\usepackage{siunitx}
% parse-numbers=false : accepte aussi \SI{2,5 \times 10^{-3}}{...}
\sisetup{locale=FR,output-decimal-marker={,},group-minimum-digits=4,parse-numbers=false}
\DeclareSIUnit\ohm{\ensuremath{\symup{\Omega}}} % U+2126 absent de la police
\DeclareSIUnit\division{div} % réglages d'oscilloscope (ms/div, V/div)
\DeclareSIUnit\tr{tr} % tours (vitesse de rotation, ex. \SI{1500}{\tr\per\minute})
\DeclareSIUnit\molar{M} % concentration molaire (ex. \SI{3}{\molar}, \SI{1}{\micro\molar})
\DeclareSIUnit\an{an} % année (le modèle confond parfois avec \year, un compteur TeX) — ex. \SI{5}{\kilo\meter\per\an}
\DeclareSIUnit\FCFA{FCFA} % franc CFA (ex. \SI{500}{\FCFA\per\kilo\gram})
\DeclareSIUnit\degreeBrix{\textdegree Brix} % teneur en sucre d'un sirop/jus (réfractométrie)
\DeclareSIUnit\month{mois} % le modèle utilise parfois \month, un compteur TeX, comme unité de durée
\usepackage{qrcode}
% \degree : commande fréquemment utilisée par les modèles pour un angle ou une
% température en dehors de \SI{}{} (non fournie par les paquets chargés).
\providecommand{\degree}{^{\circ}}
% \qed (fin de démonstration), utilisé par les modèles sans amsthm
\providecommand{\qed}{\hfill$\square$}
% \barre : double barre des valeurs interdites dans un tableau de variations (tkz-tab)
\providecommand{\barre}{\ensuremath{\|}}
% environnement proof (amsthm non chargé) utilisé par les modèles
\ifdefined\proof\else\newenvironment{proof}[1][Démonstration]{\par\noindent\textit{#1.}\ }{\qed\par}\fi
\usepackage{lastpage}
\usepackage{hyperref}
\hypersetup{hidelinks}

% ============================================================
% Palette « Pop » OnBuch+ — mêmes hex que la classe P (plus_ui.dart)
% ============================================================
\definecolor{poporange}{HTML}{F59321}
\definecolor{popdark}{HTML}{E07A0C}
\definecolor{popcream}{HTML}{FFF6E8}
\definecolor{popink}{HTML}{1C1714}
\definecolor{poppurple}{HTML}{7A5AE0}
\definecolor{popgreen}{HTML}{2E9E6A}
\definecolor{poppink}{HTML}{E0457B}
\definecolor{popblue}{HTML}{2F7DE1}
\definecolor{popgold}{HTML}{C98A12}
\definecolor{popmuted}{HTML}{8A7F76}
\definecolor{popline}{HTML}{EADFD2}
\definecolor{popgray}{HTML}{8A7F76}
\colorlet{popgrey}{popgray}
% Alias tolérants (le modèle nomme parfois les couleurs différemment)
\colorlet{popwhite}{white}
\colorlet{popblack}{popink}
\colorlet{popred}{poppink}
\colorlet{popyellow}{popgold}
% Teintes claires (fonds de tableaux/encadrés) — le modèle nomme parfois
% les couleurs avec un suffixe L (ex. popblueL) sans les avoir définies.
\colorlet{poporangeL}{poporange!20}
\colorlet{popdarkL}{popdark!20}
\colorlet{poppurpleL}{poppurple!20}
\colorlet{popgreenL}{popgreen!20}
\colorlet{poppinkL}{poppink!20}
\colorlet{popblueL}{popblue!20}
\colorlet{popgoldL}{popgold!20}
\colorlet{popredL}{popred!20}
\colorlet{popgrayL}{popgray!20}
% Teintes claires pour fonds de tableaux / zones
\colorlet{popblueL}{popblue!10!white}
\colorlet{poporangeL}{poporange!14!white}
\colorlet{popgreenL}{popgreen!12!white}
\colorlet{poppurpleL}{poppurple!12!white}
\colorlet{poppinkL}{poppink!12!white}
\colorlet{popgoldL}{popgold!14!white}

\pagecolor{popcream}

% ============================================================
% Typographie
% ============================================================
\defaultfontfeatures{Ligatures=TeX}
\setmainfont{PlusJakartaSans-Medium.ttf}[Path=../../../fonts/,BoldFont=PlusJakartaSans-Bold.ttf]
% Symboles, API et emojis absents de la police principale : repli sur DejaVu Sans (caractères actifs)
\newfontface\symfont{DejaVuSans.ttf}[Path=../../../fonts/]
\makeatletter
\newcommand\symactive[1]{\catcode#1=13 \begingroup\lccode`\~=#1 \lowercase{\endgroup\def~}{{\symfont\char#1}}}
\newcommand\symblank[1]{\catcode#1=13 \begingroup\lccode`\~=#1 \lowercase{\endgroup\def~}{}}
\newcommand\symhyph[1]{\catcode#1=13 \begingroup\lccode`\~=#1 \lowercase{\endgroup\def~}{\mbox{-}}}
\newcount\symc
\newcommand\symrange[2]{\symc=#1 \@whilenum\symc<#2 \do{\expandafter\symactive\expandafter{\the\symc}\advance\symc by 1 }}
\newcommand\symblankrange[2]{\symc=#1 \@whilenum\symc<#2 \do{\expandafter\symblank\expandafter{\the\symc}\advance\symc by 1 }}
\makeatother
\symrange{"0250}{"02FF}\symactive{"2713}\symactive{"2714}\symactive{"2717}\symactive{"2718}\symactive{"2610}\symactive{"2611}\symactive{"2612}
\symhyph{"2011}\symhyph{"2E1A}\symblankrange{"1F300}{"1FAFF}\symblank{"FE0F}
% Maths en sans-serif (Fira Math) avec chiffres et lettres droites de Plus
% Jakarta, pour que les formules se fondent dans le texte.
\usepackage[math-style=ISO,bold-style=ISO]{unicode-math}
\setmathfont{FiraMath-Regular.otf}[Path=../../../fonts/]
\setmathfont{PlusJakartaSans-Medium.ttf}[Path=../../../fonts/,range={up/num,up/{Latin,latin},\mathup}]
\usepackage{icomma} % virgule décimale sans espace en mode math
% Caractères mathématiques Unicode absents des polices de texte (Plus Jakarta,
% Archivo Black) : rendus par leur équivalent mathématique, en texte comme en
% formule — sinon ils s'affichent en carré vide (ex. « Arithmétique dans ℤ »).
\usepackage{newunicodechar}
\newunicodechar{ℝ}{\ensuremath{\mathbb{R}}}
\newunicodechar{ℤ}{\ensuremath{\mathbb{Z}}}
\newunicodechar{ℂ}{\ensuremath{\mathbb{C}}}
\newunicodechar{ℕ}{\ensuremath{\mathbb{N}}}
\newunicodechar{ℚ}{\ensuremath{\mathbb{Q}}}
\newunicodechar{𝔻}{\ensuremath{\mathbb{D}}}
\newunicodechar{α}{\ensuremath{\alpha}}
\newunicodechar{β}{\ensuremath{\beta}}
\newunicodechar{γ}{\ensuremath{\gamma}}
\newunicodechar{δ}{\ensuremath{\delta}}
\newunicodechar{ε}{\ensuremath{\varepsilon}}
\newunicodechar{θ}{\ensuremath{\theta}}
\newunicodechar{λ}{\ensuremath{\lambda}}
\newunicodechar{μ}{\ensuremath{\mu}}
\newunicodechar{σ}{\ensuremath{\sigma}}
\newunicodechar{τ}{\ensuremath{\tau}}
\newunicodechar{φ}{\ensuremath{\varphi}}
\newunicodechar{ψ}{\ensuremath{\psi}}
\newunicodechar{ω}{\ensuremath{\omega}}
\newunicodechar{Σ}{\ensuremath{\Sigma}}
% majuscules produites par \MakeUppercase dans les titres (α-aminés -> Α)
\newunicodechar{Α}{\ensuremath{\alpha}}
\newunicodechar{Β}{\ensuremath{\beta}}
\newunicodechar{Γ}{\ensuremath{\Gamma}}
\newunicodechar{Δ}{\ensuremath{\Delta}}
\newunicodechar{Ε}{\ensuremath{\varepsilon}}
\newunicodechar{Θ}{\ensuremath{\Theta}}
\newunicodechar{Λ}{\ensuremath{\Lambda}}
\newunicodechar{Μ}{\ensuremath{\mu}}
\newunicodechar{Τ}{\ensuremath{\tau}}
\newunicodechar{Φ}{\ensuremath{\Phi}}
\newunicodechar{Ψ}{\ensuremath{\Psi}}
\newunicodechar{Ω}{\ensuremath{\Omega}}
\newunicodechar{↦}{\ensuremath{\mapsto}}
\newunicodechar{∈}{\ensuremath{\in}}
\newunicodechar{∉}{\ensuremath{\notin}}
\newunicodechar{∧}{\ensuremath{\wedge}}
\newunicodechar{⇔}{\ensuremath{\Leftrightarrow}}
\newunicodechar{⟺}{\ensuremath{\Longleftrightarrow}}
\newunicodechar{⇒}{\ensuremath{\Rightarrow}}
\newunicodechar{⟹}{\ensuremath{\Longrightarrow}}
\newunicodechar{∘}{\ensuremath{\circ}}
\newunicodechar{∪}{\ensuremath{\cup}}
\newunicodechar{∩}{\ensuremath{\cap}}
\newunicodechar{≡}{\ensuremath{\equiv}}
\newunicodechar{⊂}{\ensuremath{\subset}}
\newunicodechar{⊥}{\ensuremath{\perp}}
\newunicodechar{∃}{\ensuremath{\exists}}
\newunicodechar{∀}{\ensuremath{\forall}}
\newunicodechar{∛}{\ensuremath{\sqrt[3]{\,}}}
\newunicodechar{∜}{\ensuremath{\sqrt[4]{\,}}}
\newunicodechar{⃗}{\ensuremath{{}^{\to}}}
\newunicodechar{ⁿ}{\ifmmode^{n}\else\textsuperscript{\ensuremath{n}}\fi}
\newunicodechar{ˣ}{\ifmmode^{x}\else\textsuperscript{\ensuremath{x}}\fi}
\newunicodechar{ᵇ}{\ifmmode^{b}\else\textsuperscript{\ensuremath{b}}\fi}
\newunicodechar{ʳ}{\ifmmode^{r}\else\textsuperscript{\ensuremath{r}}\fi}
\newunicodechar{ᵘ}{\ifmmode^{u}\else\textsuperscript{\ensuremath{u}}\fi}
\newunicodechar{ʸ}{\ifmmode^{y}\else\textsuperscript{\ensuremath{y}}\fi}
\newunicodechar{ᵃ}{\ifmmode^{a}\else\textsuperscript{\ensuremath{a}}\fi}
\newunicodechar{ᵉ}{\ifmmode^{e}\else\textsuperscript{\ensuremath{e}}\fi}
\newunicodechar{ᵏ}{\ifmmode^{k}\else\textsuperscript{\ensuremath{k}}\fi}
\newunicodechar{ᵖ}{\ifmmode^{p}\else\textsuperscript{\ensuremath{p}}\fi}
\newunicodechar{ᴺ}{\ifmmode^{N}\else\textsuperscript{\ensuremath{N}}\fi}
\newunicodechar{ᶜ}{\ifmmode^{c}\else\textsuperscript{\ensuremath{c}}\fi}
\newunicodechar{ʰ}{\ifmmode^{h}\else\textsuperscript{\ensuremath{h}}\fi}
\newunicodechar{ᵠ}{\ifmmode^{q}\else\textsuperscript{\ensuremath{q}}\fi}
\newunicodechar{ₙ}{\ifmmode_{n}\else\textsubscript{\ensuremath{n}}\fi}
\newunicodechar{ₐ}{\ifmmode_{a}\else\textsubscript{\ensuremath{a}}\fi}
\newunicodechar{ᵢ}{\ifmmode_{i}\else\textsubscript{\ensuremath{i}}\fi}
\newunicodechar{ᵣ}{\ifmmode_{r}\else\textsubscript{\ensuremath{r}}\fi}
\newunicodechar{ₖ}{\ifmmode_{k}\else\textsubscript{\ensuremath{k}}\fi}
\newunicodechar{ᵦ}{\ifmmode_{\beta}\else\textsubscript{\ensuremath{\beta}}\fi}
\newunicodechar{ₓ}{\ifmmode_{x}\else\textsubscript{\ensuremath{x}}\fi}
\newfontfamily\popdisplay{ArchivoBlack-Regular.ttf}[Path=../../../fonts/]
\newfontfamily\poplogo{SpaceGrotesk-Bold.ttf}[Path=../../../fonts/]
\newfontfamily\popsemi{PlusJakartaSans-SemiBold.ttf}[Path=../../../fonts/]
\newfontfamily\popxbold{PlusJakartaSans-ExtraBold.ttf}[Path=../../../fonts/]
\newfontfamily\popxboldls{PlusJakartaSans-ExtraBold.ttf}[Path=../../../fonts/,LetterSpace=3.0]

\setlist[enumerate]{leftmargin=1.7em,itemsep=5pt,topsep=4pt}
\setlist[itemize]{leftmargin=1.7em,itemsep=4pt,topsep=4pt,label={\color{poporange}\textbullet}}
\setlength{\parindent}{0pt}
\setlength{\parskip}{5pt}
\renewcommand{\arraystretch}{1.25}

% pgfplots : style Pop par défaut
\pgfplotsset{
  every axis/.append style={axis line style={popink,line width=1pt},tick style={popink},
    grid style={popline,line width=0.5pt},label style={font=\small},tick label style={font=\footnotesize}},
  popcurve/.style={poporange,line width=1.8pt,no markers,smooth},
}
% Même style disponible dans un \draw TikZ simple (hors pgfplots)
\tikzset{popcurve/.style={poporange,line width=1.8pt}}
\tikzset{popgrid/.style={popline,very thin}}
\colorlet{popcurve}{poporange} % le nom du style est parfois utilisé comme couleur

% ============================================================
% Pill / squiggle
% ============================================================
\newcommand{\poppill}[3]{%
  \tikz[baseline=-0.55ex]\node[draw=popink,line width=1.2pt,rounded corners=2.6pt,%
    fill=#2,inner xsep=6.5pt,inner ysep=3pt]{%
    \popxboldls\fontsize{7}{7}\selectfont\color{#3}\MakeUppercase{#1}};%
}
\newcommand{\popsquiggle}[1]{%
  \tikz[baseline=-0.4ex]{
    \draw[#1,line width=2.6pt,line cap=round]
      plot [smooth] coordinates {(0,0.09) (0.34,0.01) (0.68,0.21) (1.02,0.01) (1.36,0.10)};
  }%
}
% Mot-clé surligné (marqueur orange sous le texte)
\newcommand{\cle}[1]{{\popxbold\color{popdark}#1}}

% ============================================================
% En-tête / pied
% ============================================================
\pagestyle{fancy}
\fancyhf{}
\renewcommand{\headrulewidth}{0pt}
\renewcommand{\footrulewidth}{0pt}
\lhead{\raisebox{-0.15ex}{\poplogo\fontsize{13}{13}\selectfont\color{popink}OnBuch\color{poporange}+}}
\rhead{\poppill{\DOCMATIERE\ \textbullet\ \DOCNIVEAU\ \textbullet\ Leçon \DOCLECON}{white}{popink}}
\lfoot{\popxbold\fontsize{7.6}{7.6}\selectfont\color{popmuted}\MakeUppercase{Cours signé OnBuch+}\ \textbullet\ {\popsemi\DOCTITRE}}
\rfoot{\tikz[baseline=-0.6ex]\node[rounded corners=8.5pt,fill=popink,minimum height=17pt,minimum width=17pt,inner xsep=5pt,inner sep=0]{\popxbold\fontsize{8}{8}\selectfont\color{popcream}\thepage\,/\,\pageref*{LastPage}};}

% ============================================================
% Titres
% ============================================================
\newcommand{\chaptitre}[2]{%
  \begin{tcolorbox}[enhanced,colback=poporange,colframe=popink,boxrule=2.6pt,arc=11pt,
      drop shadow=popink,left=15pt,right=15pt,top=12pt,bottom=11pt,before skip=3pt,after skip=18pt]
    \poppill{#2}{popink}{popcream}\par\vspace{9pt}
    {\raggedright\hyphenpenalty=10000\exhyphenpenalty=10000\popdisplay\fontsize{22}{24}\selectfont\color{popcream}\MakeUppercase{#1}\par}
  \end{tcolorbox}
}
% Section numérotée : pastille orange avec le numéro + titre Archivo
\newcounter{popsec}
\newcommand{\coursec}[1]{%
  \stepcounter{popsec}\setcounter{popsub}{0}%
  \par\vspace{16pt}\noindent
  \begin{minipage}{\linewidth}
  \tikz[baseline=-0.7ex]\node[draw=popink,line width=1.6pt,fill=poporange,rounded corners=4pt,minimum size=22pt,inner sep=0]{\popdisplay\fontsize{12}{12}\selectfont\color{popcream}\thepopsec};%
  \hspace{8pt}{\popdisplay\fontsize{15}{17}\selectfont\color{popink}\MakeUppercase{#1}}\par
  \vspace{4pt}\noindent{\color{poporange}\rule{36pt}{3.6pt}}
  \end{minipage}\par\nopagebreak\vspace{8pt}
}
\newcounter{popsub}
\newcommand{\courssub}[1]{%
  \stepcounter{popsub}%
  \par\vspace{9pt}\noindent{\popxbold\fontsize{11.5}{13}\selectfont\color{popink}{\color{poporange}\thepopsec.\thepopsub}\hspace{6pt}#1}\par\nopagebreak\vspace{4pt}
}

% ============================================================
% Boîtes pédagogiques — même recette : fond blanc, bordure épaisse colorée,
% ombre dure encre, badge plein. Titre optionnel [#1].
% ============================================================
\tcbset{popbase/.style={breakable,enhanced,colback=white,boxrule=2.3pt,arc=9pt,
  shadow={3pt}{-3pt}{0pt}{popink},left=14pt,right=14pt,top=11pt,bottom=11pt,before skip=11pt,after skip=15pt,
  fontupper=\popsemi\fontsize{10.6}{14.5}\selectfont\color{popink},
  fontlower=\popsemi\fontsize{10.6}{14.5}\selectfont\color{popink}}}
% Coche dessinée (le glyphe ✓ n'existe pas dans les polices du document)
\newcommand{\popcheck}[1]{\tikz[baseline=-0.5ex]\draw[#1,line width=1.6pt,line cap=round,line join=round](-0.13,0.0)--(-0.04,-0.09)--(0.13,0.1);}
\providecommand{\coche}{\popcheck{popgreen}}
\providecommand{\resultat}[1]{\par\smallskip\noindent\fcolorbox{popgreen}{popgreenL}{\parbox{\dimexpr\linewidth-2\fboxsep-2\fboxrule\relax}{\textbf{Résultat.} #1}}\par\smallskip}
\newcommand{\popboxhead}[3]{\poppill{#1}{#2}{white}\ifx\relax#3\relax\else\hspace{6pt}{\popxbold\fontsize{10.6}{12}\selectfont\color{popink}#3}\fi\par\vspace{7pt}}

\newtcolorbox{aretenir}[1][]{popbase,colframe=popgreen,before upper={\popboxhead{À retenir}{popgreen}{#1}}}
% Texte en anglais (document de lecture, dialogue, extrait) : cadre
% d'étude. Un titre optionnel (source, auteur) s'affiche dans l'en-tête.
\newtcolorbox{textebox}[1][]{popbase,colframe=popblue,before upper={\popboxhead{Text}{popblue}{#1}\begin{otherlanguage}{english}},after upper={\end{otherlanguage}}}
% Marques ✓ / ✗ (forme correcte / fautive) souvent utilisées par les modèles
\providecommand{\cross}{\textcolor{poppink}{\ensuremath{\times}}}
\providecommand{\xmark}{\cross}\providecommand{\crossmark}{\cross}\providecommand{\cmark}{\ensuremath{\checkmark}}
% Ligne de réponse / justification à compléter (inventée par les modèles)
\providecommand{\justification}[1]{\par\noindent\textit{Justification~:} \dotfill\par}
% \textipa (package tipa non chargé) : la police ne couvre pas l'API -> simple passage
\providecommand{\textipa}[1]{#1}
% Soulignement (ulem non chargé) : \uline, \ul
\providecommand{\uline}[1]{\underline{#1}}\providecommand{\ul}[1]{\underline{#1}}
% \glossaire{\item ...} inventée par les modèles : liste de glossaire
\providecommand{\glossaire}[1]{\begin{itemize}#1\end{itemize}}
% \alert (beamer) inventée par les modèles : texte mis en évidence
\providecommand{\alert}[1]{\textbf{\textcolor{poppink}{#1}}}
% variantes ASCII d'ordinaux écrites en macro
\providecommand{\iere}{\textsuperscript{re}}\providecommand{\ieme}{\textsuperscript{e}}\providecommand{\ere}{\textsuperscript{re}}\providecommand{\eme}{\textsuperscript{e}}\providecommand{\er}{\textsuperscript{er}}\providecommand{\ie}{\textsuperscript{e}}
% Ordinaux français écrits en macro par les modèles : 1\ière, 3\ième
\newcommand{\ière}{\textsuperscript{re}}\newcommand{\ième}{\textsuperscript{e}}
% \exemple (inventée par les modèles) : étiquette « Exemple : »
\providecommand{\exemple}{\textbf{Exemple~:}\ }
% \square utilisable aussi hors mode math (cases à cocher)
\AtBeginDocument{\let\squareMath\square\renewcommand{\square}{\ensuremath{\squareMath}}}
% Mot ou phrase en anglais à l'intérieur d'un paragraphe français
\newcommand{\eng}[1]{\ifmmode\text{\textit{#1}}\else\begingroup\selectlanguage{english}\itshape #1\endgroup\fi}
\let\esp\eng % alias toléré
% flèches d'intonation inventées par les modèles
\providecommand{\textsearrow}{}\renewcommand{\textsearrow}{\ensuremath{\searrow}}\providecommand{\textnearrow}{}\renewcommand{\textnearrow}{\ensuremath{\nearrow}}
% \fr{..} : mot français (inventé par les modèles)
\providecommand{\fr}[1]{\textit{#1}}
\newtcolorbox{exemplebox}[1][]{popbase,colframe=poporange,before upper={\popboxhead{Exemple}{poporange}{#1}}}
\newtcolorbox{definition}[1][]{popbase,colframe=popblue,before upper={\popboxhead{Définition}{popblue}{#1}}}
\newtcolorbox{propriete}[1][]{popbase,colframe=poppurple,before upper={\popboxhead{Propriété}{poppurple}{#1}}}
\newtcolorbox{methode}[1][]{popbase,colframe=popgold,before upper={\popboxhead{Méthode}{popgold}{#1}}}
\newtcolorbox{attention}[1][]{popbase,colframe=poppink,before upper={\popboxhead{Attention}{poppink}{#1}}}
\newtcolorbox{experience}[1][]{popbase,colframe=popblue,colback=popblueL,before upper={\popboxhead{Expérience}{popblue}{#1}}}
% « Le savais-tu ? » : carte violette pleine (contexte, culture, Cameroun)
\newtcolorbox{savaistu}[1][]{popbase,colback=poppurple,colframe=popink,
  fontupper=\popsemi\fontsize{10.4}{14.2}\selectfont\color{white},
  before upper={\poppill{Le savais-tu\,?}{white}{poppurple}\ifx\relax#1\relax\else\hspace{6pt}{\popxbold\color{white}#1}\fi\par\vspace{7pt}}}
% Exercice résolu : énoncé en haut, \tcblower puis la solution sur fond crème
\newcounter{popexo}\newcounter{popexr}
\newtcolorbox{exoresolu}[1][]{popbase,colframe=poporange,colbacklower=popcream,
  segmentation style={popink,line width=1.2pt,dashed},
  before upper={\stepcounter{popexr}\popboxhead{Exercice résolu \thepopexr}{poporange}{#1}},
  before lower={\poppill{Solution}{popink}{popcream}\par\vspace{6pt}}}
% Exercice (non résolu en place) — la correction est dans la partie Corrigés
\newtcolorbox{exercice}[1][]{popbase,colframe=popink,
  before upper={\stepcounter{popexo}\popboxhead{Exercice \thepopexo}{popink}{#1}}}
% Corrigé
\newtcolorbox{corrige}[1][]{popbase,colframe=popgreen,colback=popgreenL,
  before upper={\popboxhead{Corrigé}{popgreen}{#1}}}
% Objectifs / prérequis / bilan
\newtcolorbox{objectifs}{popbase,colframe=popink,colback=poporangeL,
  before upper={\popboxhead{Objectifs de la leçon}{popink}{}}}
\newtcolorbox{prerequis}{popbase,colframe=popink,colback=popblueL,
  before upper={\popboxhead{Prérequis}{popblue}{}}}
\newtcolorbox{bilan}{popbase,colframe=popink,colback=white,boxrule=2.8pt,
  before upper={\popboxhead{Fiche bilan}{poporange}{L'essentiel à connaître pour l'examen}}}
% Figure encadrée avec légende
\newtcolorbox{popfigure}[1][]{enhanced,breakable=false,colback=white,colframe=popline,boxrule=1.4pt,arc=8pt,
  left=10pt,right=10pt,top=10pt,bottom=8pt,before skip=10pt,after skip=12pt,halign=center,
  lower separated=false,halign lower=center,
  fontlower=\popsemi\fontsize{9}{11.5}\selectfont\color{popmuted},
  #1}
\newcommand{\legende}[1]{\par\vspace{6pt}{\popsemi\fontsize{9}{11.5}\selectfont\color{popmuted}#1}}

% ============================================================
% Signature OnBuch+
% ============================================================
\newcommand{\poplogoinline}[1]{{\poplogo\fontsize{#1}{#1}\selectfont\color{popink}OnBuch\color{poporange}+}}
\newcommand{\signatureonbuch}{%
  \begin{tcolorbox}[enhanced,colback=popink,colframe=popink,boxrule=2.2pt,arc=10pt,
      drop shadow=poporange,left=14pt,right=14pt,top=10pt,bottom=10pt,before skip=0pt,after skip=0pt]
    \tikz[baseline=-0.7ex]\node[circle,fill=popgreen,draw=popcream,line width=1.3pt,minimum size=22pt,inner sep=0]{\popcheck{white}};%
    \hspace{9pt}\begin{minipage}[c]{0.62\linewidth}
      {\popxbold\fontsize{12}{13}\selectfont\color{popcream}Signé\ \textbullet\ {\poplogo OnBuch\color{poporange}+}}\par\vspace{2pt}
      {\raggedright\popsemi\fontsize{8}{10.5}\selectfont\color{popline}\MakeUppercase{Équipe pédagogique OnBuch+}\\\MakeUppercase{Conforme au programme officiel MINESEC}\par}
    \end{minipage}\hfill
    {\poplogo\fontsize{22}{22}\selectfont\color{popcream}OnBuch\color{poporange}+}
  \end{tcolorbox}%
}

% ============================================================
% Page de garde
% #1 titre · #2 matière · #3 niveau · #4 module · #5 n° leçon · #6 durée ·
% #7 description / objectifs (texte ou itemize)
% ============================================================
\newcommand{\pagegarde}[7]{%
  \thispagestyle{empty}
  \newgeometry{a4paper,margin=1.7cm,top=1.6cm,bottom=1.4cm}%
  \begin{tikzpicture}[remember picture,overlay]
    \fill[poporange!18] ([xshift=-3.2cm,yshift=-3.6cm]current page.north east) circle (4.6cm);
    \fill[poppurple!14] ([xshift=2.4cm,yshift=7.6cm]current page.south west) circle (3.8cm);
  \end{tikzpicture}%
  {\poplogo\fontsize{30}{30}\selectfont\color{popink}OnBuch\color{poporange}+}\hfill
  \raisebox{0.6ex}{\poppill{Cours signé \textbullet\ Premium}{popink}{popcream}}\par
  \vspace{1pt}\popsquiggle{poppurple}\par
  \vspace{8pt}
  \begin{tcolorbox}[enhanced,colback=poporange,colframe=popink,boxrule=2.8pt,arc=13pt,
      drop shadow=popink,left=18pt,right=18pt,top=12pt,bottom=13pt,after skip=10pt]
    \poppill{#2 \textbullet\ #3}{popink}{popcream}\hspace{5pt}\poppill{#4}{white}{popink}\par\vspace{8pt}
    {\popdisplay\fontsize{14}{14}\selectfont\color{popink}LEÇON #5}\par\vspace{4pt}
    {\raggedright\hyphenpenalty=10000\exhyphenpenalty=10000\popdisplay\fontsize{25}{27}\selectfont\color{popcream}\MakeUppercase{#1}\par}
    \vspace{8pt}
    {\popxbold\fontsize{9.5}{9.5}\selectfont\color{popink}\MakeUppercase{Durée conseillée : #6}}
  \end{tcolorbox}
  \begin{tcolorbox}[enhanced,colback=white,colframe=popink,boxrule=2.2pt,arc=10pt,
      drop shadow=popink,left=16pt,right=16pt,top=10pt,bottom=10pt,after skip=8pt]
    \poppill{Dans cette leçon}{poppurple}{white}\par\vspace{5pt}
    \popsemi\fontsize{9.3}{12.6}\selectfont\color{popink}#7
  \end{tcolorbox}
  \noindent\begin{minipage}[t]{0.487\linewidth}
    \begin{tcolorbox}[enhanced,colback=popgreen,colframe=popink,boxrule=2.2pt,arc=10pt,
        drop shadow=popink,left=11pt,right=11pt,top=10pt,bottom=10pt,height=3.1cm,valign=center]
      \tcbox[colback=white,colframe=popink,boxrule=1.3pt,arc=5pt,boxsep=0pt,nobeforeafter,
        left=3pt,right=3pt,top=3pt,bottom=3pt,valign=center]{\qrcode[height=1.9cm]{https://play.google.com/store/apps/details?id=com.luvvix.onbuch&hl=fr}}%
      \hspace{8pt}\begin{minipage}[c]{0.5\linewidth}
        {\popdisplay\fontsize{11}{12.5}\selectfont\color{white}\MakeUppercase{Télécharge OnBuch}}\par\vspace{4pt}
        {\raggedright\popsemi\fontsize{8.4}{11}\selectfont\color{white}Gratuit \textbullet\ Google Play\\Tuteur IA, annales, concours, résultats\par}
      \end{minipage}
    \end{tcolorbox}
  \end{minipage}\hfill
  \begin{minipage}[t]{0.487\linewidth}
    \begin{tcolorbox}[enhanced,colback=poppurple,colframe=popink,boxrule=2.2pt,arc=10pt,
        drop shadow=popink,left=11pt,right=11pt,top=10pt,bottom=10pt,height=3.1cm,valign=center]
      \tikz[baseline=-0.7ex]\node[circle,fill=white,draw=popink,line width=1.3pt,minimum size=1.6cm,inner sep=0]{\popdisplay\fontsize{24}{24}\selectfont\color{poppurple}+};%
      \hspace{8pt}\begin{minipage}[c]{0.6\linewidth}
        {\popdisplay\fontsize{11}{12.5}\selectfont\color{white}\MakeUppercase{Passe à OnBuch+}}\par\vspace{4pt}
        {\raggedright\popsemi\fontsize{8.4}{11}\selectfont\color{white}Tous les cours signés, Tuteur IA illimité, fiches PDF — onglet OnBuch+ de l'app\par}
      \end{minipage}
    \end{tcolorbox}
  \end{minipage}
  \vspace{8pt}
  \popcontenu
  \vfill
  \signatureonbuch
  \restoregeometry
  \newpage
}

% Rangée « Ce cours contient » (page de garde)
\newcommand{\poptile}[3]{% #1 couleur #2 gros texte #3 légende
  \begin{tcolorbox}[enhanced,colback=#1,colframe=popink,boxrule=2pt,arc=9pt,shadow={2.5pt}{-2.5pt}{0pt}{popink},
      left=6pt,right=6pt,top=9pt,bottom=9pt,halign=center,valign=center,height=1.75cm,width=\linewidth,nobeforeafter]
    {\popdisplay\fontsize{12.5}{13}\selectfont\color{white}\MakeUppercase{#2}}\par\vspace{3pt}
    {\centering\popsemi\fontsize{7.8}{9.5}\selectfont\color{white}#3\par}
  \end{tcolorbox}}
\newcommand{\popcontenu}{%
  \par\noindent{\popxboldls\fontsize{7.5}{7.5}\selectfont\color{popmuted}CE COURS SIGNÉ CONTIENT}\par\vspace{7pt}
  \noindent\begin{minipage}[t]{0.235\linewidth}\poptile{popblue}{Cours}{complet, illustré, pas à pas}\end{minipage}\hfill
  \begin{minipage}[t]{0.235\linewidth}\poptile{poporange}{Méthodes}{et exercices résolus}\end{minipage}\hfill
  \begin{minipage}[t]{0.235\linewidth}\poptile{poppink}{Exercices}{type examen + corrigés}\end{minipage}\hfill
  \begin{minipage}[t]{0.235\linewidth}\poptile{popgreen}{Fiche bilan}{l'essentiel à retenir}\end{minipage}\par}

% Sommaire
\newtcolorbox{sommaire}{enhanced,colback=white,colframe=popink,boxrule=2.2pt,arc=10pt,
  drop shadow=popink,left=18pt,right=18pt,top=13pt,bottom=13pt,before skip=6pt,after skip=20pt,
  before upper={\poppill{Sommaire}{poppurple}{white}\par\vspace{9pt}\popsemi\fontsize{10.6}{14.5}\selectfont\color{popink}}}

% Signature de fin de cours — poussée tout en bas de la dernière page
\newcommand{\findecours}{%
  \par\vfill\nopagebreak
  \begin{center}\popsquiggle{poporange}\end{center}\vspace{4pt}
  \signatureonbuch
}
\AtBeginDocument{\def\llbracket{\mathopen{[\mkern-3.5mu[}}\def\rrbracket{\mathclose{]\mkern-3.5mu]}}} % intervalles d'entiers (glyphe ⟦ absent de la police)
% Glyphes absents de Fira Math : redéfinis à partir de glyphes disponibles
\AtBeginDocument{%
  \def\setminus{\mathbin{\backslash}}\let\smallsetminus\setminus
  \def\vdots{\mathord{\vbox{\baselineskip3.2pt\lineskiplimit0pt\kern4pt\hbox{.}\hbox{.}\hbox{.}}}}%
  \def\ddots{\mathinner{\mkern1mu\raise6.5pt\hbox{.}\mkern2mu\raise3.6pt\hbox{.}\mkern2mu\raise0.7pt\hbox{.}\mkern1mu}}%
  \def\triangle{\mathord{\scalebox{0.9}{$\symup{\Delta}$}}}%
  \def\nabla{\mathord{\raisebox{\depth}{\scalebox{1}[-1]{$\symup{\Delta}$}}}}%
  \def\checkmark{\popcheck{popdark}}%
  \def\star{\mathbin{\ast}}\def\bigstar{\mathord{\ast}}%
  \def\diamond{\mathbin{\scalebox{0.6}{$\Diamond$}}}%
  \def\bigcup{\mathop{\mathchoice{\raisebox{-0.3em}{\scalebox{1.7}{$\cup$}}}{\scalebox{1.25}{$\cup$}}{\cup}{\cup}}\displaylimits}%
  \def\bigcap{\mathop{\mathchoice{\raisebox{-0.3em}{\scalebox{1.7}{$\cap$}}}{\scalebox{1.25}{$\cap$}}{\cap}{\cap}}\displaylimits}%
  \def\lesseqqgtr{\mathrel{\lessgtr}}\def\bigoplus{\mathop{\oplus}\displaylimits}%
}
\providecommand{\ssi}{si et seulement si} % abréviation fréquente
\AtBeginDocument{\providecommand{\wideparen}{\overparen}} % arc de cercle

\begin{document}

\pagegarde{Lesson 52 — Speaking: Exchanges information on using ICTs to enhance work or learning}{Anglais}{Tle A}{Chapter 5 : Media and Communication}{EN52}{2 h}{Cette leçon d'expression orale apprend aux élèves de Terminale A à échanger des informations en anglais sur l'utilisation des TIC pour améliorer le travail et l'apprentissage. À travers des dialogues modèles, des formules fonctionnelles et des jeux de rôle guidés, les élèves développent leur aisance à poser des questions, répondre, relancer et conclure une conversation.
\vspace{4pt}
\begin{multicols}{2}\small
\begin{itemize}
\item À la fin de cette leçon, je sais utiliser le vocabulaire de base des TIC (devices, applications, Internet) dans une conversation.
\item À la fin de cette leçon, je sais demander une information sur l'utilisation des technologies (How do you...? Which app...?).
\item À la fin de cette leçon, je sais donner une information et expliquer l'utilité d'un outil numérique (I use... to...; It helps me...).
\item À la fin de cette leçon, je sais exprimer mon accord, mon désaccord et mon opinion sur les TIC.
\item À la fin de cette leçon, je sais relancer une conversation et demander des précisions (Really? Could you tell me more?).
\item À la fin de cette leçon, je sais conclure poliment un échange oral.
\end{itemize}\end{multicols}}

\begin{sommaire}
\begin{multicols}{2}
\begin{enumerate}
\item Warm-up: ICTs in our daily lives
\item Model dialogue 1: ICTs for learning
\item Functional language: asking, informing, follow-up, closing
\item Model dialogue 2: ICTs at work
\item Pronunciation, stress and intonation
\item Guided role-plays and free practice
\item Activité d'intégration
\item Méthodes et exercices résolus
\item Exercices
\item Corrigés des exercices
\item Fiche bilan
\end{enumerate}
\end{multicols}
\end{sommaire}

\begin{objectifs}
\begin{itemize}
\item À la fin de cette leçon, je sais utiliser le vocabulaire de base des TIC (devices, applications, Internet) dans une conversation.
\item À la fin de cette leçon, je sais demander une information sur l'utilisation des technologies (How do you...? Which app...?).
\item À la fin de cette leçon, je sais donner une information et expliquer l'utilité d'un outil numérique (I use... to...; It helps me...).
\item À la fin de cette leçon, je sais exprimer mon accord, mon désaccord et mon opinion sur les TIC.
\item À la fin de cette leçon, je sais relancer une conversation et demander des précisions (Really? Could you tell me more?).
\item À la fin de cette leçon, je sais conclure poliment un échange oral.
\item À la fin de cette leçon, je sais prononcer correctement le vocabulaire des TIC avec la bonne accentuation et la bonne intonation.
\item À la fin de cette leçon, je sais mener un court dialogue en binôme sur les TIC et l'apprentissage.
\end{itemize}
\end{objectifs}

\begin{prerequis}
\begin{itemize}
\item Vocabulaire des médias et de la communication (Chapter 5, leçons précédentes)
\item Questions en Wh- et en yes/no (révisées en Première)
\item Présent simple pour exprimer les habitudes (I use, I download)
\item Verbes modaux can/could pour la capacité et la politesse
\item Expressions d'opinion vues en Première (I think, in my opinion)
\end{itemize}
\end{prerequis}

\newpage

\coursec{Warm-up: ICTs in our daily lives}

\courssub{Lead-in: a connected day in Cameroon}

Imaginez ces trois scènes. Amina, élève de Terminale à Yaoundé, prépare son exposé d'histoire avec sa camarade de Buea : elles échangent des messages vocaux sur WhatsApp le soir et rédigent leur plan ensemble sur Google Docs. Pendant la récréation, deux élèves comparent leurs applications préférées pour réviser l'anglais : l'un jure par YouTube, l'autre par un dictionnaire en ligne. Au ministère, un jeune employé explique à son collègue comment la visioconférence a transformé ses réunions : plus besoin de se déplacer de Douala à Yaoundé pour une simple réunion d'une heure. Même un conducteur de bendskin à Bamenda utilise son smartphone pour recevoir des paiements par Mobile Money et localiser ses clients.

Partout au Cameroun, au Nigeria ou au Ghana, savoir \emph{parler en anglais} des technologies de l'information et de la communication (TIC) est devenu indispensable : à l'école, au travail, dans les examens et dans les concours. Cette leçon vous apprend à \cle{échanger des informations} sur l'utilisation des technologies pour \cle{travailler} et \cle{apprendre}.

\textbf{Problématique de la leçon :} \emph{How can we talk naturally, in English, about the digital tools we use every day for work and learning?}

\begin{methode}[Objectif de la séance]
À la fin de cette leçon, vous serez capable de :
\begin{enumerate}
\item nommer en anglais les outils numériques que vous utilisez chaque jour ;
\item poser des questions simples sur l'usage des TIC (\eng{Do you use a smartphone? What do you use it for?}) ;
\item répondre de façon adaptée et donner des précisions (fréquence, but, lieu) ;
\item utiliser les formules fonctionnelles pour demander, donner une information, relancer et conclure ;
\item mener un court dialogue entre un journaliste et un lycéen sur les TIC.
\end{enumerate}
\end{methode}

\courssub{Brainstorming: which digital tools do you use every day?}

Avant d'apprendre du vocabulaire nouveau, partons de \emph{votre} expérience. Répondez mentalement à ces trois questions d'amorce, en anglais :

\begin{itemize}
\item \eng{Do you use a smartphone?} — \eng{Yes, I do. / No, I don't.}
\item \eng{What do you use it for?} — \eng{I use it for chatting, watching videos and doing my homework.}
\item \eng{How often do you go online?} — \eng{I go online every day. / I go online twice a week.}
\end{itemize}

Remarquez la structure de la réponse : \eng{I use it for + nom ou verbe en -ing}. Cette tournure, très fréquente, sera réutilisée dans tous les dialogues de la leçon.

\begin{propriete}[La structure « use something for + -ing »]
Pour expliquer à quoi sert un outil, l'anglais utilise la structure :
\begin{center}
\eng{to use + objet + for + verbe en -ing}
\end{center}
\begin{itemize}
\item \eng{I use my phone for checking my emails.} « J'utilise mon téléphone pour consulter mes e-mails. »
\item \eng{We use WhatsApp for sharing documents.} « Nous utilisons WhatsApp pour partager des documents. »
\item \eng{She uses YouTube for learning English.} « Elle utilise YouTube pour apprendre l'anglais. »
\end{itemize}
Après \eng{for}, le verbe est toujours à la forme en \emph{-ing}, jamais à l'infinitif : on dit \eng{for learning}, et non \eng{for to learn}.
\end{propriete}

\courssub{Reading: a short text to observe the vocabulary}

Lisez le texte suivant. Soulignez tous les mots qui désignent des outils numériques ou des actions liées à Internet.

\begin{textebox}[ICTs in a student's life]
ICTs have changed the way we learn. With a simple smartphone, a student in Bamenda can watch a lesson on YouTube, check a word in an online dictionary and chat with classmates on WhatsApp. Before an exam, many students create revision groups on social media: they upload their notes, download past question papers and share useful links. At home, they connect to Wi-Fi or buy a data bundle to browse the Internet when the connection is slow.

ICTs are also changing the world of work. In offices in Douala and Yaoundé, employees send emails instead of letters, organise video conferences with partners abroad and use spreadsheets to manage budgets. A young secretary can attend a meeting in Buea without leaving her desk in Limbe.

However, technology is only a tool. A smartphone can help you learn English, but it can also waste your time if you spend hours watching funny videos. The key question is not “Do you use ICTs?” but “How do you use them?” Used wisely, a simple app can open the doors of knowledge. Used badly, the same app becomes a distraction. That is why teachers now encourage students to use e-learning platforms, educational videos and e-books, and to see their phones as small schools in their pockets.
\end{textebox}

\textbf{Glossaire :}

\begin{center}
\begin{tabularx}{\linewidth}{|l|l|X|}
\hline
\rowcolor{popblueL} \textbf{Mot anglais} & \textbf{Nature} & \textbf{Traduction} \\
\hline
to check a word & verbe & vérifier un mot \\
\hline
classmates & nom & camarades de classe \\
\hline
past question papers & nom & anciens sujets d'examen \\
\hline
abroad & adverbe & à l'étranger \\
\hline
to waste time & verbe & perdre son temps \\
\hline
wisely & adverbe & sagement, judicieusement \\
\hline
distraction & nom & distraction, diversion \\
\hline
spreadsheet & nom & tableur (logiciel de calcul) \\
\hline
\end{tabularx}
\end{center}

\textbf{Questions de compréhension :}
\begin{enumerate}
\item \eng{What can a student in Bamenda do with a smartphone?}
\item \eng{What do students share in their revision groups?}
\item \eng{How do office workers use ICTs, according to the text?}
\item \eng{What danger of the smartphone does the text mention?}
\end{enumerate}

\begin{corrige}[Questions de compréhension]
\begin{enumerate}
\item \eng{He or she can watch a lesson on YouTube, check a word in an online dictionary and chat with classmates on WhatsApp.}
\item \eng{They upload their notes, download past question papers and share useful links.}
\item \eng{They send emails, organise video conferences and use spreadsheets to manage budgets.}
\item \eng{A smartphone can waste your time and become a distraction if you use it badly.}
\end{enumerate}
\end{corrige}

\courssub{Key vocabulary: the words you need}

\begin{definition}[ICT]
\cle{ICT} est le sigle de \emph{Information and Communication Technology} (« technologies de l'information et de la communication »). Il désigne l'ensemble des outils numériques utilisés pour \emph{communiquer} et \emph{traiter l'information} : téléphones, ordinateurs, Internet, applications, réseaux sociaux. Au pluriel, on dit souvent \eng{ICTs} : \eng{ICTs are changing education in Africa.}
\end{definition}

Voici le vocabulaire clé de la leçon, à connaître activement :

\begin{center}
\begin{tabularx}{\linewidth}{|l|X|X|}
\hline
\rowcolor{popblueL} \textbf{English} & \textbf{Français} & \textbf{Exemple en contexte} \\
\hline
a device & un appareil (appareil électronique) & \eng{A smartphone is a useful device.} \\
\hline
an app / an application & une application & \eng{WhatsApp is my favourite app.} \\
\hline
to browse the Internet & naviguer sur Internet & \eng{I browse the Internet every evening.} \\
\hline
to download & télécharger (vers son appareil) & \eng{I downloaded an e-book yesterday.} \\
\hline
to upload & mettre en ligne, envoyer (vers Internet) & \eng{She uploaded her notes to the group.} \\
\hline
Wi-Fi & le Wi-Fi (la connexion sans fil) & \eng{Is there Wi-Fi in your school?} \\
\hline
a data bundle & un forfait de données & \eng{I bought a data bundle this morning.} \\
\hline
social media & les réseaux sociaux & \eng{Many students use social media to revise.} \\
\hline
an e-learning platform & une plateforme d'apprentissage en ligne & \eng{Our school has an e-learning platform.} \\
\hline
video conferencing & la visioconférence & \eng{We use video conferencing for meetings.} \\
\hline
a spreadsheet & un tableur & \eng{He uses a spreadsheet for the budget.} \\
\hline
\end{tabularx}
\end{center}

\begin{attention}[Piège : « information » est indénombrable]
En français, on dit « des informations ». En anglais, \eng{information} est \emph{indénombrable} : il ne prend \textbf{jamais} de \emph{-s} et ne s'emploie jamais avec \emph{a/an}.
\begin{itemize}
\item Correct : \eng{I need some information.} « J'ai besoin d'informations. »
\item Incorrect : \eng{I need informations.} / \eng{an information}
\end{itemize}
Pour compter, on dit \eng{a piece of information} (« une information ») ou \eng{two pieces of information}. Même piège avec \eng{news} (indénombrable) : \eng{The news is good.} « Les nouvelles sont bonnes. »
\end{attention}

\begin{attention}[download ou upload ?]
Les francophones confondent souvent les deux verbes. Retenez la direction de la flèche :
\begin{itemize}
\item \eng{to download} = recevoir \emph{depuis} Internet vers votre appareil : \eng{I downloaded a video.} « J'ai téléchargé une vidéo. »
\item \eng{to upload} = envoyer \emph{vers} Internet : \eng{I uploaded my photo.} « J'ai mis ma photo en ligne. »
\end{itemize}
Astuce : \emph{down} = vers le bas (vers vous), \emph{up} = vers le haut (vers le réseau).
\end{attention}

\courssub{Mind map: organising the vocabulary}

Pour mémoriser ce lexique, rien de tel qu'une carte mentale. Observez comment les mots se rangent en quatre familles autour du mot central \eng{ICTs}.

\begin{popfigure}
\begin{center}
\begin{tikzpicture}[mindmap, grow cyclic,
  every node/.style={concept, text=white, font=\small},
  root concept/.append style={concept color=popblue, font=\bfseries},
  level 1/.append style={level distance=3.2cm, sibling angle=90},
  level 2/.append style={level distance=2.2cm, sibling angle=45, font=\scriptsize}]
\node[root concept]{ICTs}
  child[concept color=poporange]{ node {Devices}
    child { node {phone} }
    child { node {laptop} }
    child { node {tablet} } }
  child[concept color=popgreen]{ node {Apps}
    child { node {WhatsApp} }
    child { node {Zoom} }
    child { node {YouTube} } }
  child[concept color=poppurple]{ node {Uses}
    child { node {learning} }
    child { node {working} }
    child { node {chatting} } }
  child[concept color=poppink]{ node {Places}
    child { node {school} }
    child { node {home} }
    child { node {office} } };
\end{tikzpicture}
\end{center}
\legende{Carte mentale du vocabulaire des TIC : quatre familles de mots à réutiliser dans vos dialogues.}
\end{popfigure}

\begin{methode}[Utiliser la carte mentale pour parler]
\begin{enumerate}
\item Choisissez un \emph{device} : \eng{my phone}, \eng{my laptop}.
\item Ajoutez une \emph{app} : \eng{WhatsApp}, \eng{Google}, \eng{YouTube}.
\item Précisez l'\emph{usage} avec \eng{for + -ing} : \eng{for learning English}, \eng{for sending emails}.
\item Précisez le \emph{lieu} et la \emph{fréquence} : \eng{at home}, \eng{every day}.
\end{enumerate}
Résultat : \eng{I use my phone and WhatsApp for learning English at home every day.} — une phrase complète, correcte et naturelle.
\end{methode}

\courssub{ICTs for work vs ICTs for learning}

Tous les outils numériques ne servent pas au même but. Distinguons deux grandes catégories que vous devrez comparer dans vos dialogues.

\begin{center}
\begin{tabularx}{\linewidth}{|X|X|}
\hline
\rowcolor{popblueL} \textbf{ICTs for work (travail)} & \textbf{ICTs for learning (études)} \\
\hline
emails (les e-mails professionnels) & online dictionaries (dictionnaires en ligne) \\
\hline
video conferencing (la visioconférence) & educational videos (vidéos éducatives) \\
\hline
spreadsheets (les tableurs) & e-books (les livres numériques) \\
\hline
printers and scanners (imprimantes et scanners) & e-learning platforms (plateformes en ligne) \\
\hline
\eng{I use emails to contact my boss.} & \eng{I use an online dictionary to check new words.} \\
\hline
\end{tabularx}
\end{center}

\begin{exemplebox}[Comparer les deux usages]
\begin{itemize}
\item \eng{At work, my father uses video conferencing for his meetings.} « Au travail, mon père utilise la visioconférence pour ses réunions. »
\item \eng{At school, I use YouTube to watch English lessons.} « À l'école, j'utilise YouTube pour regarder des leçons d'anglais. »
\item \eng{Some apps are useful for both work and learning.} « Certaines applications sont utiles à la fois pour le travail et les études. »
\end{itemize}
Remarquez : \eng{both ... and ...} signifie « à la fois ... et ... ».
\end{exemplebox}

\begin{exoresolu}[Classify the tools]
Énoncé : classez les outils suivants en deux colonnes, \eng{work} ou \eng{learning}, puis écrivez une phrase complète pour deux d'entre eux : \emph{emails, online dictionary, spreadsheet, educational videos, video conferencing, e-books}.
\tcblower
Solution détaillée :
\begin{itemize}
\item \textbf{Work :} emails, spreadsheet, video conferencing.
\item \textbf{Learning :} online dictionary, educational videos, e-books.
\end{itemize}
Phrases possibles :
\begin{itemize}
\item \eng{I use an online dictionary for checking difficult words.} « J'utilise un dictionnaire en ligne pour vérifier les mots difficiles. »
\item \eng{My uncle uses video conferencing for his meetings with partners in Lagos.} « Mon oncle utilise la visioconférence pour ses réunions avec des partenaires à Lagos. »
\end{itemize}
Méthode : demandez-vous « qui utilise cet outil, et pour quoi faire ? ». Si la réponse concerne un employé, un bureau, une réunion, c'est \emph{work} ; si elle concerne un élève, une leçon, un examen, c'est \emph{learning}.
\end{exoresolu}

\courssub{Functional language: asking, informing, following up, closing}

Pour mener un échange naturel, il faut maîtriser les formules de base de l'interaction. Voici les outils linguistiques pour structurer votre dialogue.

\begin{propriete}[Formules fonctionnelles pour l'échange d'informations]
\begin{center}
\begin{tabularx}{\linewidth}{|X|X|}
\hline
\rowcolor{popblueL} \textbf{Fonction} & \textbf{Expressions utiles} \\
\hline
\emph{Demander une information} & \eng{Do you use...?} / \eng{What do you use... for?} / \eng{How often do you...?} / \eng{Which app do you prefer?} / \eng{Can you tell me about...?} \\
\hline
\emph{Donner une information} & \eng{I use... for + -ing.} / \eng{I usually...} / \eng{My favourite app is...} / \eng{It helps me to...} \\
\hline
\emph{Relancer / Approfondir} & \eng{Really?} / \eng{That's interesting.} / \eng{Why do you like it?} / \eng{And what about...?} / \eng{Can you give me an example?} \\
\hline
\emph{Conclure / Remercier} & \eng{Thank you for your time.} / \eng{That's all for today.} / \eng{Thanks for sharing your experience.} / \eng{Have a nice day!} \\
\hline
\end{tabularx}
\end{center}
\end{propriete}

\courssub{Model dialogue 1: ICTs for learning}

Écoutez (ou lisez) ce dialogue entre un journaliste (\eng{Journalist}) et un élève de Terminale (\eng{Student}). Repérez les formules fonctionnelles et la structure \eng{use ... for + -ing}.

\begin{textebox}[Interview: ICTs for learning]
\textbf{Journalist:} Good morning. I'm doing a report on students and technology. \textbf{Do you use a smartphone?}

\textbf{Student:} \textbf{Yes, I do.} I use it every day.

\textbf{Journalist:} \textbf{What do you use it for?}

\textbf{Student:} \textbf{I use it for browsing the Internet} and \textbf{for chatting with my classmates on WhatsApp.} I also use it \textbf{for downloading past question papers.}

\textbf{Journalist:} \textbf{That's interesting.} \textbf{How often do you go online?}

\textbf{Student:} \textbf{I go online every evening} after school, \textbf{especially when I have Wi-Fi.}

\textbf{Journalist:} \textbf{Which app helps you most for your studies?}

\textbf{Student:} \textbf{YouTube, definitely.} I watch educational videos for English and Science. It's \textbf{very useful.}

\textbf{Journalist:} \textbf{Thank you for your time.} That was very helpful.

\textbf{Student:} \textbf{You're welcome. Have a nice day!}
\end{textebox}

\begin{methode}[Analyse du dialogue 1]
\begin{enumerate}
\item \textbf{Opening} : Le journaliste se présente et pose une question fermée (\eng{Do you use...?}).
\item \textbf{Details} : Questions ouvertes (\eng{What... for?}, \eng{How often...?}, \eng{Which...?}).
\item \textbf{Responses} : L'élève utilise le \emph{Present Simple} pour les habitudes (\eng{I use, I go}) et la structure \eng{use ... for + -ing}.
\item \textbf{Follow-up} : \eng{That's interesting.} montre l'écoute active.
\item \textbf{Closing} : Formules de politesse standardisées.
\end{enumerate}
\end{methode}

\courssub{Model dialogue 2: ICTs at work}

Ce dialogue met en scène un employé de bureau (\eng{Employee}) et un journaliste (\eng{Journalist}). Notez le vocabulaire spécifique au travail.

\begin{textebox}[Interview: ICTs at work]
\textbf{Journalist:} Good afternoon. Can you tell me how you use technology at work?

\textbf{Employee:} Certainly. \textbf{I use a laptop and a smartphone.} They are essential devices.

\textbf{Journalist:} \textbf{What do you use them for?}

\textbf{Employee:} \textbf{I use my laptop for sending emails} and \textbf{for working on spreadsheets.} I use my phone \textbf{for video conferencing} when I'm not at my desk.

\textbf{Journalist:} \textbf{Really? Do you often have video conferences?}

\textbf{Employee:} \textbf{Yes, twice a week.} We connect with our partners in Lagos and Nairobi. It saves time and money.

\textbf{Journalist:} \textbf{That's very efficient.} \textbf{Do you use social media for work?}

\textbf{Employee:} \textbf{Not really.} We use a professional platform for sharing documents. Social media is for my private life.

\textbf{Journalist:} \textbf{Thanks for sharing your experience.}

\textbf{Employee:} \textbf{My pleasure. Goodbye.}
\end{textebox}

\courssub{Pronunciation, stress and intonation}

La prononciation correcte des mots-clés et l'intonation des phrases sont essentielles pour être compris. Entraînez-vous à voix haute.

\begin{propriete}[Accentuation des mots clés (Word Stress)]
En anglais, l'accent tonique (la syllabe forte) change le sens ou la nature du mot. Marquez-le bien.
\begin{center}
\begin{tabularx}{\linewidth}{|l|l|X|}
\hline
\rowcolor{popblueL} \textbf{Mot} & \textbf{Prononciation (approx. fr.)} & \textbf{Règle / Remarque} \\
\hline
device & « di-\textbf{vaïss} » & Accent sur la 2\up{e} syllabe (verbe/nom) \\
\hline
download & « \textbf{daoun}-lode » & Accent sur la 1\up{re} (verbe particulaire) \\
\hline
upload & « \textbf{ap}-lode » & Accent sur la 1\up{re} (verbe particulaire) \\
\hline
Wi-Fi & « \textbf{ouaï}-\textbf{faï} » & Deux accents forts (sigle) \\
\hline
application & « a-pli-\textbf{ké}-chon » & Accent sur l'antépénultième \\
\hline
technology & « tek-\textbf{no}-lo-dji » & Accent sur \emph{no} ; \emph{ch} = /k/ \\
\hline
conference & « \textbf{kon}-fe-rens » & Accent sur la 1\up{re} \\
\hline
spreadsheet & « \textbf{spred}-chiit » & Mot composé : accent sur 1\up{er} élément \\
\hline
\end{tabularx}
\end{center}
\end{propriete}

\begin{propriete}[Intonation des phrases (Sentence Intonation)]
L'intonation (la mélodie de la phrase) signale le type de question.
\begin{itemize}
\item \textbf{Questions fermées (Yes/No) :} Intonation \textbf{montante} à la fin. \eng{Do you use a smartphone? ↗}
\item \textbf{Questions ouvertes (Wh-) :} Intonation \textbf{descendante} à la fin. \eng{What do you use it for? ↘}
\item \textbf{Listes :} Intonation montante sur les éléments non finaux, descendante sur le dernier. \eng{I use it for chatting ↗, watching videos ↗ and doing homework ↘.}
\item \textbf{Affirmations / Réponses :} Intonation descendante. \eng{I use it every day. ↘}
\end{itemize}
\end{propriete}

\begin{attention}[Piège de prononciation : le \emph{ch} anglais]
Le \emph{ch} anglais se prononce « tch » (\eng{chat} « tchat », \eng{church} « tcheurtch »), mais dans les mots d'origine grecque comme \eng{technology}, \eng{character}, \eng{chemistry}, \eng{mechanic}, il se prononce « k ». Ne dites donc pas « tèk-no-lo-tchi » mais « tèk-no-lo-dji ».
\end{attention}

\begin{experience}[Entraînement à la prononciation : \emph{Shadowing}]
En binôme : l'élève A lit une phrase du dialogue 1 ou 2 ; l'élève B la répète \emph{immédiatement} en imitant l'intonation et l'accentuation (technique du \emph{shadowing}). Changez de rôle. Concentrez-vous sur : \eng{What do you use it for? ↘}, \eng{That's interesting. ↘}, \eng{Really? ↗}.
\end{experience}

\courssub{Guided role-plays and free practice}

C'est le moment de passer à l'oral. Vous allez réaliser la tâche finale : un jeu de rôle d'interview.

\begin{experience}[Tâche finale : Jeu de rôle « Journaliste / Lycéen »]
\textbf{Consigne :} En binôme, préparez et jouez une interview de 2 minutes sur le thème : \emph{« How do you use ICTs for learning? »}

\textbf{Rôle A (Journaliste) :} Vous travaillez pour \eng{Cameroon Radio Television (CRTV) English Service}. Préparez 5 questions en utilisant les formules fonctionnelles (ouverte/fermée, relance).
\textbf{Rôle B (Lycéen) :} Vous êtes en Terminale A. Préparez vos réponses en utilisant le vocabulaire de la leçon (\eng{device, app, for + -ing, frequency, place}).

\textbf{Étapes de préparation (10 min) :}
\begin{enumerate}
\item \textbf{Individuel :} Complétez votre fiche de préparation.
  \begin{itemize}
  \item \textit{Journaliste :} \eng{Question 1: Do you...? Question 2: What... for? Question 3: How often...? Question 4: Which app...? Question 5: Do you think...?}
  \item \textit{Lycéen :} \eng{I use my [device] and [app] for [V-ing] at [place] [frequency]. My favourite app is... because...}
  \end{itemize}
\item \textbf{Binôme :} Répétez le dialogue à voix basse, corrigez la prononciation (accentuation, intonation Wh-/Yes-No).
\item \textbf{Public :} Présentez votre interview devant la classe. Les autres élèves notent : vocabulaire utilisé, structure \eng{for + -ing}, intonation, fluidité.
\end{enumerate}

\textbf{Grille d'auto-évaluation (à remplir après le passage) :}
\begin{center}
\begin{tabularx}{\linewidth}{|X|c|c|c|}
\hline
\rowcolor{popblueL} \textbf{Critère} & \textbf{Acquis} & \textbf{En cours} & \textbf{Non acquis} \\
\hline
J'utilise 5 mots de vocabulaire ICT minimum & & & \\
\hline
J'utilise correctement \eng{use ... for + -ing} & & & \\
\hline
Je pose/réponds aux questions (Wh- / Yes-No) & & & \\
\hline
Mon intonation est correcte (Wh- ↘, Yes-No ↗) & & & \\
\hline
Je utilise des formules de relance / conclusion & & & \\
\hline
\end{tabularx}
\end{center}
\end{experience}

\begin{exercice}[Entraînement écrit : Prepare your script]
Écrivez le script complet de votre interview (10-12 lignes) dans votre cahier. Utilisez au moins 6 mots du tableau de vocabulaire clé et 3 formules fonctionnelles différentes. Soulignez les mots de vocabulaire ICT et encadrez les formules fonctionnelles.
\end{exercice}

\begin{corrige}[Exercice : Exemple de script attendu]
\textbf{Journalist:} Good morning. \fbox{Can you tell me} how you use technology for school? \\
\textbf{Student:} \fbox{Certainly.} \textbf{I use my \underline{smartphone}} and \textbf{\underline{laptop}} every day. \\
\textbf{Journalist:} \fbox{What do you use them for?} \\
\textbf{Student:} \textbf{I use my phone for \underline{browsing the Internet} and for \underline{chatting} on WhatsApp.} I use my laptop for \textbf{\underline{downloading} \underline{past question papers}} and \textbf{\underline{writing} essays}. \\
\textbf{Journalist:} \fbox{That's interesting.} \fbox{How often do you go online?} \\
\textbf{Student:} \textbf{I go online every evening} when I have \textbf{\underline{Wi-Fi}} or a \textbf{\underline{data bundle}}. \\
\textbf{Journalist:} \fbox{Which app helps you most?} \\
\textbf{Student:} \textbf{YouTube.} I watch \textbf{\underline{educational videos}} for English. \fbox{It's very useful.} \\
\textbf{Journalist:} \fbox{Thank you for your time.} \\
\textbf{Student:} \fbox{You're welcome. Goodbye!}
\end{corrige}

\begin{savaistu}[WhatsApp et le Baccalauréat]
Au Cameroun, de nombreux élèves de Terminale créent des groupes WhatsApp de révision avant le Baccalauréat : ils y partagent des sujets d'annales, des fiches de grammaire et des enregistrements audio pour s'entraîner à la compréhension orale. Le téléphone portable, souvent critiqué, devient alors une véritable petite école de poche — à condition de l'utiliser avec discipline !
\end{savaistu}

\begin{aretenir}[Ce qu'il faut retenir de cette leçon]
\begin{itemize}
\item \eng{ICT} = \emph{Information and Communication Technology} : les outils numériques pour communiquer et traiter l'information.
\item Vocabulaire clé : \eng{device, app, to browse the Internet, to download, to upload, Wi-Fi, data bundle, social media, e-learning platform, video conferencing, spreadsheet}.
\item Structure d'usage : \eng{to use + objet + for + -ing} (jamais \eng{for to + verbe}).
\item \eng{information} est indénombrable : jamais de \emph{-s} (pas \eng{informations}), on dit \eng{a piece of information}.
\item Distinction : \eng{ICTs for work} (emails, video conferencing, spreadsheets) vs \eng{ICTs for learning} (online dictionaries, educational videos, e-books).
\item Formules fonctionnelles : Questions fermées (intonation ↗), questions ouvertes (intonation ↘), relance (\eng{Really?}, \eng{That's interesting.}), conclusion (\eng{Thank you for your time.}).
\item Prononciation : \eng{technology} (« tèk-no-lo-dji », \emph{ch}=k), \eng{download/upload} (accent 1\up{re} syllabe), \eng{device} (accent 2\up{e} syllabe).
\end{itemize}
\end{aretenir}

\coursec{Model dialogue 1: ICTs for learning}

Dans la section précédente (\emph{Warm-up: ICTs in our daily lives}), nous avons découvert le vocabulaire de base des technologies de l'information et de la communication : \eng{smartphone}, \eng{app}, \eng{website}, \eng{to download}, \eng{online}. Nous avons constaté que ces outils font partie de notre quotidien, à l'école comme à la maison. Il est maintenant temps de passer à l'action : comment parler de ces outils en anglais ? Comment demander à un camarade quelle application il utilise, et comment lui expliquer la nôtre ? C'est exactement ce que nous allons apprendre grâce à un premier dialogue modèle entre deux élèves, Kemi et Tabi.

\courssub{Listening and reading: the model dialogue}

Fermez d'abord vos cahiers et écoutez le dialogue lu par l'enseignant (ou joué par deux élèves volontaires). Écoutez une première fois pour saisir l'idée générale : qui parle ? De quoi parlent-ils ? Puis ouvrez vos cahiers et lisez le dialogue en suivant une deuxième écoute.

\begin{textebox}[Model dialogue — “Revising with an app”]
\textbf{Kemi:} Hi Tabi! You look busy. What are you doing on your phone?

\textbf{Tabi:} Hi Kemi! I'm revising English vocabulary with an app called Quizlet.

\textbf{Kemi:} Really? How does it work?

\textbf{Tabi:} It uses flashcards. It helps me memorise new words quickly.

\textbf{Kemi:} That sounds useful! Do you use any other apps?

\textbf{Tabi:} Yes, I watch English lessons on YouTube and I use an online dictionary.

\textbf{Kemi:} Wow! Could you show me how to download Quizlet?

\textbf{Tabi:} Sure, it's very easy. Let's do it after class.
\end{textebox}

\textbf{Glossaire (mots difficiles) :}

\begin{center}
\begin{tabularx}{\linewidth}{|l|l|X|}
\hline
\rowcolor{popblueL} \textbf{English} & \textbf{Nature} & \textbf{Français} \\
\hline
busy & adjectif & occupé \\
\hline
to revise & verbe & réviser \\
\hline
an app & nom & une application \\
\hline
a flashcard & nom & une carte mémoire (mot d'un côté, traduction de l'autre) \\
\hline
to memorise & verbe & mémoriser, apprendre par cœur \\
\hline
quickly & adverbe & rapidement, vite \\
\hline
useful & adjectif & utile \\
\hline
an online dictionary & nom & un dictionnaire en ligne \\
\hline
to download & verbe & télécharger \\
\hline
after class & locution & après le cours \\
\hline
\end{tabularx}
\end{center}

\begin{attention}[Prononciation des mots nouveaux]
Pas d'alphabet phonétique ici : on note la prononciation en lettres françaises.
\begin{itemize}
\item \eng{busy} se prononce « bi-zi » (et non « bou-ssi » : attention à l'orthographe trompeuse !).
\item \eng{memorise} se prononce « me-mo-raïz ».
\item \eng{download} se prononce « da-oun-loud ».
\item \eng{useful} se prononce « ious-foul ».
\item \eng{dictionary} se prononce « dik-chè-neu-ri » (3 syllabes, accent sur la première).
\end{itemize}
\end{attention}

\courssub{Global comprehension: who, what, which tools?}

Avant d'étudier les formules, vérifions la compréhension globale du dialogue. Répondez mentalement, puis à voix haute, aux questions suivantes :

\begin{enumerate}
\item Who are the two speakers? (\emph{Qui sont les deux locuteurs ?})
\item What is Tabi doing on her phone? (\emph{Que fait Tabi sur son téléphone ?})
\item Which app does she use to revise vocabulary? (\emph{Quelle application utilise-t-elle pour réviser le vocabulaire ?})
\item How does Quizlet work? (\emph{Comment Quizlet fonctionne-t-il ?})
\item Which other ICT tools does Tabi mention? (\emph{Quels autres outils TIC mentionne-t-elle ?})
\item What does Kemi ask at the end? (\emph{Que demande Kemi à la fin ?})
\end{enumerate}

\begin{corrige}[Compréhension globale]
\begin{enumerate}
\item The speakers are Kemi and Tabi, two classmates. (\emph{Deux camarades de classe.})
\item She is revising English vocabulary. (\emph{Elle révise du vocabulaire anglais.})
\item She uses an app called Quizlet. (\emph{Une application appelée Quizlet.})
\item It uses flashcards to help her memorise new words quickly. (\emph{Il utilise des cartes mémoire pour l'aider à mémoriser vite de nouveaux mots.})
\item She also watches English lessons on YouTube and uses an online dictionary. (\emph{YouTube et un dictionnaire en ligne.})
\item She asks Tabi to show her how to download Quizlet. (\emph{Elle demande à Tabi de lui montrer comment télécharger Quizlet.})
\end{enumerate}
\end{corrige}

\courssub{Spotting the functional formulas}

Un dialogue d'échange d'informations n'est pas une suite de phrases au hasard : il suit une \cle{structure} précise. Observons d'abord les répliques, puis dégageons la règle.

\begin{exemplebox}[Observer avant de nommer]
\begin{itemize}
\item \eng{Hi Tabi!} → salutation : c'est l'\textbf{ouverture} (opening).
\item \eng{What are you doing on your phone?} / \eng{How does it work?} / \eng{Do you use any other apps?} → questions : Kemi \textbf{demande une information} (asking).
\item \eng{I'm revising English vocabulary with an app called Quizlet.} / \eng{It uses flashcards.} → réponses : Tabi \textbf{donne une information} (informing).
\item \eng{Really?} / \eng{That sounds useful!} / \eng{Wow!} → Kemi \textbf{réagit} (reacting) et relance la conversation.
\item \eng{Let's do it after class.} → proposition qui \textbf{clôt} l'échange (closing).
\end{itemize}
\end{exemplebox}

\begin{propriete}[Structure d'un dialogue d'échange d'informations]
Un dialogue oral pour échanger des informations suit généralement cinq étapes :
\begin{enumerate}
\item \textbf{Opening} : saluer et engager la conversation (\eng{Hi!}, \eng{You look busy.}).
\item \textbf{Asking} : poser des questions ouvertes (\eng{What...?}, \eng{How...?}) ou fermées (\eng{Do you...?}).
\item \textbf{Informing} : donner l'information demandée, souvent avec une raison ou un détail (\eng{I use Quizlet to memorise vocabulary.}).
\item \textbf{Reacting} : montrer de l'intérêt pour relancer (\eng{That sounds great!}, \eng{Really?}).
\item \textbf{Closing} : conclure poliment (\eng{Let's do it after class.}, \eng{Thanks!}).
\end{enumerate}
La \textbf{réaction} est essentielle en anglais : un locuteur qui ne réagit pas paraît froid ou impoli, contrairement au français où l'on peut enchaîner les questions sans commentaire.
\end{propriete}

\begin{popfigure}
\begin{center}
\begin{tikzpicture}[node distance=0.55cm,
box/.style={draw=popblue, fill=popblueL, rounded corners, thick, minimum width=2.6cm, minimum height=1cm, align=center, font=\small},
arr/.style={-{Stealth[length=3mm]}, thick, poporange}]
\node[box, align=center] (open){\textbf{Opening}\\ \eng{Hi Tabi!}};
\node[box, right=of open, align=center] (ask){\textbf{Asking}\\ \eng{How does it work?}};
\node[box, right=of ask, align=center] (inform){\textbf{Informing}\\ \eng{It uses flashcards.}};
\node[box, right=of inform, align=center] (react){\textbf{Reacting}\\ \eng{That sounds useful!}};
\node[box, right=of react, align=center] (close){\textbf{Closing}\\ \eng{Let's do it after class.}};
\draw[arr] (open) -- (ask);
\draw[arr] (ask) -- (inform);
\draw[arr] (inform) -- (react);
\draw[arr] (react) -- (close);
\draw[arr, dashed, popgreen] (react.south) to[bend right=35] node[below, font=\scriptsize, popgreen]{relancer : nouvelle question} (ask.south);
\end{tikzpicture}
\end{center}
\legende{Structure d'un dialogue d'échange d'informations : la réaction permet de relancer le cycle question–réponse.}
\end{popfigure}

\courssub{Key formulas in the dialogue}

\begin{center}
\begin{tabularx}{\linewidth}{|l|X|X|}
\hline
\rowcolor{popblueL} \textbf{Fonction} & \textbf{English (dialogue)} & \textbf{Français} \\
\hline
Demander & \eng{What are you doing on your phone?} & Que fais-tu sur ton téléphone ? \\
\hline
Demander & \eng{How does it work?} & Comment ça marche ? \\
\hline
Demander & \eng{Do you use any other apps?} & Utilises-tu d'autres applications ? \\
\hline
Donner une information & \eng{I'm revising English vocabulary with an app called Quizlet.} & Je révise du vocabulaire anglais avec une application appelée Quizlet. \\
\hline
Donner une information & \eng{It helps me memorise new words quickly.} & Cela m'aide à mémoriser de nouveaux mots rapidement. \\
\hline
Réagir & \eng{Really?} & Vraiment ? \\
\hline
Réagir & \eng{That sounds useful!} & Ça a l'air utile ! \\
\hline
Demander un service & \eng{Could you show me how to download Quizlet?} & Pourrais-tu me montrer comment télécharger Quizlet ? \\
\hline
Conclure & \eng{Sure, it's very easy. Let's do it after class.} & Bien sûr, c'est très facile. Faisons ça après le cours. \\
\hline
\end{tabularx}
\end{center}

\begin{exemplebox}[Exemples traduits à retenir par cœur]
\begin{itemize}
\item \eng{It helps me memorise new words.} « Cela m'aide à mémoriser de nouveaux mots. »
\item \eng{Could you show me how to download it?} « Pourrais-tu me montrer comment le télécharger ? »
\item \eng{Which apps do you use to learn English?} « Quelles applications utilises-tu pour apprendre l'anglais ? »
\item \eng{I use Quizlet to memorise vocabulary.} « J'utilise Quizlet pour mémoriser le vocabulaire. »
\item \eng{That sounds great!} « Ça a l'air génial ! »
\end{itemize}
\end{exemplebox}

\begin{attention}[Pièges fréquents des francophones]
\begin{itemize}
\item \eng{It helps me \textbf{to} memorise} est accepté, mais la forme la plus naturelle est sans \eng{to} : \eng{It helps me memorise}. Ne dites jamais \emph{“It helps me for memorise”} : la préposition française « pour » ne se traduit pas ici.
\item \eng{How does it work?} : n'oubliez pas l'auxiliaire \eng{does} (3\textsuperscript{e} personne). Faute typique : \emph{“How it works?”}.
\item \eng{That sounds useful!} : \eng{to sound} signifie ici « avoir l'air », pas « sonner ». Ne traduisez pas « ça sonne utile ».
\item \eng{an app called Quizlet} : « une application \emph{appelée} Quizlet ». N'utilisez pas \eng{named} à l'oral courant, et jamais \emph{“called as Quizlet”}.
\item \eng{Could you show me...?} est plus poli que \eng{Can you show me...?} : préférez-le pour demander un service.
\end{itemize}
\end{attention}

\courssub{Learning the dialogue: method and choral repetition}

\begin{methode}[Comment apprendre un dialogue en cinq étapes]
\begin{enumerate}
\item \textbf{Écouter} le dialogue deux fois sans le lire : saisir l'idée générale (qui, quoi, quels outils).
\item \textbf{Lire à voix haute} en suivant le texte : souligner les mots difficiles et vérifier leur prononciation.
\item \textbf{Repérer les formules} : classer chaque réplique (demander, informer, réagir, conclure).
\item \textbf{Répéter en binôme} : d'abord en lisant, puis en cachant le texte, en soignant l'intonation (voix qui monte pour les questions fermées, qui descend pour les questions en \eng{wh-}).
\item \textbf{Remplacer les informations} par les siennes : changer les applications, les raisons et la conclusion pour créer son propre dialogue.
\end{enumerate}
\end{methode}

\textbf{Reprise en chœur (choral drilling).} La classe répète d'abord chaque réplique après l'enseignant, en insistant sur :
\begin{itemize}
\item l'accentuation : \eng{QUIZlet}, \eng{FLASHcards}, \eng{MEMorise}, \eng{DICtionary} (accent sur la première syllabe) ;
\item l'intonation montante de \eng{Really?} « ri-li ? » (la voix monte, comme en français) ;
\item l'intonation descendante de \eng{What are you doing?} (la voix descend à la fin des questions en \eng{wh-}) ;
\item l'enthousiasme de \eng{That sounds useful!} : phrase exclamative, voix haute puis descendante.
\end{itemize}
Ensuite, la classe se divise en deux groupes : le groupe A joue Kemi, le groupe B joue Tabi, puis on échange les rôles. Enfin, travail en binôme : chaque paire joue le dialogue sans regarder le texte.

\begin{exoresolu}[Classer les formules]
Énoncé — Voici cinq répliques du dialogue. Indiquez pour chacune sa fonction : \emph{asking}, \emph{informing}, \emph{reacting} ou \emph{closing}.
\begin{enumerate}
\item \eng{Do you use any other apps?}
\item \eng{It helps me memorise new words quickly.}
\item \eng{Wow!}
\item \eng{Let's do it after class.}
\item \eng{How does it work?}
\end{enumerate}
\tcblower
Solution détaillée :
\begin{enumerate}
\item \textbf{Asking} : question fermée en \eng{Do you...?} pour obtenir une information.
\item \textbf{Informing} : Tabi donne une explication sur l'utilité de l'application.
\item \textbf{Reacting} : exclamation qui montre l'intérêt et encourage le locuteur à continuer.
\item \textbf{Closing} : proposition qui met fin à l'échange en fixant une action future.
\item \textbf{Asking} : question ouverte en \eng{How...?} qui demande un mode de fonctionnement.
\end{enumerate}
Astuce : les répliques de réaction sont souvent courtes (\eng{Really?}, \eng{Wow!}, \eng{Great!}) et ne contiennent aucune information nouvelle.
\end{exoresolu}

\courssub{Guided substitution: make it your own dialogue}

Dernière étape : la \cle{substitution guidée}. On garde la structure du dialogue (opening → asking → informing → reacting → closing), mais on remplace les applications et les raisons par des exemples personnels. Voici une banque d'idées :

\begin{center}
\begin{tabularx}{\linewidth}{|X|X|}
\hline
\rowcolor{popgreenL} \textbf{ICT tool (à substituer)} & \textbf{Reason (à substituer)} \\
\hline
\eng{an app called Duolingo} & \eng{to practise grammar every day} \\
\hline
\eng{WhatsApp groups} & \eng{to discuss homework with classmates} \\
\hline
\eng{YouTube channels} & \eng{to improve my listening skills} \\
\hline
\eng{Google Translate} & \eng{to understand difficult words} \\
\hline
\eng{an online dictionary} & \eng{to check the pronunciation of words} \\
\hline
\eng{educational podcasts} & \eng{to learn while travelling by bendskin} \\
\hline
\end{tabularx}
\end{center}

\begin{experience}[Substitution drill en binôme]
Travail par deux. L'élève A joue le rôle de Kemi, l'élève B celui de Tabi. Remplacez \eng{Quizlet} par un outil de la banque ci-dessus (ou le vôtre) et la raison \eng{to memorise vocabulary} par une raison personnelle. Exemple de départ :

\eng{A: Hi B! You look busy. What are you doing on your phone?}

\eng{B: Hi A! I'm practising grammar with an app called Duolingo.}

\eng{A: Really? How does it work?}

Jouez le dialogue complet, puis échangez les rôles avec un autre outil. Deux binômes volontaires présenteront leur dialogue devant la classe.
\end{experience}

\begin{exercice}[Pour s'entraîner à la maison]
\begin{enumerate}
\item Recopiez le dialogue modèle en soulignant en rouge les questions et en vert les réactions.
\item Écrivez trois questions qu'un élève peut poser à un camarade sur les TIC (\eng{Which apps...?}, \eng{How often...?}, \eng{Why...?}).
\item Préparez votre propre version du dialogue (8 répliques minimum) avec une application que vous utilisez réellement, prête à être jouée au début de la section 4.
\end{enumerate}
\end{exercice}

\begin{aretenir}[L'essentiel de la section]
\begin{itemize}
\item Un dialogue d'échange d'informations suit cinq étapes : \textbf{opening → asking → informing → reacting → closing}.
\item Pour demander : \eng{Which apps do you use?}, \eng{How does it work?}
\item Pour informer : \eng{I use Quizlet to memorise vocabulary.}, \eng{It helps me learn quickly.}
\item Pour réagir et relancer : \eng{Really?}, \eng{That sounds useful!}, \eng{Wow!}
\item Pour conclure : \eng{Let's do it after class.}
\item Apprendre un dialogue = écouter, lire, repérer, répéter en binôme, substituer.
\end{itemize}
\end{aretenir}

\coursec{Functional language: asking, informing, follow-up, closing}

Dans la section précédente, nous avons écouté et lu un premier dialogue modèle dans lequel deux élèves, Amina et Peter, parlaient de l'utilisation des TIC pour apprendre (\eng{ICTs for learning}). Vous avez remarqué que leur conversation n'était pas une suite de phrases isolées : chacun \cle{posait des questions}, \cle{donnait des informations}, \cle{réagissait} aux propos de l'autre, puis \cle{concluait} poliment l'échange. Cette section a pour but de démonter ce mécanisme : nous allons isoler les \cle{formules fonctionnelles} (\eng{functional language}) qui font vivre une conversation en anglais, les classer, les expliquer, puis vous entraîner à les employer. À la fin de cette section, vous serez capable de mener un court dialogue complet, de la première question à la formule de congé.

\begin{definition}[Functional language / Speech acts]
En linguistique, on appelle \textbf{formules fonctionnelles} (ou \eng{speech acts} / \eng{language functions}) les morceaux de langue pré-fabriqués qui servent à accomplir une action précise dans la communication : demander, informer, exprimer son accord, réagir, prendre congé. Les maîtriser, c'est posséder la « boîte à outils » de l'oral.
\end{definition}

\courssub{Observing the language: what makes a conversation work?}

Relisons ensemble cet extrait du dialogue modèle 1 et observons les phrases en gras.

\begin{exemplebox}[Extract from Model Dialogue 1]
\textbf{Amina:} \textbf{Which apps do you use for learning}, Peter?\\
\textbf{Peter:} Well, \textbf{I use my phone to check} my English vocabulary every morning. And you?\\
\textbf{Amina:} \textbf{I use WhatsApp to discuss} homework with my classmates. \textbf{It helps me save time}.\\
\textbf{Peter:} \textbf{Really? That's interesting!} \textbf{How often do you use it}?\\
\textbf{Amina:} Almost every day. \textbf{I think ICTs make learning easier}.\\
\textbf{Peter:} \textbf{I agree}. Anyway, \textbf{I have to go now. It was nice talking to you!}
\end{exemplebox}

Observons : chaque réplique joue un rôle précis. On peut classer toutes ces formules en quatre grandes fonctions : \cle{demander} (\eng{asking}), \cle{informer} (\eng{informing}), \cle{réagir et relancer} (\eng{reacting and following up}), \cle{conclure} (\eng{closing}). C'est la « grammaire de la conversation » : un locuteur qui maîtrise ces quatre blocs peut tenir un échange naturel sur n'importe quel sujet, y compris les TIC au travail et à l'école.

\courssub{Asking for information: les questions de l'échange}

Pour obtenir une information sur l'usage des TIC, on utilise deux grands types de questions.

\textbf{1. Les questions fermées (oui/non)} : elles commencent par un auxiliaire (\eng{Do}, \eng{Does}, \eng{Can}, \eng{Have}...).

\eng{Do you often study online?} « Est-ce que tu étudies souvent en ligne ? »

\eng{Can you access the Internet at school?} « Peux-tu accéder à Internet à l'école ? »

\textbf{2. Les questions ouvertes (en \eng{Wh-})} : elles commencent par un mot interrogatif (\eng{Which}, \eng{What}, \eng{How}, \eng{How often}, \eng{Why}) suivi de l'auxiliaire, puis du sujet, puis du verbe.

\eng{Which apps do you use for learning?} « Quelles applications utilises-tu pour apprendre ? »

\eng{How do you use the Internet for your work?} « Comment utilises-tu Internet pour ton travail ? »

\eng{What do you use WhatsApp for?} « Pour quoi utilises-tu WhatsApp ? »

\begin{propriete}[La structure de la question en anglais]
Ordre obligatoire : \textbf{(mot interrogatif) + auxiliaire + sujet + verbe + complément}.\\
\eng{What do you use it for?} = « Pour quoi l'utilises-tu ? »\\
Remarquez la préposition \eng{for} placée \textbf{à la fin} de la question : c'est la structure naturelle en anglais parlé. \eng{For what do you use it?} est correct mais très formel et presque jamais utilisé à l'oral.
\end{propriete}

\begin{attention}[Piège : l'ordre des mots et l'article]
Le francophone traduit souvent mot à mot : « Quelles applications tu utilises ? » donne \eng{Which apps you use?} — c'est \textbf{faux}. L'auxiliaire \eng{do} est obligatoire : \eng{Which apps \textbf{do} you use?} De même, ne dites pas \eng{How you use Internet?} mais \eng{How \textbf{do} you use \textbf{the} Internet?} (Internet prend généralement \eng{the} à l'oral).
\end{attention}

\courssub{Informing: donner une information sur ses habitudes}

Pour répondre et décrire ses usages, l'anglais dispose de structures très fréquentes qu'il faut mémoriser comme des blocs.

\begin{propriete}[Exprimer l'utilité d'un outil : « I use + objet + to + verbe »]
Pour expliquer à quoi sert un outil, on utilise la structure \textbf{sujet + use + objet + to + base verbale} :\\
\eng{I use Google to do my research.} « J'utilise Google pour faire mes recherches. »\\
\eng{We use Zoom for our meetings.} « Nous utilisons Zoom pour nos réunions. » (variante : \eng{use + objet + for + nom / gérondif})\\
\eng{I use my phone to check my emails.} « J'utilise mon téléphone pour consulter mes e-mails. »\\
Le \eng{to} exprime ici le \textbf{but} (il correspond au « pour » français suivi de l'infinitif).
\end{propriete}

Pour ajouter un commentaire sur l'avantage de l'outil, deux formules sont précieuses :

\eng{It helps me save time.} « Cela m'aide à gagner du temps. » (forme courante : \eng{help + objet + base verbale})

\eng{It saves me a lot of time.} « Cela me fait gagner beaucoup de temps. »

\eng{Thanks to the Internet, I can find information quickly.} « Grâce à Internet, je peux trouver des informations rapidement. »

\begin{exemplebox}[Exemples traduits]
\eng{What do you use it for?} « Pour quoi l'utilises-tu ? »\\
\eng{I use it to send documents to my boss.} « Je l'utilise pour envoyer des documents à mon patron. »\\
\eng{It saves me a lot of time.} « Cela me fait gagner beaucoup de temps. »\\
\eng{Thanks to my tablet, I can read my lessons anywhere.} « Grâce à ma tablette, je peux lire mes leçons partout. »
\end{exemplebox}

\courssub{Giving an opinion, agreeing and disagreeing}

Un échange d'informations devient vite un échange d'opinions. Voici les formules de base, de la plus neutre à la plus nuancée.

\begin{itemize}
\item \eng{I think ICTs make learning easier.} « Je pense que les TIC rendent l'apprentissage plus facile. »
\item \eng{In my opinion, every student should have an email address.} « À mon avis, chaque élève devrait avoir une adresse e-mail. »
\item \eng{From my point of view, smartphones are essential.} « De mon point de vue, les smartphones sont indispensables. »
\item \eng{I agree.} « Je suis d'accord. »
\item \eng{I couldn't agree more.} « Je suis tout à fait d'accord. » (fort)
\item \eng{I don't quite agree because not everyone has a smartphone.} « Je ne suis pas tout à fait d'accord car tout le monde n'a pas de smartphone. »
\item \eng{I see what you mean, but...} « Je comprends ce que tu veux dire, mais... » (désaccord poli)
\end{itemize}

\begin{attention}[Faux ami de construction : « je suis d'accord »]
Ne traduisez \textbf{jamais} « je suis d'accord » par \eng{I am agree} : c'est une erreur grave et très fréquente. En anglais, \eng{agree} est un \textbf{verbe}, pas un adjectif. On dit donc simplement \eng{I agree}, \eng{I don't agree} ou, pour nuancer poliment, \eng{I don't quite agree}. De même : « il n'est pas d'accord » = \eng{He doesn't agree} (et non \eng{He isn't agree}).
\end{attention}

\courssub{Reacting and following up: relancer la conversation}

Un bon locuteur ne se contente pas de répondre : il \cle{réagit} (montre qu'il écoute) et \cle{relance} (pose une question de suivi). C'est ce qui rend un dialogue vivant et naturel.

\begin{itemize}
\item \textbf{Réactions courtes (show interest)} : \eng{Really?} « Vraiment ? » — \eng{That's interesting!} « C'est intéressant ! » — \eng{Wow!} — \eng{I see.} « Je vois / Je comprends. » — \eng{Right.} « D'accord / C'est ça. »
\item \textbf{Relances et demandes de précisions (follow-up)} : \eng{Could you tell me more?} « Pourrais-tu m'en dire plus ? » — \eng{What do you mean?} « Que veux-tu dire ? » — \eng{How often do you use it?} « À quelle fréquence l'utilises-tu ? » — \eng{Why do you prefer this app?} « Pourquoi préfères-tu cette appli ? »
\end{itemize}

\begin{methode}[Comment relancer une conversation en deux temps (Technique du « Rebond »)]
\begin{enumerate}
\item Commencez par un \textbf{mot de réaction} qui montre votre intérêt : \eng{Really?}, \eng{Wow!}, \eng{Interesting!}, \eng{That's great!}, \eng{Right.}
\item Enchaînez immédiatement avec une \textbf{question ouverte} en \eng{How}, \eng{What}, \eng{Why} ou \eng{Which} : \eng{How often do you use it?}, \eng{What do you use it for?}, \eng{Why do you prefer this app?}
\end{enumerate}
Exemple complet : \eng{Really? That's interesting! Why do you prefer Zoom to WhatsApp for your meetings?} « Vraiment ? C'est intéressant ! Pourquoi préfères-tu Zoom à WhatsApp pour tes réunions ? »
\end{methode}

\courssub{Closing: conclure poliment l'échange}

En anglais, on ne quitte jamais une conversation brutalement : on annonce la fin (\eng{signalling the end}), on remercie ou on exprime le plaisir d'avoir parlé, puis on prend congé.

\begin{itemize}
\item \eng{Well, I have to go now.} « Bon, je dois y aller maintenant. » (signal de fin)
\item \eng{Thank you for sharing.} « Merci d'avoir partagé (ces informations / ton avis). »
\item \eng{It was nice talking to you.} « C'était agréable de parler avec toi. »
\item \eng{Let's keep in touch online!} « Restons en contact en ligne ! »
\item \eng{See you later! / Catch you later!} « À plus tard ! / À la prochaine ! » (informel)
\end{itemize}

\courssub{Summary table and mind map}

\begin{center}
\begin{tabularx}{\linewidth}{|X|X|X|}
\hline
\rowcolor{popblueL} Fonction & Formule anglaise & Traduction française \\
\hline
Asking (closed) & Do you often study online? & Est-ce que tu étudies souvent en ligne ? \\
\hline
Asking (open) & Which apps do you use for learning? & Quelles applications utilises-tu pour apprendre ? \\
\hline
Asking (open) & How often do you go online? & À quelle fréquence te connectes-tu ? \\
\hline
Informing & I use my phone to check my emails. & J'utilise mon téléphone pour consulter mes e-mails. \\
\hline
Informing (benefit) & It helps me save time. / It saves me time. & Cela m'aide à gagner du temps. / Cela me fait gagner du temps. \\
\hline
Opinion & In my opinion, ICTs make work easier. & À mon avis, les TIC facilitent le travail. \\
\hline
Agreeing & I agree. / I couldn't agree more. & Je suis d'accord. / Je suis tout à fait d'accord. \\
\hline
Disagreeing politely & I don't quite agree. / I see your point, but... & Je ne suis pas tout à fait d'accord. / Je vois ton point, mais... \\
\hline
Reacting & Really? That's interesting! & Vraiment ? C'est intéressant ! \\
\hline
Follow-up & Could you tell me more? / What do you mean? & Pourrais-tu m'en dire plus ? / Que veux-tu dire ? \\
\hline
Closing & It was nice talking to you. See you! & C'était un plaisir de parler avec toi. À bientôt ! \\
\hline
\end{tabularx}
\end{center}

\begin{popfigure}
\begin{tikzpicture}[node distance=0.4cm]
\node[draw=popblue, fill=popblueL, rounded corners, minimum width=2.5cm, minimum height=1.5cm, align=center, font=\small\bfseries] (ask) at (0,0) {ASKING\\ \normalfont\footnotesize Which apps do you use?\\ \normalfont\footnotesize How often...?\\ \normalfont\footnotesize What... for?};
\node[draw=popgreen, fill=popgreenL, rounded corners, minimum width=2.5cm, minimum height=1.5cm, align=center, font=\small\bfseries, right=of ask] (inf) {INFORMING\\ \normalfont\footnotesize I use Zoom for...\\ \normalfont\footnotesize It saves me time.\\ \normalfont\footnotesize Thanks to... I can...};
\node[draw=poporange, fill=poporangeL, rounded corners, minimum width=2.5cm, minimum height=1.5cm, align=center, font=\small\bfseries, right=of inf] (rea) {REACTING \& FOLLOW-UP\\ \normalfont\footnotesize Really? Interesting!\\ \normalfont\footnotesize Tell me more.\\ \normalfont\footnotesize I agree / I don't quite agree.};
\node[draw=poppurple, fill=poppurpleL, rounded corners, minimum width=2.5cm, minimum height=1.5cm, align=center, font=\small\bfseries, right=of rea] (clo) {CLOSING\\ \normalfont\footnotesize Thanks for sharing.\\ \normalfont\footnotesize Nice talking to you.\\ \normalfont\footnotesize Let's keep in touch!};
\draw[-{Stealth}, thick, popmuted] (ask) -- (inf);
\draw[-{Stealth}, thick, popmuted] (inf) -- (rea);
\draw[-{Stealth}, thick, popmuted] (rea) -- (clo);
\end{tikzpicture}
\legende{Les quatre fonctions d'un échange oral : demander, informer, réagir/relancer, conclure.}
\end{popfigure}

\courssub{Guided practice: Mini role-play (Speaking activity)}

\begin{experience}[Speed-dating ICT : 4 minutes par binôme]
\textbf{Consigne :} En binôme, menez un dialogue complet de 2 minutes en suivant le schéma \textbf{Ask $\rightarrow$ Inform $\rightarrow$ React/Follow-up $\rightarrow$ Close}. Changez de partenaire 3 fois.
\begin{enumerate}
\item \textbf{Élève A} : Pose une question ouverte sur l'usage du téléphone pour les études/travail.
\item \textbf{Élève B} : Répond avec \eng{I use... to...} + un avantage (\eng{It saves me time}).
\item \textbf{Élève A} : Réagit (\eng{Really?}) + relance (\eng{How often...?} ou \eng{Why...?}).
\item \textbf{Élève B} : Répond (fréquence + opinion \eng{I think...}).
\item \textbf{Élève A} : Exprime accord/désaccord poli.
\item \textbf{Ensemble} : Concluez poliment (\eng{Nice talking to you}).
\end{enumerate}
\textbf{Aide-mémoire au tableau :} \eng{Which / What / How often} -- \eng{I use... to... / It helps...} -- \eng{Really? Why...?} -- \eng{Nice talking to you.}
\end{experience}

\begin{savaistu}[ICTs in Cameroon: WhatsApp, Mobile Money \& the Digital Divide]
Au Cameroun, l'usage des TIC est marqué par la prédominance du mobile. \textbf{WhatsApp} est l'outil n°1 pour l'éducation (groupes de classe, envoi de photos de notes, révisions bac) et le commerce informel. \textbf{Mobile Money} (MTN MoMo, Orange Money) a révolutionné le paiement sans compte bancaire. Cependant, la \textbf{fracture numérique} (\eng{digital divide}) persiste : coût de la data, couverture 4G inégale (zones rurales vs Yaoundé/Douala), électricité instable. Dans un échange oral, mentionner ces réalités (\eng{Data is expensive}, \eng{Network is poor in my village}) montre une maîtrise du vocabulaire \textbf{contextualisé}, très valorisée au Baccalauréat.
\end{savaistu}

\courssub{Exercice résolu et entraînement écrit}

\begin{exoresolu}[Reconstituer un dialogue cohérent]
Énoncé — Remettez les répliques suivantes dans l'ordre pour reconstituer un dialogue logique entre deux collègues de Douala.\\
a. \eng{Really? How often do you use it?}\\
b. \eng{I have to go now. It was nice talking to you!}\\
c. \eng{How do you use the Internet for your work, Sandra?}\\
d. \eng{Almost every day. It saves me a lot of time.}\\
e. \eng{I use my laptop to send reports to my manager.}
\tcblower
Solution détaillée — On suit la logique des quatre fonctions. 1) On commence par une question (\eng{asking}) : réplique \textbf{c}. 2) On donne l'information demandée avec la structure \eng{use + to} (\eng{informing}) : réplique \textbf{e}. 3) On réagit et on relance avec une question en \eng{How often} (\eng{reacting + follow-up}) : réplique \textbf{a}. 4) On répond à la relance en précisant la fréquence et l'avantage : réplique \textbf{d}. 5) On conclut poliment (\eng{closing}) : réplique \textbf{b}. Ordre correct : \textbf{c – e – a – d – b}. Dialogue complet : \eng{How do you use the Internet for your work, Sandra? — I use my laptop to send reports to my manager. — Really? How often do you use it? — Almost every day. It saves me a lot of time. — I have to go now. It was nice talking to you!}
\end{exoresolu}

\begin{exercice}[À vous de jouer : Complétez le dialogue]
Complétez ce dialogue avec les formules fonctionnelles de la leçon. Écrivez les phrases complètes en anglais.
\begin{enumerate}
\item \textbf{Paul:} \underline{\hspace{8cm}} (demande quelle application Nadia utilise pour apprendre l'anglais)
\item \textbf{Nadia:} I use an app to learn ten new words every day.
\item \textbf{Paul:} \underline{\hspace{8cm}} (réagis et demande pourquoi elle a choisi cette application)
\item \textbf{Nadia:} Because it is free and easy to use. I think ICTs make learning easier.
\item \textbf{Paul:} \underline{\hspace{8cm}} (exprime ton accord)
\item \textbf{Nadia:} \underline{\hspace{8cm}} (remercie et prends congé)
\end{enumerate}
\end{exercice}

\begin{corrige}[Exercice « À vous de jouer »]
\begin{enumerate}
\item \textbf{Paul:} \eng{Which app do you use to learn English?} (ou \eng{What app do you use for learning English?})
\item \textbf{Paul:} \eng{Really? That's interesting! Why did you choose this app?} (ou \eng{Why do you like it?})
\item \textbf{Paul:} \eng{I agree.} (ou \eng{I quite agree with you.} / \eng{I couldn't agree more.})
\item \textbf{Nadia:} \eng{Thank you for sharing. It was nice talking to you. See you!}
\end{enumerate}
\textbf{Attention aux erreurs à éviter :} pas de \eng{I am agree}, pas de question sans auxiliaire (\eng{Which app you use?} est faux), n'oubliez pas \eng{the} devant \eng{Internet} si utilisé.
\end{corrige}

\begin{aretenir}[L'essentiel de la section]
Un échange oral réussi suit quatre étapes : \cle{demander} (questions avec auxiliaire : \eng{Which apps do you use?}), \cle{informer} (structure \eng{I use + objet + to + verbe} / \eng{It saves me time}), \cle{réagir et relancer} (mot de réaction + question ouverte : \eng{Really? How often...?}), \cle{conclure} (\eng{It was nice talking to you.}). Retenez aussi les deux pièges majeurs : \eng{I agree} (jamais \eng{I am agree}) et l'auxiliaire obligatoire dans les questions (\eng{do/does/did}). Vous êtes maintenant prêts pour le deuxième dialogue modèle, consacré aux TIC au travail (\eng{ICTs at work}).
\end{aretenir}

\coursec{Model dialogue 2: ICTs at work}

Dans la section précédente, nous avons appris les \cle{formules fonctionnelles} pour demander une information (\eng{May I ask you...?}), la donner (\eng{Well, actually...}), relancer (\eng{Could you tell me more?}) et conclure (\eng{Thank you for your time}). Nous allons maintenant voir ces formules \emph{en action} dans un contexte nouveau : le \cle{monde du travail}. Après les TIC pour apprendre (section 2), observons comment les TIC servent à \emph{travailler}, à travers une interview menée à Douala.

\courssub{Observing the model dialogue}

Lisez d'abord le dialogue à voix haute, en faisant attention à la politesse du journaliste et aux réponses complètes de l'employée. Ne cherchez pas encore à tout traduire : repérez simplement (1) les formules déjà étudiées, (2) les outils technologiques mentionnés, (3) un avantage et un inconvénient des TIC au travail.

\begin{textebox}[Model dialogue — A survey on ICTs at work (Douala)]
\textbf{Journalist:} Excuse me, madam. I'm doing a survey on ICTs at work. May I ask you a few questions?\\
\textbf{Employee:} Of course, go ahead.\\
\textbf{Journalist:} How do you use technology in your job?\\
\textbf{Employee:} Well, I use my computer to write reports and send emails to my colleagues. We also have online meetings on Zoom every Monday.\\
\textbf{Journalist:} Do ICTs make your work easier?\\
\textbf{Employee:} Absolutely! Communication is much faster now. But sometimes power cuts disturb our work.\\
\textbf{Journalist:} I see. Thank you very much for your time.\\
\textbf{Employee:} You're welcome.
\end{textebox}

\textbf{Glossaire (mots difficiles)}

\begin{center}
\begin{tabularx}{\linewidth}{|l|l|X|}
\hline
\rowcolor{popblueL} English & Nature & Français \\
\hline
survey & nom & enquête, sondage \\
\hline
to go ahead & expression & allez-y, je vous en prie \\
\hline
report & nom & rapport, compte rendu \\
\hline
colleague & nom & collègue \\
\hline
online meeting & nom & réunion en ligne \\
\hline
power cut & nom & coupure d'électricité (délestage) \\
\hline
to disturb & verbe & perturber, déranger \\
\hline
\end{tabularx}
\end{center}

\textbf{Questions de compréhension (oral)}

\begin{enumerate}
\item What is the journalist doing? \emph{(He is doing a survey on ICTs at work.)}
\item What does the employee use her computer for? \emph{(To write reports and send emails.)}
\item How often do they have online meetings? \emph{(Every Monday.)}
\item What advantage of ICTs does she mention? \emph{(Communication is much faster.)}
\item What problem does she mention? \emph{(Power cuts disturb their work.)}
\end{enumerate}

\courssub{Professional vocabulary: ICTs at work}

Enrichissons maintenant le lexique du dialogue avec le vocabulaire professionnel essentiel de cette leçon. Notez la prononciation en lettres françaises.

\begin{definition}[ICT vocabulary for the workplace]
\begin{itemize}
\item \eng{to send emails} « tou send i-mèils » = envoyer des courriels
\item \eng{to attend online meetings} « tou a-tènd on-laïne mi-tingz » = participer à des réunions en ligne
\item \eng{video conferencing} « vi-di-o kon-feu-ren-ssing » = la visioconférence
\item \eng{to work remotely} « tou weurk ri-mo-teu-li » = travailler à distance, télétravailler
\item \eng{a spreadsheet} « eu chprèd-chiite » = un tableur (Excel)
\item \eng{a printer} « eu prin-teur » = une imprimante
\item \eng{a scanner} « eu sska-neur » = un scanner, un numériseur
\end{itemize}
\end{definition}

\begin{exemplebox}[Exemples en contexte]
\eng{I send emails to my colleagues every morning.} « J'envoie des courriels à mes collègues chaque matin. »

\eng{Our manager attends online meetings with the head office in Yaoundé.} « Notre directeur participe à des réunions en ligne avec le siège à Yaoundé. »

\eng{Many banks now use video conferencing to train their staff.} « Beaucoup de banques utilisent désormais la visioconférence pour former leur personnel. »

\eng{She works remotely from Buea twice a week.} « Elle télétravaille depuis Buea deux fois par semaine. »

\eng{The accountant uses a spreadsheet to calculate salaries.} « Le comptable utilise un tableur pour calculer les salaires. »

\eng{Please scan this document with the scanner and print it with the printer.} « Veuillez numériser ce document avec le scanner et l'imprimer avec l'imprimante. »
\end{exemplebox}

\begin{popfigure}
\begin{center}
\begin{tikzpicture}[mindmap, grow cyclic,
  every node/.style={concept, font=\small, align=center},
  root concept/.append style={concept color=popblue, font=\bfseries},
  level 1/.append style={level distance=3.4cm, sibling angle=90},
  level 1 concept/.append style={concept color=poporangeL, text=popdark}]
\node[root concept]{ICTs at work}
  child { node[align=center]{Communication\\ \eng{send emails}\\ \eng{online meetings}\\ \eng{video conferencing}} }
  child { node[align=center]{Documents\\ \eng{write reports}\\ \eng{spreadsheet}\\ \eng{printer, scanner}} }
  child { node[align=center]{Organisation\\ \eng{work remotely}\\ \eng{share files}\\ \eng{save time}} }
  child { node[align=center]{Problems\\ \eng{power cuts}\\ \eng{poor connection}} };
\end{tikzpicture}
\end{center}
\legende{Carte mentale du lexique : les TIC au travail}
\end{popfigure}

\courssub{Spotting advantages and disadvantages}

Le dialogue exprime clairement un \cle{avantage} (\eng{Communication is much faster now}) et un \cle{inconvénient} (\eng{power cuts disturb our work}). Pour parler des TIC au travail, il faut savoir classer les arguments.

\begin{propriete}[Exprimer avantages et inconvénients]
Pour présenter un avantage : \eng{ICTs allow us to...} (les TIC nous permettent de...), \eng{Thanks to ICTs, we can...} (grâce aux TIC, nous pouvons...), \eng{The main advantage is...} (le principal avantage est...).

Pour présenter un inconvénient : \eng{However, ...} (cependant...), \eng{The main drawback is...} (le principal inconvénient est...), \eng{Unfortunately, ...} (malheureusement...).

Après \eng{allow} et \eng{enable}, on utilise la structure : \textbf{allow + quelqu'un + to + base verbale}. \eng{ICTs allow us to work from home.} « Les TIC nous permettent de travailler depuis la maison. »
\end{propriete}

\begin{popfigure}
\begin{center}
\begin{tikzpicture}
\node[draw=popgreen, fill=popgreenL, rounded corners, text=popdark, minimum width=6.2cm, minimum height=0.9cm, font=\bfseries] (adv) at (0,2.6) {Advantages (+)};
\node[draw=popgreen!60, fill=white, rounded corners, align=left, text=popdark, minimum width=6.2cm] (advlist) at (0,0.9) {\eng{fast communication}\\ \eng{saving time}\\ \eng{remote work (working from home)}\\ \eng{easy access to information}};
\node[draw=poppink, fill=poppinkL, rounded corners, text=popdark, minimum width=6.2cm, minimum height=0.9cm, font=\bfseries] (dis) at (7.4,2.6) {Disadvantages (−)};
\node[draw=poppink!60, fill=white, rounded corners, align=left, text=popdark, minimum width=6.2cm] (dislist) at (7.4,0.9) {\eng{power cuts}\\ \eng{poor Internet connection}\\ \eng{distractions (social media)}\\ \eng{expensive equipment}};
\end{tikzpicture}
\end{center}
\legende{Advantages vs Disadvantages of ICTs at work}
\end{popfigure}

\begin{exemplebox}[Avantages et inconvénients en phrases complètes]
\eng{Thanks to emails, communication is much faster than before.} « Grâce aux courriels, la communication est bien plus rapide qu'avant. »

\eng{Video conferencing saves time: we don't need to travel to Yaoundé for every meeting.} « La visioconférence fait gagner du temps : nous n'avons pas besoin de nous rendre à Yaoundé pour chaque réunion. »

\eng{Remote working allows employees to work from home.} « Le télétravail permet aux employés de travailler depuis chez eux. »

\eng{However, power cuts sometimes disturb our work.} « Cependant, les coupures d'électricité perturbent parfois notre travail. »

\eng{The main drawback is the poor Internet connection in some areas.} « Le principal inconvénient est la mauvaise connexion Internet dans certaines zones. »
\end{exemplebox}

\begin{attention}[Faux amis et noms indénombrables]
\textbf{1.} \eng{an information} n'existe pas : \emph{information} est \cle{indénombrable} en anglais. Dites \eng{a piece of information} (une information) ou \eng{some information} (des informations). De même, \emph{advice} est indénombrable : \eng{some advice} (des conseils), \eng{a piece of advice} (un conseil) — jamais \eng{an advice}.

\textbf{2.} \eng{to assist} ne veut pas dire « assister à » (au sens de « être présent ») mais « aider ». Pour « assister à une réunion », dites \eng{to attend a meeting}.

\textbf{3.} \eng{actually} signifie « en fait, réellement », pas « actuellement ». Pour « actuellement », dites \eng{currently} ou \eng{at present}.

\textbf{4.} On dit \eng{to work remotely} ou \eng{to work from home}, jamais \eng{to do home office} (calque de l'allemand, à éviter).
\end{attention}

\begin{exoresolu}[Compléter un mini-dialogue]
Énoncé — Complétez ce dialogue avec les expressions suivantes : \eng{power cuts}, \eng{a piece of information}, \eng{online meetings}, \eng{work remotely}, \eng{spreadsheet}.

\eng{A: Do you use technology in your office?}

\eng{B: Yes. I use a (1) \dots\ to prepare the budget, and we have (2) \dots\ every Friday.}

\eng{A: Do you ever (3) \dots\ from home?}

\eng{B: Sometimes, but (4) \dots\ make it difficult. By the way, could you give me (5) \dots\ about your Internet provider?}

\tcblower
Solution détaillée :
\begin{enumerate}
\item \eng{spreadsheet} — on prépare un budget avec un \emph{tableur}.
\item \eng{online meetings} — « nous avons des \emph{réunions en ligne} chaque vendredi » ; le pluriel convient avec \eng{every Friday}.
\item \eng{work remotely} — après \eng{Do you ever...}, il faut la base verbale : « télétravailler depuis chez soi ».
\item \eng{power cuts} — sujet pluriel suivi de \eng{make} (sans -s) : « les coupures d'électricité rendent cela difficile ».
\item \eng{a piece of information} — \emph{information} est indénombrable ; pour parler d'une seule information, on utilise \eng{a piece of}.
\end{enumerate}
\end{exoresolu}

\begin{savaistu}[Remote working around the world]
Au Royaume-Uni et aux États-Unis, le télétravail (\eng{remote working} ou \eng{working from home}, souvent abrégé \eng{WFH}) s'est généralisé depuis 2020 grâce à des outils comme \emph{Zoom} et \emph{Microsoft Teams}. Beaucoup d'entreprises ont adopté le \eng{hybrid work} (travail hybride) : quelques jours au bureau, quelques jours à la maison. Au Cameroun, le télétravail se développe aussi dans les banques et les entreprises de Douala et de Yaoundé, même si les coupures d'électricité (\eng{power cuts}) et la qualité de la connexion restent des défis quotidiens.
\end{savaistu}

\courssub{Guided role-play: the journalist and the employee}

\begin{experience}[Role-play — An interview on ICTs at work]
Travaillez par deux. L'élève A joue le \textbf{journaliste}, l'élève B joue un(e) \textbf{employé(e)} (banque, école, hôpital ou cybercafé de votre quartier). Le journaliste \emph{doit} poser les questions de la trame ci-dessous ; l'employé(e) répond avec des phrases complètes et ajoute au moins un avantage et un inconvénient. Puis échangez les rôles.

\textbf{Trame de questions imposées (journaliste) :}
\begin{enumerate}
\item \eng{Excuse me. May I ask you a few questions about ICTs at work?}
\item \eng{How do you use technology in your job?}
\item \eng{Do you ever work remotely or attend online meetings?}
\item \eng{What are the advantages of ICTs for you?}
\item \eng{Do you face any problems, like power cuts or a poor Internet connection?}
\item \eng{Thank you very much for your time.}
\end{enumerate}

\textbf{Aide pour l'employé(e) :} commencez par \eng{Well, ...} ou \eng{Actually, ...} ; utilisez \eng{I use... to...}, \eng{ICTs allow us to...}, \eng{However, ...} ; terminez par \eng{You're welcome.}
\end{experience}

\begin{methode}[Réussir le jeu de rôle]
\begin{enumerate}
\item \textbf{Préparez} : choisissez un métier et notez trois outils TIC que ce métier utilise (\eng{computer, printer, scanner, spreadsheet...}).
\item \textbf{Écoutez} la question avant de répondre : répondez précisément à ce qui est demandé.
\item \textbf{Développez} chaque réponse avec un détail (\eng{every Monday}, \eng{twice a week}) et un connecteur (\eng{also}, \eng{but}, \eng{however}).
\item \textbf{Soignez la prononciation} : accentuez le mot important de la phrase (\eng{Communication is much FASter now}) et montez le ton à la fin des questions en \eng{yes/no}.
\item \textbf{Concluez poliment} : \eng{Thank you for your time. — You're welcome.}
\end{enumerate}
\end{methode}

\begin{exercice}[Your turn — Build your own dialogue]
En vous inspirant du dialogue modèle, écrivez puis jouez un dialogue de huit répliques entre un journaliste et un travailleur de votre choix (enseignant, infirmier, commerçant au marché de Mokolo). Votre dialogue doit contenir : deux questions du journaliste, deux outils TIC, un avantage (\eng{faster communication} ou \eng{saving time}) et un inconvénient (\eng{power cuts} ou \eng{poor Internet connection}).
\end{exercice}

\begin{corrige}[Exercice — Build your own dialogue]
Corrigé-type (modèle de production attendue) :

\eng{A: Excuse me, sir. May I ask you a few questions about ICTs at work?}

\eng{B: Of course, go ahead.}

\eng{A: How do you use technology in your shop?}

\eng{B: Well, I use my phone to send messages to customers and a spreadsheet to manage sales.}

\eng{A: What is the main advantage for you?}

\eng{B: Communication is much faster and I save time. However, power cuts disturb my work.}

\eng{A: Thank you very much for your time.}

\eng{B: You're welcome.}

Critères de réussite : formules de politesse respectées, vocabulaire de la leçon utilisé correctement, au moins un avantage et un inconvénient exprimés, phrases complètes et prononciation claire.
\end{corrige}

\begin{aretenir}[Ce qu'il faut retenir]
\begin{itemize}
\item Vocabulaire du travail : \eng{to send emails, to attend online meetings, video conferencing, to work remotely, spreadsheet, printer, scanner}.
\item Avantages des TIC : \eng{faster communication, saving time, working from home} ; inconvénients : \eng{power cuts, poor Internet connection, distractions}.
\item Structure clé : \eng{allow + quelqu'un + to + verbe} (\eng{ICTs allow us to work faster}).
\item Pièges : \eng{information} et \eng{advice} sont indénombrables (\eng{a piece of information}) ; \eng{attend a meeting} = assister à une réunion ; \eng{actually} = en fait.
\item Un bon dialogue d'interview suit toujours la même architecture : salutation polie, questions ouvertes, réponses développées, conclusion courtoise.
\end{itemize}
\end{aretenir}

\coursec{Pronunciation, stress and intonation}

Dans la section précédente, nous avons écouté et joué le dialogue \emph{ICTs at work}, dans lequel deux collègues échangent des informations sur les outils numériques. Vous avez remarqué que les mots du monde des TIC (\eng{computer}, \eng{Internet}, \eng{download}) reviennent sans cesse. Mais bien parler anglais, ce n'est pas seulement utiliser les bons mots : c'est aussi les \cle{accentuer} correctement et faire \cle{monter} ou \cle{descendre} la voix au bon moment. C'est l'objet de cette section : prononciation, accentuation (\emph{stress}) et intonation.

\courssub{Word stress: stressing the right syllable}

\courssub{Observing first: listen and repeat}

Avant toute règle, lisons à voix haute ce court échange. Les majuscules indiquent les syllabes \textbf{accentuées}, c'est-à-dire prononcées plus fort, plus longtemps et plus haut que les autres.

\begin{textebox}[Read aloud — stressed syllables in CAPITALS]
— Do YOU use the InTERnet for your HOMEwork?

— Yes, I DO. I WATCH videos on YouTUBE and I DOWNload E-books.

— REAlly ↗? Which APpliCAtion do you preFER ↘?

— I preFER KHAN AcadeMY. The TECHnology is SIMple and the InforMAtion is CLEAR.

— That's GREAT ↘! My COMputer is OLD, but my SMARTphone works WELL.
\end{textebox}

\begin{definition}[Le stress (accentuation) en anglais]
En anglais, chaque mot de plus d'une syllabe possède une syllabe \cle{accentuée} (\emph{stressed syllable}), prononcée plus fort, plus longue et plus haute. Les autres syllabes sont \emph{faibles}, souvent réduites à un son « e » sourd (comme dans le français « le »). Contrairement au français, où l'accent tombe presque toujours sur la \textbf{dernière} syllabe du groupe de mots, l'anglais accentue une syllabe \textbf{à l'intérieur du mot}, et sa place est propre à chaque mot : il faut l'apprendre avec le mot.
\end{definition}

\begin{propriete}[Accentuation des mots clés des TIC]
Voici les mots essentiels de ce chapitre, avec leur syllabe accentuée (en majuscules) et la prononciation notée en lettres françaises.

\begin{center}
\begin{tabularx}{\linewidth}{|l|X|X|}\hline
\rowcolor{popblueL} Mot & Syllabe accentuée & Prononciation (lettres françaises) \\ \hline
COMputer & 1re syllabe & « keum-PIOU-teur » \\ \hline
InTERnet & 2e syllabe & « ine-TER-net » \\ \hline
APpliCAtion & 1re et 3e syllabes & « a-pli-KÉ-chn » \\ \hline
TECHnology & 2e syllabe & « tek-NO-lo-dji » \\ \hline
DOWNload (nom) / downLOAD (verbe) & 1re (nom) / 2e (verbe) & « DAONE-lode » / « daone-LODE » \\ \hline
InforMAtion & 3e syllabe & « ine-feur-MÉ-chn » \\ \hline
\end{tabularx}
\end{center}

Remarquez le cas de \eng{download} : comme beaucoup de mots anglais à deux syllabes, le \textbf{nom} est accentué sur la première syllabe (\eng{a DOWNload}) et le \textbf{verbe} sur la deuxième (\eng{to downLOAD a file}). Même phénomène avec \eng{an IMport / to imPORT}, \eng{a REcord / to reCORD}.
\end{propriete}

\begin{attention}[Le piège de l'accent final français]
Les francophones accentuent naturellement la dernière syllabe : ils disent « in-ter-NÈTE », « in-for-ma-CIONE », « tech-no-lo-GUIE ». Résultat : l'interlocuteur anglophone peine à reconnaître le mot. Entraînez-vous à « frapper » la bonne syllabe et à avaler les autres : \eng{COM-p'ter}, \eng{in-TER-net}, \eng{in-for-MA-tion}. Taper doucement du pied ou de la main sur la syllabe accentuée aide beaucoup.
\end{attention}

\courssub{Intonation: the music of English questions}

\courssub{Observation: two questions, two melodies}

Écoutez (ou lisez) ces deux questions tirées de nos dialogues :

\eng{Do you study online ↗?} « Est-ce que tu étudies en ligne ? »

\eng{Which app do you prefer ↘?} « Quelle application préfères-tu ? »

Dans la première, la voix \textbf{monte} sur le dernier mot ; dans la seconde, elle \textbf{descend}. Ce n'est pas un hasard : c'est une règle fondamentale de l'anglais parlé.

\begin{methode}[Choisir la bonne intonation pour une question]
\begin{enumerate}
\item La question attend-elle \eng{yes} ou \eng{no} (question fermée, commençant par \eng{do}, \eng{does}, \eng{can}, \eng{have}, \eng{is}...) ? Alors la voix \textbf{monte} à la fin : \eng{Do you use a smartphone ↗?}
\item La question commence-t-elle par un mot interrogatif en \emph{Wh-} (\eng{what}, \eng{which}, \eng{where}, \eng{when}, \eng{why}, \eng{how}) ? Alors la voix \textbf{descend} à la fin : \eng{What do you use it for ↘?}
\item En cas de doute, demandez-vous : « Puis-je répondre simplement oui ou non ? » Si oui, intonation montante ; sinon, descendante.
\end{enumerate}
\end{methode}

\begin{popfigure}
\begin{tikzpicture}[>=Stealth, scale=1]
% Question fermée : intonation montante
\node[font=\bfseries, popblue] at (3,2.6) {Do you study online?};
\draw[very thick, popblue, ->] (0.6,0.4) .. controls (2,0.5) and (3.4,0.8) .. (5.4,2.0);
\node[popblue, align=center, font=\small] at (3,-0.4) {yes/no question : la voix monte ↗};
% Question en Wh- : intonation descendante
\node[font=\bfseries, poporange] at (10.6,2.6) {Which app do you prefer?};
\draw[very thick, poporange, ->] (8.2,2.0) .. controls (9.6,1.8) and (11.4,1.0) .. (13.0,0.4);
\node[poporange, align=center, font=\small] at (10.6,-0.4) {Wh- question : la voix descend ↘};
\end{tikzpicture}
\legende{Deux mélodies pour deux types de questions : montante pour les questions fermées, descendante pour les questions en Wh-.}
\end{popfigure}

\begin{exemplebox}[Questions montantes et descendantes]
\begin{itemize}
\item \eng{Do you have a laptop ↗?} « As-tu un ordinateur portable ? »
\item \eng{Can you send me the file ↗?} « Peux-tu m'envoyer le fichier ? »
\item \eng{Is the connection good in Buea ↗?} « La connexion est-elle bonne à Buea ? »
\item \eng{What do you use your tablet for ↘?} « Pour quoi utilises-tu ta tablette ? »
\item \eng{How often do you check your e-mails ↘?} « À quelle fréquence consultes-tu tes e-mails ? »
\item \eng{Why do you prefer this website ↘?} « Pourquoi préfères-tu ce site web ? »
\end{itemize}
\end{exemplebox}

\courssub{Intonation of reactions: Really ↗? That's GREAT ↘!}

L'intonation porte aussi les \textbf{émotions}. Une même phrase peut changer de sens selon la mélodie.

\begin{exemplebox}[Really : surprise ou doute ?]
\eng{Really ↗?} dit avec une intonation \textbf{montante} = « Vraiment ? » : surprise sincère, intérêt, envie d'en savoir plus.

\eng{Really ↘.} dit avec une intonation \textbf{descendante} = doute ou ironie : « Ah bon... j'en doute. »

\eng{That's GREAT ↘!} avec une voix qui monte puis retombe fortement = enthousiasme authentique : « C'est génial ! »
\end{exemplebox}

\begin{propriete}[Réactions enthousiastes : la mélodie de l'exclamation]
Dans les exclamations positives (\eng{Great!}, \eng{Wonderful!}, \eng{Excellent!}, \eng{That's amazing!}), la voix part \textbf{haut} sur la syllabe accentuée puis \textbf{descend} nettement. Une exclamation dite à voix plate sonne faux ou ironique en anglais. Pour relancer la conversation avec intérêt, utilisez au contraire la montée : \eng{Really ↗? And then ↗?}
\end{propriete}

\courssub{Difficult sounds for French speakers}

Trois sons des mots de ce chapitre posent problème aux francophones. Voici comment les produire.

\begin{propriete}[Trois sons à travailler]
\begin{itemize}
\item \textbf{Le « th »} de \eng{technology}, \eng{the}, \eng{think} : placez le bout de la \textbf{langue entre les dents} et soufflez doucement. Pour \eng{the} et \eng{this}, faites vibrer les cordes vocales (son sonore) ; pour \eng{think} et \eng{technology}, soufflez sans vibrer (son sourd).
\item \textbf{Le « w »} de \eng{work}, \eng{web}, \eng{website} : arrondissez fortement les \textbf{lèvres} comme pour souffler une bougie, puis ouvrez-les rapidement. Ne le confondez pas avec « v » : \eng{west} n'est pas « vest ».
\item \textbf{Le « r » anglais} de \eng{Internet}, \eng{work}, \eng{prefer} : il ne se \textbf{roule pas}. La langue se recourbe vers le palais \textbf{sans le toucher}. Dites « ine-TER-net » avec un « r » très doux, presque un « w » léger.
\end{itemize}
\end{propriete}

\begin{attention}[« Ze » au lieu de « the » : l'erreur n° 1 des francophones]
Beaucoup de francophones prononcent \eng{the} comme « ze », \eng{think} comme « sink » ou « fink », \eng{technology} comme « té-kno-lo-gui ». Ces substitutions peuvent créer des malentendus (\eng{I think so} / « I sink so »). Remède : langue entre les dents, souffle, et entraînement quotidien sur la série \eng{the — this — that — think — technology — three}. De même, ne prononcez jamais le « r » anglais à la française (roulé) : dans \eng{Internet}, gardez un « r » mou.
\end{attention}

\begin{center}
\begin{tabularx}{\linewidth}{|X|X|X|}\hline
\rowcolor{popblueL} Français (habitude) & English (cible) & Remarque \\ \hline
« ze In-ter-NÈTE » & the InTERnet & Langue entre les dents ; accent sur TER \\ \hline
« ouèrk » avec « w » faible & work & Lèvres bien arrondies, « r » non roulé \\ \hline
« té-kno-lo-GUIE » & TECHnology & Accent sur NO ; « th » soufflé \\ \hline
Voix plate pour toutes les questions & Do you study online ↗? & Monter la voix pour yes/no \\ \hline
« Vraiment ? » plat & Really ↗? & La montée exprime la surprise \\ \hline
\end{tabularx}
\end{center}

\begin{exoresolu}[Questions : monter ou descendre ?]
Énoncé : pour chaque question, indiquez si l'intonation finale monte (↗) ou descend (↘), puis prononcez-la correctement.

1. \eng{Do you download e-books?}
2. \eng{Which website do you visit most?}
3. \eng{Can ICTs improve your marks?}
4. \eng{How do you use your smartphone at school?}
5. \eng{Is your computer new?}

\tcblower

Solution détaillée :

1. \eng{Do you download e-books ↗?} — question fermée (réponse \eng{yes/no}, auxiliaire \eng{do}) : la voix \textbf{monte}.
2. \eng{Which website do you visit most ↘?} — question en \emph{Wh-} (\eng{which}) : la voix \textbf{descend}.
3. \eng{Can ICTs improve your marks ↗?} — question fermée avec \eng{can} : la voix \textbf{monte}.
4. \eng{How do you use your smartphone at school ↘?} — mot interrogatif \eng{how} : la voix \textbf{descend}.
5. \eng{Is your computer new ↗?} — question fermée avec \eng{is} : la voix \textbf{monte}.

Astuce de vérification : si la réponse possible est « oui » ou « non », montez ; si la réponse est une information, descendez.
\end{exoresolu}

\begin{experience}[Choral repetition and pair work with self-correction]
\begin{enumerate}
\item \textbf{Chœur} : la classe répète ensemble après le professeur les mots accentués : \eng{COMputer — InTERnet — APpliCAtion — TECHnology — DOWNload — InforMAtion}, en tapant dans les mains sur la syllabe accentuée.
\item \textbf{Chœur} : répétition des questions avec de grands gestes de la main qui montent (↗) ou descendent (↘) : \eng{Do you use a smartphone ↗? — What do you use it for ↘?}
\item \textbf{Binômes} : l'élève A lit le dialogue de la boîte \emph{Read aloud} ; l'élève B écoute et lève la main à chaque erreur d'accentuation ou d'intonation. Puis on échange les rôles.
\item \textbf{Auto-correction} : chaque élève s'enregistre (téléphone ou magnétophone de la classe), se réécoute et note deux points à améliorer : un stress, une intonation.
\end{enumerate}
\end{experience}

\begin{aretenir}[L'essentiel de la prononciation]
\begin{itemize}
\item Chaque mot anglais a sa syllabe accentuée : \eng{COMputer}, \eng{InTERnet}, \eng{TECHnology}, \eng{InforMAtion}. Apprenez le stress avec le mot.
\item Question fermée (yes/no) = intonation \textbf{montante} ↗ ; question en \emph{Wh-} = intonation \textbf{descendante} ↘.
\item \eng{Really ↗?} = surprise et intérêt ; \eng{Really ↘.} = doute ou ironie.
\item « th » : langue entre les dents ; « w » : lèvres arrondies ; « r » anglais : jamais roulé.
\item Une bonne prononciation rend vos échanges sur les TIC clairs, naturels et agréables à écouter.
\end{itemize}
\end{aretenir}

\coursec{Guided role-plays and free practice}

Vous savez maintenant prononcer les mots-clés des TIC avec le bon accent et poser vos questions avec une intonation montante ou descendante correcte. Il est temps de mettre tout cela en pratique~: dans cette dernière étape de la leçon, vous allez \cle{échanger des informations} réelles avec vos camarades, d'abord avec un support très guidé, puis de façon de plus en plus libre. C'est exactement la progression exigée à l'épreuve orale du Baccalauréat~: de la reproduction guidée à l'interaction spontanée.

\courssub{How to succeed in a role-play}

Avant de commencer, observons ce que font les bons candidats à l'oral. Comparez ces deux extraits~:

\begin{exemplebox}[Observer : deux performances]
\textbf{Candidate A (faible)~:}\\
— Do you use a computer? — Yes. — Do you like it? — Yes. — Goodbye.

\textbf{Candidate B (excellente)~:}\\
— Good morning, Amina! Could you tell me which app you use to learn English?\\
— Well, I use Quizlet to memorize vocabulary. What about you?\\
— Really? That's interesting! Personally, I prefer watching videos on YouTube because it helps my listening.\\
— I see your point. Anyway, thanks for sharing. Have a nice day!
\end{exemplebox}

La candidate B salue, pose des questions ouvertes, \cle{réagit} aux réponses (\eng{Really? That's interesting!}), développe ses propres réponses et conclut poliment. Voici la méthode complète.

\begin{methode}[Réussir un jeu de rôle au Bac : les 5 étapes]
\begin{enumerate}
\item \textbf{Greet} — Saluez votre partenaire~: \eng{Hello! / Good morning! How are you today?}
\item \textbf{Ask} — Posez au moins \textbf{trois questions} en utilisant des formules variées~: \eng{Could you tell me...? / What about you? / Do you often...? / Which app do you use to...?}
\item \textbf{Give information} — Donnez des informations \textbf{développées} avec la structure \eng{I use + noun + to + verb}~: \eng{I use my phone to download past papers.}
\item \textbf{React} — Réagissez à chaque réponse de votre partenaire~: \eng{Really? / That's interesting! / I see. / I don't quite agree because...}
\item \textbf{Close} — Concluez poliment~: \eng{Thank you for your time. / It was nice talking to you. / See you soon!}
\end{enumerate}
\end{methode}

\begin{popfigure}
\begin{center}
\begin{tikzpicture}[
  node distance=0.45cm,
  box/.style={draw=popblue, fill=popblueL, rounded corners, thick, align=center, minimum width=4.6cm, minimum height=0.85cm, font=\small},
  arr/.style={-{Stealth[length=2.5mm]}, thick, popdark}
]
\node[box, align=center] (greet){\textbf{1. Greet}\\ \eng{Hello! How are you?}};
\node[box, below=of greet, align=center] (ask){\textbf{2. Ask 3 questions}\\ \eng{Could you tell me...?}};
\node[box, below=of ask, align=center] (give){\textbf{3. Give information}\\ \eng{I use ... to ...}};
\node[box, below=of give, align=center] (react){\textbf{4. React}\\ \eng{Really? That's interesting!}};
\node[box, below=of react, fill=popgreenL, draw=popgreen, align=center] (close){\textbf{5. Close politely}\\ \eng{Thanks for sharing. Goodbye!}};
\draw[arr] (greet) -- (ask);
\draw[arr] (ask) -- (give);
\draw[arr] (give) -- (react);
\draw[arr] (react) -- (close);
\end{tikzpicture}
\end{center}
\legende{Organigramme du jeu de rôle~: les cinq étapes d'un échange réussi.}
\end{popfigure}

\begin{attention}[Ne répondez jamais par un seul mot~!]
L'erreur la plus fréquente des candidats francophones est la réponse minimale~: \eng{Yes.} ou \eng{No.} Cela tue l'interaction et fait perdre des points. Développez toujours en trois temps~: \textbf{réponse + détail + exemple}.\\
\eng{Yes, I do. I use it every evening to revise my lessons.}\\
« Oui. Je l'utilise tous les soirs pour réviser mes leçons. »\\
Autre piège~: ne traduisez pas mot à mot \eng{I am agree} — on dit \eng{I agree} (« je suis d'accord » se dit sans \eng{be}).
\end{attention}

\courssub{Role-play 1 (guided): Comparing favourite apps for learning English}

Travaillez en binôme. L'élève A et l'élève B comparent leurs applications préférées pour apprendre l'anglais. Suivez la trame de six répliques ci-dessous, puis recommencez en changeant les applications et les raisons.

\begin{textebox}[Role-play 1 — My favourite app for learning English]
\textbf{A:} Hello, Brenda! How are you today?\\
\textbf{B:} I'm fine, thanks. And you?\\
\textbf{A:} Very well. Could you tell me which app you use to learn English?\\
\textbf{B:} Well, I use Duolingo to practise grammar every morning before school. What about you?\\
\textbf{A:} Really? That's interesting! Personally, I use YouTube to watch English lessons because it improves my listening.\\
\textbf{B:} I see. In my opinion, using two apps is even better. Anyway, thank you for sharing, and see you in class!\\
\textbf{A:} Thank you too. Have a nice day!
\end{textebox}

\begin{center}
\begin{tabularx}{\linewidth}{|X|X|}
\hline
\rowcolor{popblueL} \textbf{English} & \textbf{Français} \\
\hline
“Could you tell me which app you use?” & « Pourriez-vous me dire quelle application vous utilisez ? » \\
\hline
“I use Duolingo to practise grammar.” & « J'utilise Duolingo pour m'entraîner en grammaire. » \\
\hline
“What about you?” & « Et toi / vous ? » \\
\hline
“It improves my listening.” & « Cela améliore ma compréhension orale. » \\
\hline
“In my opinion, using two apps is even better.” & « À mon avis, utiliser deux applications, c'est encore mieux. » \\
\hline
“Thank you for sharing.” & « Merci d'avoir partagé (ces informations). » \\
\hline
\end{tabularx}
\end{center}

\courssub{Role-play 2 (semi-guided): Interviewing a teacher about ICTs in class}

Cette fois, seules les questions sont fournies~; les réponses sont libres. L'élève A joue le rôle d'un journaliste du journal du lycée~; l'élève B joue le rôle d'un enseignant du lycée de Buea qui utilise les TIC en classe.

\begin{experience}[Interview : ICTs in the classroom]
L'élève A pose au moins cinq des questions suivantes, dans l'ordre de son choix, et réagit après chaque réponse (\eng{Really? / That's surprising! / I see.})~:
\begin{enumerate}
\item \eng{Good morning, Sir. May I ask you a few questions about ICTs in your classroom?}
\item \eng{How often do you use a computer or a projector in your lessons?}
\item \eng{Which digital tools do your students use to do their homework?}
\item \eng{Do you think ICTs make learning easier? Why?}
\item \eng{What problems do you face when you use technology in class?}
\item \eng{What advice can you give to students who want to use ICTs to improve their English?}
\item \eng{Thank you very much for your time, Sir. Have a nice day!}
\end{enumerate}
L'élève B répond librement, en développant chaque réponse avec \eng{I use ... to ...} ou \eng{In my opinion, ...}. Puis inversez les rôles.
\end{experience}

\begin{exoresolu}[Transformer une réponse minimale en réponse développée]
\textbf{Énoncé.} Pendant l'interview, l'enseignant répond à la question \eng{Do you think ICTs make learning easier?} simplement par~: \eng{Yes.} Réécrivez sa réponse en trois temps (réponse + détail + exemple). \tcblower
\textbf{Solution.} On applique la structure~: réponse (\eng{Yes, I do.}) + opinion/détail (\eng{In my opinion, ICTs make learning easier.}) + exemple concret avec \eng{to + verb}.\\
\eng{Yes, I do. In my opinion, ICTs make learning easier. For example, my students use their phones to watch English videos and to check difficult words in an online dictionary.}\\
« Oui. À mon avis, les TIC rendent l'apprentissage plus facile. Par exemple, mes élèves utilisent leur téléphone pour regarder des vidéos en anglais et pour vérifier les mots difficiles dans un dictionnaire en ligne. »
\end{exoresolu}

\courssub{Role-play 3 (free practice): Debate — Are ICTs good or bad for students?}

Dernière étape~: un débat libre en binôme. L'élève A défend les avantages des TIC, l'élève B défend les inconvénients. \textbf{Obligation}~: chacun doit utiliser \textbf{au moins deux formules d'accord ou de désaccord} du tableau ci-dessous.

\begin{exemplebox}[Formules d'opinion, d'accord et de désaccord (traduites)]
\eng{In my opinion, ICTs make learning easier.} « À mon avis, les TIC rendent l'apprentissage plus facile. »\\
\eng{I completely agree with you.} « Je suis tout à fait d'accord avec toi. »\\
\eng{That's true, but we must also consider the cost of data.} « C'est vrai, mais il faut aussi considérer le coût des données. »\\
\eng{I don't quite agree because many students waste time on social media.} « Je ne suis pas tout à fait d'accord parce que beaucoup d'élèves perdent du temps sur les réseaux sociaux. »\\
\eng{You have a point, but I think ICTs help us to work faster.} « Tu marques un point, mais je pense que les TIC nous aident à travailler plus vite. »
\end{exemplebox}

\begin{center}
\begin{tabularx}{\linewidth}{|X|X|X|}
\hline
\rowcolor{popblueL} \textbf{Fonction} & \textbf{English} & \textbf{Français} \\
\hline
Donner son opinion & “In my opinion, ... / I think (that) ...” & « À mon avis... / Je pense que... » \\
\hline
Être d'accord & “I (completely) agree. / That's true.” & « Je suis (tout à fait) d'accord. / C'est vrai. » \\
\hline
Être partiellement d'accord & “You have a point, but...” & « Tu marques un point, mais... » \\
\hline
Ne pas être d'accord & “I don't quite agree because... / I'm afraid I disagree.” & « Je ne suis pas tout à fait d'accord parce que... / Je crains de ne pas être d'accord. » \\
\hline
Conclure le débat & “To sum up, ... / Anyway, thanks for the discussion.” & « Pour résumer... / En tout cas, merci pour cette discussion. » \\
\hline
\end{tabularx}
\end{center}

\begin{experience}[Débat en binôme : Are ICTs good or bad for students?]
\begin{enumerate}
\item Préparez deux arguments chacun (une minute de préparation).
\item L'élève A commence~: \eng{In my opinion, ICTs are good for students because...}
\item Échangez au moins \textbf{quatre répliques chacun}, en réagissant à chaque argument de votre partenaire.
\item Concluez ensemble~: \eng{To sum up, we agree that ICTs are useful if we use them wisely.}
\end{enumerate}
\end{experience}

\courssub{Self-assessment grid}

Après chaque jeu de rôle, évaluez-vous honnêtement avec la grille suivante. Cochez \eng{Yes} ou \eng{Not yet} et fixez-vous un objectif pour la prochaine prise de parole.

\begin{center}
\begin{tabularx}{\linewidth}{|X|l|l|}
\hline
\rowcolor{popblueL} \textbf{Question d'auto-évaluation} & \textbf{Yes} & \textbf{Not yet} \\
\hline
“Did I greet my partner politely?” (Ai-je salué poliment ?) & & \\
\hline
“Did I ask at least 3 questions?” (Ai-je posé au moins 3 questions ?) & & \\
\hline
“Did I use ‘I use ... to ...’ at least once?” (Ai-je utilisé la structure « I use ... to ... » ?) & & \\
\hline
“Did I react to my partner's answers?” (Ai-je réagi aux réponses ?) & & \\
\hline
“Did I use 2 expressions of agreement or disagreement?” (Ai-je utilisé 2 formules d'accord/désaccord ?) & & \\
\hline
“Did I close the conversation politely?” (Ai-je conclu poliment ?) & & \\
\hline
\end{tabularx}
\end{center}

\begin{aretenir}[Les critères de réussite à l'oral]
Un bon échange oral repose sur quatre piliers~: \textbf{saluer}, \textbf{questionner} (au moins 3 questions variées), \textbf{développer} ses réponses (jamais un seul mot), \textbf{conclure poliment}. La prononciation et l'intonation travaillées dans la section précédente donnent à votre échange son naturel.
\end{aretenir}

\courssub{Class correction: common mistakes heard during the role-plays}

Pendant les jeux de rôle, le professeur note les erreurs entendues sans interrompre, puis la classe les corrige collectivement au tableau. Voici les erreurs les plus fréquentes des francophones dans cette situation de communication.

\begin{center}
\begin{tabularx}{\linewidth}{|X|X|X|}
\hline
\rowcolor{popblueL} \textbf{Erreur entendue} & \textbf{Correction} & \textbf{Explication} \\
\hline
“I am agree.” & “I agree.” & \eng{Agree} est un verbe~: pas de \eng{be}. \\
\hline
“I use my phone for learn English.” & “I use my phone to learn English.” & Après \eng{use ...}, le but s'exprime avec \eng{to + verb}. \\
\hline
“Do you think that ICTs are good? — Yes.” & “Yes, I do. I think they help us to...” & Jamais de réponse d'un seul mot~: développez. \\
\hline
“Which app do you use for to study?” & “Which app do you use to study?” & On ne dit jamais \eng{for to}. \\
\hline
“I don't quite am agree.” & “I don't quite agree.” & Une seule négation, sans \eng{be}. \\
\hline
“Thank you for your answer. Bye!” (intonation plate) & “Thank you for your answer. Goodbye!” (intonation descendante) & La conclusion se dit avec une intonation descendante et chaleureuse. \\
\hline
\end{tabularx}
\end{center}

\begin{exercice}[À vous de jouer]
En binôme, rejouez le rôle-play 1 en remplaçant les applications par deux outils que vous utilisez réellement (WhatsApp, Google, un dictionnaire en ligne, la radio en anglais). Votre dialogue doit contenir~: une salutation, trois questions, deux réactions, une formule d'opinion et une conclusion polie. Puis remplissez la grille d'auto-évaluation.
\end{exercice}

\begin{corrige}[Exercice — éléments attendus]
Il n'existe pas de réponse unique, mais un dialogue réussi doit respecter les cinq étapes de l'organigramme. Exemple de production attendue~:\\
\textbf{A:} Hi, Junior! How are you?\\
\textbf{B:} Fine, thanks. Could you tell me which tool you use to learn English?\\
\textbf{A:} Well, I use WhatsApp to chat in English with my cousin in Bamenda. What about you?\\
\textbf{B:} Really? That's clever! I use Google to check the meaning of new words.\\
\textbf{A:} In my opinion, chatting is the best way to improve.\\
\textbf{B:} I don't quite agree, because videos also train our ears.\\
\textbf{A:} You have a point. Anyway, thanks for sharing. See you tomorrow!\\
Vérifiez~: 3 questions ou plus, \eng{I use ... to ...}, 2 formules d'accord/désaccord, salutation et conclusion.
\end{corrige}

\begin{savaistu}[À l'épreuve orale du Baccalauréat]
À l'oral du Bac, l'examinateur valorise la \cle{fluidité} et l'\cle{interaction} bien plus que la perfection grammaticale. Un candidat qui hésite, corrige une petite erreur et continue à parler naturellement obtient une meilleure note qu'un candidat qui récite un texte mémorisé sans écouter son interlocuteur. Les examinateurs camerounais apprécient particulièrement les candidats qui réagissent (\eng{Really? That's interesting!}) et qui concluent poliment~: ce sont des marqueurs de compétence communicative authentique.
\end{savaistu}

\newpage

\coursec{Activité d'intégration}

\courssub{Objectifs de l'activité}

À l'issue de cette activité d'intégration, l'élève sera capable de :

\begin{itemize}
\item mobiliser, dans un échange oral authentique, l'ensemble des \cle{formules fonctionnelles} étudiées dans la leçon : demander une information (\eng{Which apps do you use...?}), donner une information (\eng{I use... to...}), relancer (\eng{Really? Could you tell us more?}) et conclure (\eng{Thank you for joining us today});
\item employer correctement le \cle{vocabulaire des TIC} (apps, platforms, devices, online resources) et la structure \eng{I use + nom + to + base verbale};
\item exprimer un \cle{avantage} et un \cle{inconvénient} des TIC pour l'apprentissage et le travail;
\item réagir aux propos d'un interlocuteur (\eng{I agree.} / \eng{I see your point, but...});
\item soigner sa \cle{prononciation}, son \cle{accentuation} et son \cle{intonation} (montante pour les questions yes/no, descendante pour les questions en \eng{wh-});
\item évaluer la production orale de ses camarades à l'aide d'une grille critériée.
\end{itemize}

\courssub{Situation}

\begin{textebox}[CRTV Radio — “The Youth Hour”, Yaoundé]
\textbf{Announcement heard on CRTV Radio, Monday 7.30 a.m.}

“Good morning, Cameroon! This is \emph{The Youth Hour} on CRTV Radio. This week, our topic is: \emph{ICTs in my school}. How do students and teachers use phones, computers and the internet to learn and to work better? We are looking for secondary schools ready to send a team — one presenter, one student and one teacher — for a three-minute live interview. The best interviews will be broadcast on Friday and published on our website. Is your school ready? Get your microphones ready... and speak English!”
\end{textebox}

\begin{definition}[Glossaire utile]
\begin{itemize}
\item \cle{presenter} (n.) : animateur / animatrice d'une émission.
\item \cle{live interview} (n.) : interview en direct.
\item \cle{to broadcast} (v.) : diffuser (à la radio ou à la télévision).
\item \cle{ready} (adj.) : prêt. Prononciation : « ré-di ».
\end{itemize}
\end{definition}

Votre classe de Terminale A du Lycée de Nkol-Eton (Yaoundé) a été sélectionnée pour participer à l'émission \emph{The Youth Hour}. Par groupes de trois, vous allez préparer et jouer une \cle{interview radio de trois minutes}, en direct devant la classe, sur le thème : \emph{“ICTs in my school — using technology to learn and to work better.”} Les autres groupes joueront le rôle du jury de la CRTV et évalueront votre prestation avec la grille vue en classe.

\courssub{Consignes}

\textbf{Tâche 1 — Comprendre la situation (5 min).} Lisez l'annonce de la CRTV à voix haute à tour de rôle. Répondez oralement en anglais : \eng{What is the topic of the programme? Who must be in the team? How long is the interview?} Soulignez dans le texte deux mots du champ lexical des médias.

\textbf{Tâche 2 — Distribuer les rôles (2 min).} Dans chaque groupe de trois, choisissez : un \cle{presenter} (l'animateur), un \cle{student} (un élève de Terminale) et un \cle{teacher} (un enseignant d'anglais ou d'informatique). Chacun note son rôle en tête de sa feuille de préparation.

\textbf{Tâche 3 — Préparer les questions de l'animateur (10 min).} L'animateur rédige \textbf{au moins quatre questions} variées sur l'utilisation des TIC pour apprendre et travailler. Il doit utiliser au moins trois formes différentes :
\begin{itemize}
\item une question en \eng{wh-} : \eng{Which apps do you use to do your homework?}
\item une question en \eng{How...?} : \eng{How do you use the internet to prepare your lessons?}
\item une question d'opinion : \eng{Do you think smartphones help students in class?}
\end{itemize}

\textbf{Tâche 4 — Préparer les réponses des invités (10 min).} L'élève et l'enseignant préparent leurs réponses. Chaque réponse doit contenir :
\begin{enumerate}
\item la structure \eng{I use + outil + to + verbe} : \eng{I use WhatsApp to share exercises with my classmates.};
\item \textbf{un avantage} : \eng{It is fast and free.};
\item \textbf{un inconvénient} : \eng{However, the connection is sometimes slow.};
\item \textbf{une réaction} aux propos de l'autre invité : \eng{I agree with the teacher.} ou \eng{Really? In my case, it is different.}
\end{enumerate}

\textbf{Tâche 5 — Préparer l'ouverture et la conclusion (5 min).} L'animateur rédige une phrase d'ouverture (\eng{Good morning, dear listeners! Welcome to The Youth Hour...}) et une phrase de conclusion (\eng{Thank you very much for joining us today. That was The Youth Hour on CRTV Radio!}).

\textbf{Tâche 6 — Répéter en soignant la prononciation (5 min).} Répétez l'interview dans votre groupe. Attention à l'\cle{intonation} : voix qui \emph{monte} à la fin des questions yes/no (\eng{Do you think smartphones help students?} « smar-te-phone elp stiou-dents » ↗), voix qui \emph{descend} à la fin des questions en \eng{wh-} (\eng{Which apps do you use?} ↘). Accentuez le mot important de chaque phrase.

\textbf{Tâche 7 — Jouer l'interview en direct (3 min par groupe).} Chaque groupe passe devant la classe. Les autres groupes remplissent la grille d'évaluation : nombre de questions posées (minimum 4), formules fonctionnelles utilisées, structure \eng{I use... to...}, avantage et inconvénient exprimés, réactions (\eng{I agree} / \eng{Really?}), intonation, conclusion de l'animateur.

\textbf{Tâche 8 — Enregistrer les meilleures interviews.} Le groupe ayant obtenu la meilleure moyenne est enregistré au téléphone ; l'enregistrement sera écouté en classe de révision.

\courssub{Ressources à mobiliser}

\begin{propriete}[Formules fonctionnelles de la leçon]
\begin{center}
\begin{tabularx}{\linewidth}{|l|X|}\hline
\rowcolor{popblueL} \textbf{Fonction} & \textbf{Formules} \\ \hline
Ouvrir & \eng{Good morning, dear listeners! Welcome to...} \\ \hline
Demander & \eng{Which apps do you use...? How do you...? Do you think...? Could you tell us more about...?} \\ \hline
Informer & \eng{I use... to... / In my school, we... / For example,...} \\ \hline
Avantage / inconvénient & \eng{It is fast and free. / However, it can be expensive. / The problem is that...} \\ \hline
Réagir & \eng{I agree. / I see your point, but... / Really? That's interesting!} \\ \hline
Conclure & \eng{Thank you for joining us today. / That's all for The Youth Hour!} \\ \hline
\end{tabularx}
\end{center}
\end{propriete}

\begin{definition}[Vocabulaire des TIC à réemployer]
\cle{apps} (applications), \cle{smartphone} (« smarte-phone »), \cle{laptop} (ordinateur portable), \cle{platform} (plateforme), \cle{online resources} (ressources en ligne), \cle{to search for information} (chercher des informations), \cle{to share documents} (partager des documents), \cle{to prepare lessons} (préparer les leçons), \cle{connection} (connexion), \cle{data bundle} (forfait internet).
\end{definition}

\begin{attention}[Erreurs fréquentes à éviter]
\begin{itemize}
\item Ne dites pas \eng{I use my phone for to learn.} La structure correcte est \eng{I use my phone \textbf{to} learn.} ou \eng{I use my phone \textbf{for learning}.} — jamais les deux ensemble.
\item \eng{Informations} est \textbf{invariable} en anglais : \eng{I search for information} (sans « s »).
\item Prononcez \eng{the} « ze » devant consonne et « zi » devant voyelle : \eng{the internet} « zi in-ter-nette ».
\item N'oubliez pas le \textbf{-s} de la 3\textsuperscript{e} personne : \eng{My teacher use\textbf{s} a laptop.}
\end{itemize}
\end{attention}

\courssub{Proposition de production et corrigé commenté}

\begin{exoresolu}[Interview modèle — “ICTs in my school”]
Énoncé : rédiger puis jouer une interview radio de trois minutes respectant toutes les consignes (4 questions minimum, structure \eng{I use... to...}, avantage, inconvénient, réactions, conclusion).
\tcblower
\textbf{Production modèle (script) :}

\textbf{Presenter:} Good morning, dear listeners! Welcome to \emph{The Youth Hour} on CRTV Radio. Today, our topic is “ICTs in my school”. With me are Amina, a student in Terminale A, and Mr Ndi, her English teacher. Amina, which apps do you use to learn English?

\textbf{Student:} Well, I use WhatsApp to share exercises with my classmates, and I use an online dictionary to check new words. It is fast and free. However, my data bundle finishes very quickly!

\textbf{Presenter:} Really? That's interesting! And you, Mr Ndi, how do you use ICTs to prepare your lessons?

\textbf{Teacher:} I use my laptop to download texts and videos for my students. For example, last week I showed a short video about Cameroonian wildlife. The students loved it. But the connection at school is sometimes very slow.

\textbf{Student:} I agree with Mr Ndi. When the connection is slow, we cannot watch the videos.

\textbf{Presenter:} I see. Do you think smartphones help students in class, Amina?

\textbf{Student:} Yes and no. They help us to search for information quickly, but some students use them to play games during the lesson.

\textbf{Teacher:} I see your point, Amina, but with clear rules, smartphones can be very useful.

\textbf{Presenter:} Thank you very much, Amina and Mr Ndi, for joining us today. That was \emph{The Youth Hour} on CRTV Radio. Goodbye, dear listeners!

\textbf{Commentaire des choix de langue :}
\begin{itemize}
\item \textbf{Ouverture et conclusion} : l'animateur utilise les formules rituelles de la radio (\eng{Welcome to...}, \eng{Thank you for joining us}), ce qui donne un cadre authentique à l'échange.
\item \textbf{Questions variées} : quatre questions, trois formes différentes (\eng{Which apps...?}, \eng{How do you...?}, \eng{Do you think...?}), avec intonation descendante pour les deux premières et montante pour la troisième.
\item \textbf{Structure cible} : \eng{I use WhatsApp \textbf{to} share}, \eng{I use my laptop \textbf{to} download} — infinitif de but correctement employé, sans « for to ».
\item \textbf{Avantages et inconvénients} : \eng{fast and free} / \eng{data bundle finishes quickly} / \eng{connection is slow} — chaque invité équilibre son propos avec \eng{however} et \eng{but}.
\item \textbf{Réactions} : \eng{Really? That's interesting!} (relance de l'animateur), \eng{I agree with Mr Ndi}, \eng{I see your point, but...} (désaccord poli de l'enseignant).
\item \textbf{Grammaire} : \eng{information} sans « s », \eng{students use\textbf{d}} évité au profit du présent simple pour les habitudes, \eng{some students use} avec verbe au pluriel.
\end{itemize}
\end{exoresolu}

\courssub{Bilan de l'activité}

Cette activité d'intégration a permis de réinvestir, dans une situation de communication réaliste et motivante (une émission de la CRTV), l'ensemble des acquis de la leçon : les \cle{formules fonctionnelles} pour ouvrir, questionner, informer, relancer et conclure ; la structure \eng{I use... to...} ; le lexique des TIC ; l'expression nuancée d'un avantage et d'un inconvénient ; ainsi que la \cle{prononciation} et l'\cle{intonation} propres aux questions anglaises. La grille d'évaluation remplie par les pairs a développé l'écoute active et l'esprit critique des élèves. Pour aller plus loin, l'enregistrement de la meilleure interview pourra être réécouté en leçon de révision, et chaque élève pourra rédiger à l'écrit un court compte rendu de son interview (\eng{In our interview, the presenter asked four questions...}), préparant ainsi la transition vers la compétence \emph{writing}.

\newpage

\coursec{Méthodes et exercices résolus}

\coursec{Lesson 52 — Speaking: Exchanges information on using ICTs to enhance work or learning}

\courssub{Warm-up: ICTs in our daily lives}

\begin{experience}[Brainstorming et vocabulaire de base]
En binôme, regardez l'image projetée (ou décrite) : un élève à Buea avec une tablette, un enseignant à Yaoundé qui projette un cours, un employé de microfinance à Douala sur son smartphone, un cybercafé à Bamenda.
\begin{enumerate}
\item Citez \cle{cinq outils TIC} (ICT tools) que vous voyez ou imaginez (smartphone, laptop, tablet, projector, router, USB drive, printer).
\item Citez \cle{trois usages} pour les études (download past papers, watch video lessons, join WhatsApp study group, type assignments, check results online).
\item Citez \cle{trois usages} pour le travail (send emails, use spreadsheet for accounts, video conference with head office, mobile money transactions, digital CV).
\end{enumerate}
\emph{Partagez vos idées à l'oral. Le professeur note le vocabulaire clé au tableau : hardware, software, app, platform, bandwidth, connection, download, upload, login, password, screenshot, spreadsheet, video call.}
\end{experience}

\begin{definition}[Vocabulaire clé : ICTs at school and at work]
\begin{center}
\begin{tabularx}{\linewidth}{|X|X|X|}
\hline
\rowcolor{popblueL} \textbf{English} & \textbf{Français} & \textbf{Exemple / Collocation} \\
\hline
ICTs (Information and Communication Technologies) & TIC (Technologies de l'Information et de la Communication) & \eng{ICTs enhance learning.} \\
\hline
smartphone / mobile phone & téléphone intelligent / portable & \eng{I use my smartphone to revise.} \\
\hline
laptop / notebook & ordinateur portable & \eng{The school provides laptops.} \\
\hline
tablet & tablette tactile & \eng{She reads e-books on her tablet.} \\
\hline
cybercafé / internet café & cybercafé & \eng{Many students go to the cybercafé in Bamenda.} \\
\hline
broadband / high-speed connection & connexion haut débit & \eng{We need a stable broadband connection.} \\
\hline
platform / learning platform & plateforme (d'apprentissage) & \eng{The GCE Board uses an online platform.} \\
\hline
spreadsheet (Excel) & tableur & \eng{He records transactions in a spreadsheet.} \\
\hline
video conference / video call & visioconférence / appel vidéo & \eng{We had a video call with Douala.} \\
\hline
mobile money (MTN MoMo, Orange Money) & argent mobile & \eng{Workers receive salaries via mobile money.} \\
\hline
to download / to upload & télécharger / envoyer (fichier) & \eng{Download the past papers from the website.} \\
\hline
to log in / to log out & se connecter / se déconnecter & \eng{Log in with your student ID.} \\
\hline
assignment / homework & devoir / travail à faire & \eng{Submit your assignment online.} \\
\hline
past papers & annales (sujets d'examen) & \eng{Practise with past papers.} \\
\hline
e-learning / online learning & apprentissage en ligne & \eng{E-learning is growing in Cameroon.} \\
\hline
digital skills / computer literacy & compétences numériques / culture informatique & \eng{Digital skills are essential today.} \\
\hline
\end{tabularx}
\end{center}
\emph{Note : \eng{homework}, \eng{equipment}, \eng{information}, \eng{software} sont \cle{indénombrables} (pas de \eng{-s} au pluriel).}
\end{definition}

\courssub{Model dialogue 1: ICTs for learning}

\begin{textebox}[Dialogue modèle — Two students at GBHS Yaoundé]
\textbf{Divine:} Good morning, Paul. Do you have a minute? I'm doing a survey for the school magazine.

\textbf{Paul:} Good morning, Divine. Sure, go ahead.

\textbf{Divine:} \eng{How do you use ICTs for your studies?}

\textbf{Paul:} I mainly use my smartphone. I download past papers from the GCE Board website and I watch video lessons on YouTube.

\textbf{Divine:} \eng{How often do you do that?}

\textbf{Paul:} Almost every evening, especially when exams are coming. I also join a WhatsApp group where we share notes.

\textbf{Divine:} \eng{That's interesting. Could you tell me more about the WhatsApp group?}

\textbf{Paul:} Yes. Our class representative posts the homework there. If someone doesn't understand, we explain in voice notes.

\textbf{Divine:} \eng{Why do you prefer the phone to a computer?}

\textbf{Paul:} Because it's cheaper and I always have it with me. Internet bundles are affordable.

\textbf{Divine:} \eng{Thank you very much for your time, Paul.}

\textbf{Paul:} You're welcome. Good luck with the magazine!
\end{textebox}

\begin{exoresolu}[Analyse du dialogue modèle 1]
Énoncé — Identifiez dans le dialogue : (a) la formule d'ouverture, (b) trois questions (type), (c) une relance, (d) la formule de conclusion, (e) deux connecteurs de développement.
\tcblower
Solution détaillée :
\begin{itemize}
\item \textbf{(a) Ouverture :} \eng{Good morning, Paul. Do you have a minute? I'm doing a survey\ldots} (politesse + raison).
\item \textbf{(b) Questions :}
      \begin{enumerate}
      \item \eng{How do you use ICTs for your studies?} (question \eng{How} + présent simple \eng{do you use} — manière).
      \item \eng{How often do you do that?} (question \eng{How often} + auxiliaire \eng{do} — fréquence).
      \item \eng{Why do you prefer the phone to a computer?} (question \eng{Why} + auxiliaire \eng{do} — raison).
      \end{enumerate}
\item \textbf{(c) Relance :} \eng{Could you tell me more about the WhatsApp group?} (question indirecte, ordre affirmatif \eng{the group is}).
\item \textbf{(d) Conclusion :} \eng{Thank you very much for your time, Paul.} (remerciement + nom).
\item \textbf{(e) Connecteurs :} \eng{because} (cause), \eng{also} (addition), \eng{especially} (précision).
\end{itemize}
\emph{Observation :} Les réponses de Paul sont \cle{développées} (phrases complètes, raisons, exemples). C'est l'attendu au Bac oral.
\end{exoresolu}

\courssub{Functional language: asking, informing, follow-up, closing}

\begin{methode}[Mener un échange d'informations sur les TIC — 5 étapes]
Pour réussir un dialogue oral sur l'utilisation des TIC (\eng{ICTs}) pour le travail ou les études, suis cette démarche :
\begin{enumerate}
\item \textbf{Ouvrir le dialogue} avec une formule polie adaptée au contexte :
      \begin{itemize}
      \item Formel (inconnu, supérieur) : \eng{Excuse me, Sir/Madam. May I ask you a few questions about ICTs in your company?}
      \item Semi-formel (camarade, collègue) : \eng{Hi [Prénom], do you have a minute? I'd like to ask you about\ldots}
      \end{itemize}
\item \textbf{Demander une information} avec une question claire (mot interrogatif + auxiliaire + sujet + verbe) :
      \begin{itemize}
      \item \eng{How do you use the Internet for your studies?} (manière)
      \item \eng{Which apps do you find most useful?} (choix limité)
      \item \eng{How often do you connect to the school platform?} (fréquence)
      \item \eng{Why are ICTs important in your work?} (raison)
      \end{itemize}
\item \textbf{Donner une information} avec des verbes précis et des connecteurs :
      \begin{itemize}
      \item \eng{I use my smartphone to download past papers.} (but : \eng{to} + infinitif)
      \item \eng{Our school uses WhatsApp groups \textbf{so that} students can share assignments.} (but : \eng{so that} + can)
      \item \eng{ICTs help us \textbf{to} work faster \textbf{and} avoid mistakes.} (verbe + infinitif + coordination)
      \end{itemize}
\item \textbf{Relancer} pour obtenir des précisions (écoute active) :
      \begin{itemize}
      \item Question de suivi : \eng{Could you tell me more about that?} / \eng{What exactly do you mean by ``e-learning''?} / \eng{Can you give me an example?}
      \item Reformulation : \eng{So, if I understand correctly, you study with videos every night?}
      \item Marqueurs d'intérêt : \eng{Really?} / \eng{I see.} / \eng{That's interesting!} / \eng{Oh, I didn't know that.}
      \end{itemize}
\item \textbf{Conclure} poliment :
      \begin{itemize}
      \item \eng{Thank you very much for the information.}
      \item \eng{That was really helpful. Goodbye!}
      \item \eng{Thanks a lot for sharing your experience. See you!}
      \end{itemize}
\end{enumerate}
\emph{Réflexe Bac :} Prépare toujours au moins \cle{trois questions} variées (\eng{What/Which, How often, Why}) et \cle{deux réponses développées} (avec \eng{because}, \eng{so that}, \eng{in order to}, \eng{for example}) avant de commencer.
\end{methode}

\begin{propriete}[Rappel grammatical : Formation des questions en anglais]
\begin{center}
\begin{tabularx}{\linewidth}{|X|X|X|}
\hline
\rowcolor{popblueL} \textbf{Type de question} & \textbf{Structure} & \textbf{Exemples (ICTs)} \\
\hline
Question \textbf{sujet} (Who/What/Which + verbe) & \textbf{Mot interrogatif (sujet) + verbe conjugué} & \eng{Who uses the computer?} — \eng{What works well?} \\
\hline
Question \textbf{objet} (Wh- + auxiliaire) & \textbf{Mot interrogatif + auxiliaire (do/does/did/can/will) + sujet + verbe base} & \eng{What do you use?} — \eng{How often does she connect?} — \eng{Why did they buy it?} \\
\hline
Question \textbf{fermée} (Yes/No) & \textbf{Auxiliaire + sujet + verbe base} & \eng{Do you have a laptop?} — \eng{Can I connect here?} — \eng{Did you send the email?} \\
\hline
Question \textbf{indirecte} (politesse) & \textbf{Could you tell me / Do you know + sujet + verbe conjugué (SANS auxiliaire do/does/did)} & \eng{Could you tell me where the router is?} — \eng{Do you know how much it costs?} \\
\hline
\end{tabularx}
\end{center}
\emph{Règle d'or :} Dans une question indirecte, l'ordre redevient \textbf{affirmatif} (sujet + verbe) et l'auxiliaire \eng{do/does/did} \textbf{disparaît}. Au présent simple 3\textsuperscript{e} pers. sg., le \textbf{-s revient sur le verbe principal} : \eng{How much does it cost?} $\rightarrow$ \eng{Do you know how much it \textbf{costs}?}
\end{propriete}

\begin{exoresolu}[Questions et réponses — Entraînement ciblé]
Énoncé — (a) Write the correct question for each answer. (b) Answer question 3 personally in a full sentence.

1. A: \ldots\ldots\ldots\ldots\ldots\ldots\ldots\ldots\ldots\ldots\ldots\ldots\ldots\ldots\ldots\ldots\ldots\ldots\ldots ? B: I use my tablet \textbf{to read e-books}. (purpose)

2. A: \ldots\ldots\ldots\ldots\ldots\ldots\ldots\ldots\ldots\ldots\ldots\ldots\ldots\ldots\ldots\ldots\ldots\ldots\ldots ? B: The secretary \textbf{connects to the platform twice a week}. (frequency, 3\textsuperscript{e} pers. sg.)

3. A: \ldots\ldots\ldots\ldots\ldots\ldots\ldots\ldots\ldots\ldots\ldots\ldots\ldots\ldots\ldots\ldots\ldots\ldots\ldots ? B: \textbf{Because} they allow fast communication and safe storage of data. (reason — ICTs at work)
\tcblower
Solution détaillée :
\begin{itemize}
\item \textbf{1.} La réponse indique un but (\eng{to read}). Question : \eng{Why do you use your tablet?} ou \eng{What do you use your tablet for?} (préposition \eng{for} à la fin). \textbf{Attention :} \eng{For what do you use\ldots?} est trop formel / vieilli.
\item \textbf{2.} \eng{Twice a week} = fréquence $\rightarrow$ \eng{How often}. Sujet = \eng{The secretary} (3\textsuperscript{e} pers. sg.) $\rightarrow$ auxiliaire \eng{does}, verbe \textbf{sans -s} : \eng{How often \textbf{does} the secretary \textbf{connect} to the platform?} Piège : ne pas écrire \eng{connects} après \eng{does}.
\item \textbf{3.} La réponse commence par \eng{Because} $\rightarrow$ Question \eng{Why}. Sujet pluriel \eng{they} $\rightarrow$ auxiliaire \eng{do} (ou \eng{are} si adjectif) : \eng{Why \textbf{are} ICTs important at work?} ou \eng{Why \textbf{do} companies \textbf{use} ICTs?}
\item \textbf{(b) Réponse personnelle (exemple) :} \eng{ICTs are crucial at work because they enable fast communication through email and video calls, and they allow secure data storage in the cloud. For instance, in my father's company in Douala, accountants use spreadsheets to manage finances instead of paper ledgers.}
\end{itemize}
\end{exoresolu}

\courssub{Model dialogue 2: ICTs at work}

\begin{textebox}[Dialogue modèle — Journalist interviews a microfinance worker in Buea]
\textbf{Journalist:} Good morning, Mr. Atanga. I'm from \eng{The Buea Times}. May I ask you about ICTs in your daily work?

\textbf{Mr. Atanga:} Good morning. Yes, certainly. What would you like to know?

\textbf{Journalist:} \eng{What ICT tools do you use at the agency?}

\textbf{Mr. Atanga:} We use desktop computers for the database, and smartphones for mobile money transactions with clients. We also have a POS terminal for card payments.

\textbf{Journalist:} \eng{How often do you rely on the Internet?}

\textbf{Mr. Atanga:} Constantly. The system is cloud-based. If the connection drops, we cannot process loans or savings withdrawals.

\textbf{Journalist:} \eng{That sounds critical. Could you explain how mobile money helps your clients?}

\textbf{Mr. Atanga:} Of course. Many clients live in remote villages. They deposit or withdraw cash via agents without travelling to Buea. It saves time and transport costs.

\textbf{Journalist:} \eng{Why would you say ICTs enhance your work overall?}

\textbf{Mr. Atanga:} Because they increase speed, reduce human error in calculations, and build trust through instant SMS receipts. Ten years ago, everything was manual ledgers.

\textbf{Journalist:} \eng{Thank you very much for your insights, Mr. Atanga.}

\textbf{Mr. Atanga:} My pleasure. Have a nice day.
\end{textebox}

\courssub{Pronunciation, stress and intonation}

\begin{methode}[Bien prononcer le vocabulaire des TIC — Accentuation et intonation]
\begin{enumerate}
\item \textbf{Accentue la bonne syllabe} (dictionnaire britannique) :
      \begin{itemize}
      \item \eng{comPUter} (« kom-pio-\textbf{TEUR} ») — accent sur 2\textsuperscript{e} syllabe.
      \item \eng{INternet} (« \textbf{IN}-teu-net ») — accent sur 1\textsuperscript{re} syllabe.
      \item \eng{inforMAtion} (« in-feur-\textbf{MÉ}-cheune ») — accent sur syllabe avant \eng{-tion}.
      \item \eng{deVElop} (« di-\textbf{VÉ}-leup ») — accent sur 2\textsuperscript{e} syllabe.
      \item \eng{techNOlogy} (« tek-\textbf{NO}-lo-dji ») — accent sur 2\textsuperscript{e} syllabe (avant \eng{-logy}).
      \item \eng{moBILE} (« mo-\textsc{bail} »), \eng{laPTOP} (« \textsc{lap}-top »), \eng{softWARE} (« soft-\textsc{weur} »), \eng{hardWARE} (« \textsc{hard}-weur »).
      \end{itemize}
\item \textbf{Soigne les sons pièges} :
      \begin{itemize}
      \item \eng{th} (\eng{the, this, that, with}) : langue entre les dents $\rightarrow$ « \textbf{zeu}, \textbf{zis}, \textbf{zat}, \textbf{wiz} » (jamais « dé, di »).
      \item \eng{work} /w$\backslash$ːrk/ $\rightarrow$ « \textbf{weurk} » (son \eng{ur} comme \eng{bird}) $\neq$ \eng{walk} /w$\backslash$ːk/ $\rightarrow$ « \textbf{wok} ».
      \item \eng{ICTs} : se dit lettre par lettre « \textbf{aï-si-tis} » (pas « ictes »).
      \item \eng{Wi-Fi} : « \textbf{waï-faï} ».
      \end{itemize}
\item \textbf{Intonation montante ↗} (voix qui monte à la fin) : questions fermées (réponse Yes/No).
      \eng{Do you have a smartphone?} ↗ \quad \eng{Can I connect here?} ↗
\item \textbf{Intonation descendante ↘} (voix qui tombe) : questions ouvertes (Wh-), affirmations, ordres.
      \eng{What do you use it for?} ↘ \quad \eng{I use it for learning.} ↘ \quad \eng{Please log in.} ↘
\item \textbf{Accent de phrase (focus)} : le mot porteur de l'information nouvelle est plus fort.
      \eng{I use my \textbf{PHONE} to study.} (pas l'ordinateur) \quad \eng{She works in \textbf{DOUALA}.} (pas Yaoundé)
\end{enumerate}
\end{methode}

\begin{exoresolu}[Intonation et accentuation — Exercices corrigés]
Énoncé — (a) Marquez l'intonation finale (↗ ou ↘). (b) Soulignez la syllabe accentuée : \eng{computer, information, technology, develop, mobile, laptop}.

1. \eng{Do you often use the Internet for your homework?}

2. \eng{Which websites do you visit?}

3. \eng{ICTs can really enhance our learning.}
\tcblower
Solution détaillée :
\begin{itemize}
\item \textbf{1.} Question fermée (auxiliaire \eng{Do}) $\rightarrow$ \textbf{↗ montante}.
\item \textbf{2.} Question ouverte (\eng{Which}) $\rightarrow$ \textbf{↘ descendante}.
\item \textbf{3.} Phrase affirmative $\rightarrow$ \textbf{↘ descendante} (sur \eng{learning}).
\item \textbf{(b) Accentuation :} \eng{com\textbf{PU}ter}, \eng{infor\textbf{MA}tion}, \eng{tech\textbf{NO}logy}, \eng{de\textbf{VE}lop}, \eng{mo\textbf{BILE}}, \eng{\textbf{LAP}top}.
      \begin{itemize}
      \item Règle \eng{-tion / -sion} : accent sur syllabe précédente (\eng{infor\textbf{MA}tion}).
      \item Règle \eng{-logy / -graphy / -metry} : accent sur syllabe précédente (\eng{tech\textbf{NO}logy}, \eng{pho\textbf{TO}graphy}).
      \item Mots composés \eng{laptop, desktop, smartphone} : accent sur 1\textsuperscript{er} élément (\eng{\textbf{LAP}top}).
      \end{itemize}
\end{itemize}
\emph{Entraînement :} Lis chaque phrase à voix haute 3 fois en exagérant la flèche d'intonation. Enregistre-toi si possible.
\end{exoresolu}

\courssub{Guided role-plays and free practice}

\begin{experience}[Jeu de rôle guidé — Carte A / Carte B]
\textbf{Consigne :} En binôme, préparez et jouez le dialogue (8--10 répliques). Changez de rôle.
\begin{center}
\begin{tabularx}{\linewidth}{|X|X|}
\hline
\rowcolor{poporangeL} \textbf{ROLE A : Journalist (School Magazine / Local Radio)} & \textbf{ROLE B : Worker (Microfinance / Cybercafé / Shop / Office)} \\
\hline
\textbf{Context:} You interview a professional about ICTs at work. & \textbf{Context:} You answer the journalist's questions honestly. \\
\hline
\textbf{Must include:} & \textbf{Must include:} \\
\begin{itemize}
\item Opening formula (polite)
\item Question 1: \eng{What ICT tools\ldots?}
\item Question 2: \eng{How often\ldots?}
\item Question 3: \eng{Why\ldots?}
\item One follow-up: \eng{Could you tell me more\ldots?}
\item Closing formula
\end{itemize} &
\begin{itemize}
\item Full sentence answers
\item \eng{Because / So that / For example}
\item Specific tools (POS, spreadsheet, smartphone, database, mobile money)
\item Frequency (daily, twice a day, constantly)
\item Polite closing
\end{itemize} \\
\hline
\end{tabularx}
\end{center}
\emph{Aide-mémoire au tableau :} \eng{tools: computer, smartphone, POS, printer, router, scanner} | \eng{verbs: process, record, transfer, withdraw, deposit, manage, communicate} | \eng{connectors: because, so that, in order to, for instance, actually}.
\end{experience}

\begin{exoresolu}[Modèle de dialogue attendu pour le jeu de rôle]
Énoncé — Exemple de dialogue complet respectant les contraintes.
\tcblower
Solution détaillée — Dialogue type (10 répliques) :

\textbf{Journalist:} Good afternoon, Madam. I'm from Radio Buea. \textbf{May I ask you a few questions about ICTs in your work?}

\textbf{Worker:} Good afternoon. Yes, of course. Go ahead.

\textbf{Journalist:} \textbf{What ICT tools do you use daily at the microfinance agency?}

\textbf{Worker:} We use desktop computers for the client database and smartphones for mobile money transactions. We also have a POS terminal for card payments.

\textbf{Journalist:} \textbf{How often do you need an Internet connection?}

\textbf{Worker:} Constantly, all day long. Our system is cloud-based. If the connection fails, we cannot approve loans or process withdrawals.

\textbf{Journalist:} \textbf{Really? Could you tell me more about the mobile money service?}

\textbf{Worker:} Certainly. Many clients are farmers in villages around Buea. They deposit or withdraw cash via local agents. It saves them a long trip to town.

\textbf{Journalist:} \textbf{Why are ICTs so important for your agency?}

\textbf{Worker:} \textbf{Because} they make transactions instant and secure. \textbf{For example}, SMS receipts build trust. \textbf{In the past}, we wrote everything in ledgers; errors were frequent.

\textbf{Journalist:} \textbf{Thank you very much for your time and clear explanations.}

\textbf{Worker:} You are welcome. Good luck with your report. Goodbye!

\textbf{Journalist:} Goodbye!
\end{exoresolu}

\begin{exercice}[Production orale libre — Préparation au Bac]
Choisissez \textbf{un} sujet. Préparez 5 minutes (mots-clés, 3 questions, 2 réponses types). Présentez le dialogue devant la classe (binôme, 2--3 min).
\begin{enumerate}
\item \textbf{Sujet A (Études) :} Vous êtes délégué(e) de classe. Interviewez un camarade sur l'utilisation de la plateforme numérique de l'école (devoirs, notes, ressources) pour un article du journal scolaire.
\item \textbf{Sujet B (Travail) :} Vous êtes en stage dans une entreprise à Douala. Interviewez un employé du service comptabilité sur l'usage des tableurs (Excel) et de la messagerie pour gagner du temps.
\item \textbf{Sujet C (Vie quotidienne) :} Dans un cybercafé à Bamenda, vous discutez avec le gérant de l'évolution de son métier : du simple accès Internet aux services (impression, formation, mobile money, e-government).
\end{enumerate}
\emph{Critères d'évaluation (grille Bac) :} Pertinence du contenu (infos échangées) / Correction linguistique (questions, temps, prépositions) / Fluidité \& Prononciation (intention, accentuation) / Interaction (relance, politesse) / Richesse lexicale (vocabulaire TIC).
\end{exercice}

\begin{corrige}[Exercice — Pistes de réponses types (Sujet A)]
\textbf{Sujet A — Extrait de dialogue possible :}

\textbf{Delegate:} Hi Grace. \textbf{Do you have a minute?} I'm writing about the school platform.

\textbf{Grace:} Hi. Sure, what do you want to know?

\textbf{Delegate:} \textbf{How do you use the platform for your studies?}

\textbf{Grace:} \textbf{I log in every evening to check the assignments posted by teachers and to download the lesson notes in PDF.}

\textbf{Delegate:} \textbf{How often do you connect?}

\textbf{Grace:} \textbf{Almost every day}, except weekends sometimes. \textbf{It's very useful because} I can revise even when I'm at home in the village.

\textbf{Delegate:} \textbf{That's great. Could you tell me if there are any problems?}

\textbf{Grace:} \textbf{Well, sometimes the connection is slow} or the server is down. But generally it helps us a lot.

\textbf{Delegate:} \textbf{Thank you very much for your help, Grace.}

\textbf{Grace:} You're welcome. See you in class!
\end{corrige}

\courssub{Erreurs fréquentes à éviter (Focus Bac)}

\begin{attention}[Les 5 erreurs qui coûtent des points au Bac oral/écrit]
\begin{enumerate}
\item \textbf{Question sans inversion ni auxiliaire} : \eng{How you use your phone?} $\times$ $\rightarrow$ \eng{How \textbf{do} you use your phone?} \checkmark. L'auxiliaire \eng{do/does/did} est \textbf{obligatoire} aux présent/prétérit simples (sauf verbe \eng{be} ou modaux).
\item \textbf{Le \eng{-s} oublié ou en trop (3\textsuperscript{e} pers. sg.)} : \eng{He use\ldots} $\times$ ; \eng{Does he uses\ldots?} $\times$. $\rightarrow$ \eng{He \textbf{uses}\ldots} \checkmark ; \eng{\textbf{Does} he \textbf{use}\ldots?} \checkmark. Le \eng{-s} \textbf{migre} sur l'auxiliaire à l'interrogatif/négatif.
\item \textbf{Faux amis et calques du français} :
      \begin{itemize}
      \item \eng{to assist} = aider / assister (soigner) $\neq$ « assister à » $\rightarrow$ \eng{to attend} (a meeting, a class).
      \item \eng{actually} = en fait / en réalité $\neq$ « actuellement » $\rightarrow$ \eng{currently} / \eng{at the moment}.
      \item \eng{to utilize} = utiliser (formel, technique) $\rightarrow$ préférer \eng{to use} (courant).
      \item \eng{formation} = formation (géologique, militaire) $\neq$ « formation professionnelle » $\rightarrow$ \eng{training} / \eng{course}.
      \end{itemize}
\item \textbf{Prépositions fautives (collocations)} :
      \begin{itemize}
      \item \eng{listen \textbf{to}} (jamais \eng{listen} seul + COD) : \eng{I listen to podcasts.}
      \item \eng{connect \textbf{to}} the Internet / a network / a server.
      \item \eng{search \textbf{for}} information / a file (mais \eng{look for}).
      \item \eng{depend \textbf{on}} (pas \eng{depend of}).
      \item \eng{on the phone} / \eng{on my laptop} / \eng{on the platform} (pas \eng{in}).
      \item \eng{at work} / \eng{at school} / \eng{at university} (lieu d'activité).
      \end{itemize}
\item \textbf{Noms indénombrables \& ordre des mots} :
      \begin{itemize}
      \item \eng{information, homework, equipment, software, advice, vocabulary} : \textbf{pas de \eng{-s}} ($\times$ \eng{informations} $\rightarrow$ \eng{pieces of information}).
      \item Adverbe de fréquence \textbf{avant} verbe principal (sauf \eng{be}) : \eng{I \textbf{often} use} / \eng{She \textbf{always} checks} / \eng{It \textbf{is} usually fast}.
      \item \eng{Goodbye} (un mot, majuscule G) ; pronom \eng{I} \textbf{toujours majuscule}.
      \end{itemize}
\end{enumerate}
\emph{Astuce révision :} Relis tes dialogues en cochant chaque règle ci-dessus. Corrige-toi avant l'examinateur !
\end{attention}

\begin{savaistu}[Culture \& Contexte : Le Cameroun et le numérique]
\begin{itemize}
\item \textbf{Silicon Mountain (Buea) :} Surnom de l'écosystème tech anglophone (start-ups, incubateurs comme ActivSpaces). Des jeunes développent des apps pour l'agriculture, la santé, la fintech.
\item \textbf{GCE Board \& Digitalisation :} Depuis 2022, l'inscription au GCE (O/L, A/L) et la consultation des résultats se font \textbf{en ligne} via le portail officiel. Les élèves doivent maîtriser \eng{login, password, upload scanned documents, print slip}.
\item \textbf{Mobile Money (MoMo, Orange Money) :} Véritable révolution bancaire. Permet aux non-bancarisés (marchands, agriculteurs, élèves) d'épargner, payer l'école, recevoir des transferts. Compétence numérique essentielle.
\item \textbf{E-government :} Portail \eng{gov.cm}, recrutement en ligne (CONCOURS), déclaration d'impôts, e-visa. L'administration exige de plus en plus des \textbf{digital skills}.
\item \textbf{Connectivité :} Projet \eng{PAVIA} (backbone fibre optique), 4G/5G (MTN, Orange, Camtel, Nexttel). Le coût des data baisse mais reste un frein en zone rurale.
\item \textbf{Compétence 21\textsuperscript{e} siècle :} \eng{Digital literacy} (savoir chercher, vérifier, communiquer, créer en ligne) est désormais au programme officiel MINESEC au même titre que lire/écrire/calculer.
\end{itemize}
\emph{Pour l'oral :} Citer un exemple local (Silicon Mountain, GCE online, Mobile Money) valorise ton propos et montre ton ancrage camerounais.
\end{savaistu}

\newpage

\coursec{Exercices}

\courssub{Je vérifie mes connaissances}

\begin{exercice}[QCM : les TIC et la communication]
Choisis la bonne réponse (a, b ou c) pour chaque question.
\begin{enumerate}
\item ICTs stands for \ldots
\begin{enumerate}
\item International Communication Tools
\item Information and Communication Technologies
\item Internet and Computer Training
\end{enumerate}
\item To ask for more details politely, you say :
\begin{enumerate}
\item \eng{Tell me more now!}
\item \eng{Could you tell me more about it?}
\item \eng{What more?}
\end{enumerate}
\item \eng{I use my tablet \ldots\ watch online lessons.}
\begin{enumerate}
\item for
\item to
\item at
\end{enumerate}
\item To close a conversation politely, you can say :
\begin{enumerate}
\item \eng{Thank you for sharing this information.}
\item \eng{Stop talking now.}
\item \eng{I am going.}
\end{enumerate}
\end{enumerate}
\end{exercice}

\begin{exercice}[Vrai ou faux ?]
Écoute ton professeur lire le court dialogue ci-dessous (ou lis-le avec un camarade), puis indique si les affirmations sont \textbf{vraies (True)} ou \textbf{fausses (False)}. Justifie chaque réponse par une courte phrase en anglais.
\begin{textebox}[A short dialogue]
\eng{— Hi Sandra! What do you use your smartphone for?}\\
\eng{— Well, I mostly use WhatsApp to chat with my classmates, and I send my homework to my teachers by email.}\\
\eng{— Do you ever use it for learning?}\\
\eng{— Yes, I watch English lessons on YouTube every evening. It really helps me improve my pronunciation.}
\end{textebox}
\begin{enumerate}
\item \eng{The girl uses WhatsApp to send emails.}
\item \eng{Sandra sends her homework by email.}
\item \eng{She never uses her phone for learning.}
\item \eng{Watching videos helps her pronunciation.}
\end{enumerate}
\end{exercice}

\begin{exercice}[Vocabulaire : le monde des TIC]
Complète chaque phrase avec le mot anglais qui convient, choisi dans la liste : \eng{laptop — password — to download — website — network — to browse}.
\begin{enumerate}
\item I forgot my \ldots, so I cannot open my email account.
\item Students \ldots\ the Internet to find information for their projects.
\item You can \ldots\ free learning applications on your phone.
\item The school \ldots\ is very slow today; nobody can connect.
\item My father bought a new \ldots\ to work from home.
\item Our school has a \ldots\ where exam results are published.
\end{enumerate}
\end{exercice}

\begin{exercice}[Grammaire : questions et prépositions]
Complète les questions suivantes avec la forme correcte de l'auxiliaire et la préposition finale si nécessaire.
\begin{enumerate}
\item What \ldots\ you \ldots\ (use) your phone \ldots ?
\item How often \ldots\ she \ldots\ (check) her emails ?
\item Which applications \ldots\ they \ldots\ (download) last week ?
\item \ldots\ you ever \ldots\ (attend) an online lesson ?
\item Who \ldots\ (teach) you how to use this software ?
\end{enumerate}
\end{exercice}

\courssub{Je m'entraîne}

\begin{exercice}[Traduction : du français à l'anglais]
Traduis les phrases suivantes en anglais correct. Attention à l'ordre des mots et à la place de la préposition dans les questions.
\begin{enumerate}
\item Pour quoi utilises-tu ton téléphone ?
\item J'utilise Internet pour faire mes recherches et réviser mes leçons.
\item Pourrais-tu m'en dire plus sur cette application ?
\item Merci beaucoup d'avoir partagé ces informations avec moi.
\item Les enseignants utilisent les TIC pour préparer leurs cours.
\end{enumerate}
\end{exercice}

\begin{exercice}[Dialogue à trous : les TIC pour apprendre]
Complète le dialogue suivant avec les formules fonctionnelles de la liste : \eng{Could you tell me more? — Thank you for sharing. — What do you use it for? — How often do you use it? — I use it to improve my vocabulary.}
\begin{textebox}[Complete the dialogue]
\eng{— Hello Kevin! I see you have a new tablet. (1) \ldots}\\
\eng{— Well, (2) \ldots\ I learn ten new English words every day.}\\
\eng{— That sounds interesting! (3) \ldots}\\
\eng{— Every evening, after my homework.}\\
\eng{— (4) \ldots\ Which application do you use?}\\
\eng{— It is a free dictionary app with games and quizzes.}\\
\eng{— Great! (5) \ldots\ I will try it too.}
\end{textebox}
\end{exercice}

\begin{exercice}[Remets le dialogue dans l'ordre]
Les répliques de ce dialogue sur les TIC au travail sont mélangées. Numérote-les de 1 à 6 pour reconstituer une conversation logique.
\begin{enumerate}
\item[\ldots] \eng{Sure! We use email to communicate with customers and video calls for meetings.}
\item[\ldots] \eng{Good morning, Mrs Etoundi! Could you tell me how your company uses ICTs?}
\item[\ldots] \eng{You are welcome. Good luck with your presentation!}
\item[\ldots] \eng{That is very useful. And what about training the workers?}
\item[\ldots] \eng{We organise online courses twice a month so that everyone can learn new skills.}
\item[\ldots] \eng{Thank you very much for this information, Mrs Etoundi.}
\end{enumerate}
\end{exercice}

\begin{exercice}[Intonation : montante ou descendante ?]
Indique si l'intonation de chaque question est \textbf{montante (rising)} ou \textbf{descendante (falling)}, puis lis chaque question à voix haute en respectant la mélodie correcte.
\begin{enumerate}
\item \eng{Do you have a computer at home?}
\item \eng{What do you use your phone for?}
\item \eng{Can you help me download this application?}
\item \eng{How often do you check your emails?}
\item \eng{Did you watch the online lesson yesterday?}
\item \eng{Why do students prefer digital books?}
\end{enumerate}
\end{exercice}

\courssub{Je me prépare au Bac}

\begin{exercice}[Compréhension de l'écrit et langue]
Lis attentivement le texte suivant, puis réponds aux questions en anglais, par des phrases complètes.
\begin{textebox}[ICTs change learning in Cameroonian schools]
\eng{In many secondary schools in Yaoundé and Douala, students are discovering new ways of learning thanks to Information and Communication Technologies. In the past, a student who missed a lesson had to copy a classmate's notes. Today, teachers send lesson summaries through WhatsApp groups, and some schools have created websites where pupils can download past examination papers.}\\
\eng{Marlyse, a Terminale student, explains: “I use my phone to watch science experiments on video. When I do not understand a word, I check an online dictionary immediately. It is faster than asking somebody.” However, not everything is perfect. The Internet connection is often slow, and many families cannot afford a computer. Teachers also warn students against spending too much time on social media instead of studying.}\\
\eng{Despite these difficulties, most students agree that ICTs make learning easier, faster and more enjoyable.}
\end{textebox}
\textbf{Comprehension questions}
\begin{enumerate}
\item How did students catch up with missed lessons in the past ?
\item Mention two ways schools use ICTs according to the text.
\item What does Marlyse do when she does not understand a word ?
\item Give two problems connected with the use of ICTs mentioned in the text.
\item Do you think ICTs are useful for your own studies ? Justify your answer in two or three sentences.
\end{enumerate}
\textbf{Language work}
\begin{enumerate}
\item Find in the text : (a) a phrasal verb meaning “to verify”, (b) an adjective meaning “more pleasant”.
\item Put into the interrogative form : \eng{Teachers send lesson summaries through WhatsApp groups.}
\item Put into the negative form : \eng{Many families can afford a computer.}
\end{enumerate}
\end{exercice}

\begin{exercice}[Traduction]
Traduis en anglais le passage suivant (environ 70 mots). Soigne les temps verbaux, les prépositions et l'ordre des mots.
\begin{textebox}[Texte à traduire]
Dans mon école, beaucoup d'élèves utilisent leur téléphone portable pour apprendre l'anglais. Ils regardent des vidéos, écoutent des chansons et téléchargent des applications gratuites. Notre professeur nous encourage à utiliser Internet pour faire des recherches, mais il nous demande de ne pas passer trop de temps sur les réseaux sociaux. Grâce aux TIC, je révise mes leçons tous les soirs et mes notes se sont améliorées.
\end{textebox}
\end{exercice}

\begin{exercice}[Jeu de rôle et composition — type Bac]
\textbf{Part A — Guided role-play (oral practice).}\\
Situation : \eng{You meet a friend who wants to improve his/her English with ICTs. Ask at least four questions about his/her habits, give at least three pieces of advice, react to his/her answers, and close the conversation politely.}\\
Prépare le dialogue par écrit (12 à 15 répliques au total), puis joue-le avec un camarade devant la classe. Tu seras évalué(e) selon la grille suivante :
\begin{center}
\begin{tabularx}{\linewidth}{|X|l|}
\hline
\rowcolor{popblueL} \textbf{Critère} & \textbf{Points} \\
\hline
Interaction : questions, relances, clôture polie & /4 \\
\hline
Vocabulaire des TIC et formules fonctionnelles & /4 \\
\hline
Prononciation, accentuation et intonation & /4 \\
\hline
Correction grammaticale (questions, prépositions) & /4 \\
\hline
Aisance et participation des deux partenaires & /4 \\
\hline
\end{tabularx}
\end{center}
\textbf{Part B — Composition (writing).}\\
Sujet : \eng{Your school wants to buy computers and connect classrooms to the Internet. Write an article for the school magazine explaining how ICTs can improve students' learning, mentioning two advantages and one possible danger, and giving your personal opinion.}\\
Rédige un article de \textbf{180 à 200 mots}, avec un titre, une introduction, deux paragraphes de développement et une conclusion.
\end{exercice}

\begin{methode}[Réussir le jeu de rôle (Part A)]
\begin{enumerate}
\item \textbf{Ouvrir} poliment : \eng{Hi Paul! I heard you want to improve your English. Could I ask you a few questions?}
\item \textbf{Poser des questions} variées (habitudes, fréquence, outils) : \eng{What do you use your phone for? How often do you watch videos in English? Have you ever tried a learning application?}
\item \textbf{Réagir} aux réponses : \eng{That sounds interesting! Really? That is very useful!}
\item \textbf{Donner des conseils} : \eng{You should download a dictionary app. Why don't you watch English lessons on YouTube? You could join an online study group.}
\item \textbf{Conclure} poliment : \eng{Thank you for sharing. Good luck with your English!}
\end{enumerate}
\end{methode}

\begin{methode}[Rédiger l'article (Part B)]
\begin{enumerate}
\item \textbf{Titre} court et accrocheur : \eng{Computers in our classrooms: a chance for every student}.
\item \textbf{Introduction} : annonce le projet de l'école et ta position générale (2--3 phrases).
\item \textbf{Développement 1} : deux avantages (accès à l'information, leçons en vidéo, téléchargement d'anciens sujets, motivation). Utilise \eng{first of all}, \eng{moreover}.
\item \textbf{Développement 2} : un danger possible (distraction des réseaux sociaux, coût, connexion lente) avec \eng{however}, \eng{on the other hand}.
\item \textbf{Conclusion} : ton opinion personnelle avec \eng{in my opinion}, \eng{to conclude}.
\item \textbf{Relecture} : vérifie les temps (présent de vérité générale), les prépositions (\eng{to connect to}, \eng{to depend on}, \eng{to spend time on}) et le nombre de mots.
\end{enumerate}
\end{methode}

\begin{attention}[Erreurs fréquentes des francophones]
\begin{itemize}
\item \emph{Informations} est un faux ami : en anglais, \eng{information} est \textbf{indénombrable}. On dit \eng{this information}, \eng{some information}, jamais \emph{informations}.
\item Ne dis pas \emph{to make research} : dis \eng{to do research} (« faire des recherches »).
\item Dans les questions avec préposition, celle-ci va en fin de phrase : \eng{What do you use your phone for?} et non \emph{For what do you use your phone?} (trop soutenu).
\item Après \eng{to use something}, le but s'exprime avec \emph{to} + base verbale : \eng{I use my tablet to learn English}, et non \emph{for learn}.
\item \eng{Social media} se construit avec \emph{on} : \eng{to spend time on social media}.
\item Prononciation : \eng{technology} se dit « tèk-no-lo-dji » (accent sur la 2\up{e} syllabe) ; \eng{download} se dit « da-oun-lode » (accent sur la 1\up{re} syllabe pour le verbe).
\end{itemize}
\end{attention}

\newpage

\coursec{Corrigés des exercices}

\begin{corrige}[Exercice 1 — QCM : les TIC et la communication]
\begin{enumerate}
\item \textbf{b} — \eng{ICTs} signifie \eng{Information and Communication Technologies} (« technologies de l'information et de la communication »). Les propositions a et c sont des inventions plausibles mais fausses.
\item \textbf{b} — \eng{Could you tell me more about it?} est la formule polie de relance : \emph{Could you} + base verbale exprime la politesse. La réponse a (\eng{Tell me more now!}) est un ordre brutal ; la réponse c (\eng{What more?}) n'est pas anglaise.
\item \textbf{b} — \eng{I use my tablet to watch online lessons.} : le but s'exprime avec \emph{to} + base verbale (« pour regarder »). \emph{For} serait suivi d'un nom (\eng{for watching} est incorrect ici).
\item \textbf{a} — \eng{Thank you for sharing this information.} est la clôture polie attendue ; b et c sont impolies et mettraient fin à l'échange de façon discourtoise.
\end{enumerate}
\end{corrige}

\begin{corrige}[Exercice 2 — Vrai ou faux ?]
\begin{enumerate}
\item \textbf{False} — \eng{She uses WhatsApp to chat with her classmates, not to send emails.} (Le dialogue confond délibérément les deux usages : c'est un piège classique de compréhension.)
\item \textbf{True} — \eng{She says: “I send my homework to my teachers by email.”}
\item \textbf{False} — \eng{She watches English lessons on YouTube every evening, so she does use her phone for learning.}
\item \textbf{True} — \eng{She says it really helps her improve her pronunciation.}
\end{enumerate}
\textbf{Point de méthode} : pour justifier, cite toujours une phrase du texte support, même reformulée ; une réponse \emph{True/False} sans justification ne vaut que la moitié des points.
\end{corrige}

\begin{corrige}[Exercice 3 — Vocabulaire : le monde des TIC]
\begin{enumerate}
\item \eng{password} — \eng{I forgot my password, so I cannot open my email account.} (« mot de passe » ; se prononce « passe-oueurd »).
\item \eng{browse} — \eng{Students browse the Internet to find information for their projects.} (\eng{to browse the Internet} = « naviguer sur Internet » ; verbe conjugué au présent simple, sans \emph{-s} car le sujet est pluriel).
\item \eng{download} — \eng{You can download free learning applications on your phone.} (\emph{can} + base verbale ; « télécharger »).
\item \eng{network} — \eng{The school network is very slow today; nobody can connect.} (« réseau »).
\item \eng{laptop} — \eng{My father bought a new laptop to work from home.} (« ordinateur portable » ; se prononce « la-ptop »).
\item \eng{website} — \eng{Our school has a website where exam results are published.} (« site web »).
\end{enumerate}
\end{corrige}

\begin{corrige}[Exercice 4 — Grammaire : questions et prépositions]
\begin{enumerate}
\item \eng{What do you use your phone for?} — présent simple : auxiliaire \emph{do} + base verbale ; la préposition \emph{for} se place en fin de question (« tu utilises ton téléphone pour quoi ? »).
\item \eng{How often does she check her emails?} — à la 3\up{e} personne du singulier, l'auxiliaire est \emph{does} et le verbe principal revient à la base verbale : on dit \emph{does she check}, jamais \emph{does she checks}.
\item \eng{Which applications did they download last week?} — le marqueur \emph{last week} impose le prétérit : auxiliaire \emph{did} + base verbale.
\item \eng{Have you ever attended an online lesson?} — \emph{ever} (« déjà, jamais » dans une question d'expérience) appelle le present perfect : \emph{have} + participe passé (\emph{attended}, verbe régulier).
\item \eng{Who taught you how to use this software?} — quand \emph{who} est sujet du verbe, il n'y a \textbf{pas d'auxiliaire} ; le verbe se conjugue directement au prétérit. \emph{Teach} est irrégulier : teach -- taught -- taught.
\end{enumerate}
\textbf{Point de vigilance} : \eng{software} est indénombrable en anglais (jamais \emph{softwares}).
\end{corrige}

\begin{corrige}[Exercice 5 — Traduction : du français à l'anglais]
\begin{enumerate}
\item \eng{What do you use your phone for?} — la préposition \emph{for} se place en fin de question dans la langue courante ; la construction \emph{For what\ldots} est possible mais très soutenue.
\item \eng{I use the Internet to do my research and revise my lessons.} — \emph{to} + base verbale exprime le but ; attention, on dit \eng{to do research} et non \emph{to make research}.
\item \eng{Could you tell me more about this application?} — \emph{Could you} + base verbale pour la demande polie ; \emph{more about} = « davantage au sujet de ».
\item \eng{Thank you very much for sharing this information with me.} — \emph{thank you for} + verbe en \emph{-ing} ; \eng{information} est indénombrable : jamais de \emph{s}.
\item \eng{Teachers use ICTs to prepare their lessons.} — pas d'article devant un pluriel de sens général (\emph{teachers} = « les enseignants en général »).
\end{enumerate}
\end{corrige}

\begin{corrige}[Exercice 6 — Dialogue à trous : les TIC pour apprendre]
\begin{enumerate}
\item \eng{What do you use it for?} — question d'ouverture logique après l'annonce de la nouvelle tablette ; \emph{it} reprend \emph{tablet}.
\item \eng{I use it to improve my vocabulary.} — seule réponse de la liste qui annonce l'apprentissage de dix mots par jour.
\item \eng{How often do you use it?} — la réponse \eng{Every evening} indique une fréquence, donc la question portait sur la fréquence.
\item \eng{Could you tell me more?} — relance polie qui introduit naturellement la question suivante sur l'application.
\item \eng{Thank you for sharing.} — formule de clôture ; la phrase \eng{I will try it too} confirme que l'échange se termine positivement.
\end{enumerate}
\textbf{Point de méthode} : dans un dialogue à trous, lis toujours la réplique \emph{qui suit} le trou : elle révèle la nature de la formule attendue (question de fréquence, remerciement, relance\ldots).
\end{corrige}

\begin{corrige}[Exercice 7 — Remets le dialogue dans l'ordre]
Ordre correct : \textbf{2 -- 1 -- 6 -- 3 -- 4 -- 5}. Dialogue reconstitué :
\begin{enumerate}
\item \eng{Good morning, Mrs Etoundi! Could you tell me how your company uses ICTs?} — salutation + demande d'information : c'est l'ouverture.
\item \eng{Sure! We use email to communicate with customers and video calls for meetings.} — \eng{Sure!} répond directement à la demande polie.
\item \eng{That is very useful. And what about training the workers?} — réaction (\eng{That is very useful}) + relance avec \eng{what about\ldots?}
\item \eng{We organise online courses twice a month so that everyone can learn new skills.} — réponse à la question sur la formation (\emph{training}).
\item \eng{Thank you very much for this information, Mrs Etoundi.} — remerciement : la conversation touche à sa fin.
\item \eng{You are welcome. Good luck with your presentation!} — réponse au remerciement + formule de clôture.
\end{enumerate}
\textbf{Indices de logique} : la salutation ouvre toujours ; \eng{Sure!} ne peut que répondre à une demande ; \eng{You are welcome} ne peut que suivre un remerciement.
\end{corrige}

\begin{corrige}[Exercice 8 — Intonation : montante ou descendante ?]
\textbf{Règle} : les questions fermées (réponse \emph{yes/no}, introduites par un auxiliaire : \emph{do, does, can, did, have\ldots}) se prononcent avec une intonation \textbf{montante} ; les questions ouvertes en \emph{wh-} (\emph{what, how, why, where, when\ldots}) ont une intonation \textbf{descendante}.
\begin{enumerate}
\item \eng{Do you have a computer at home?} — \textbf{montante} (question \emph{yes/no}).
\item \eng{What do you use your phone for?} — \textbf{descendante} (question en \emph{what}).
\item \eng{Can you help me download this application?} — \textbf{montante} (question \emph{yes/no}).
\item \eng{How often do you check your emails?} — \textbf{descendante} (question en \emph{how often}).
\item \eng{Did you watch the online lesson yesterday?} — \textbf{montante} (question \emph{yes/no}).
\item \eng{Why do students prefer digital books?} — \textbf{descendante} (question en \emph{why}).
\end{enumerate}
\textbf{Entraînement oral} : sur les questions montantes, la voix s'élève sur le dernier mot accentué (\emph{home}, \emph{application}, \emph{yesterday}) ; sur les questions descendantes, elle tombe sur le dernier mot (\emph{for}, \emph{emails}, \emph{books}).
\end{corrige}

\begin{corrige}[Exercice 9 — Compréhension de l'écrit et langue]
\textbf{Comprehension questions}
\begin{enumerate}
\item \eng{In the past, a student who missed a lesson had to copy a classmate's notes.}
\item \eng{Teachers send lesson summaries through WhatsApp groups, and some schools have created websites where pupils can download past examination papers.}
\item \eng{When she does not understand a word, she checks an online dictionary immediately.}
\item \eng{The Internet connection is often slow, and many families cannot afford a computer.} (On accepte également : \eng{students may spend too much time on social media instead of studying}.)
\item Réponse personnelle. Modèle : \eng{Yes, I think ICTs are very useful for my studies because I can watch lessons again and check difficult words in online dictionaries. They make revision faster and more enjoyable.}
\end{enumerate}
\textbf{Language work}
\begin{enumerate}
\item (a) \eng{to catch up with} = « rattraper, se mettre à jour » (\eng{students catch up with missed lessons}) ; (b) \eng{more enjoyable} = « plus agréable » (comparatif de l'adjectif long \emph{enjoyable}).
\item \eng{Do teachers send lesson summaries through WhatsApp groups?} — présent simple, sujet pluriel : auxiliaire \emph{do} en tête, le verbe reste à la base verbale.
\item \eng{Many families cannot afford a computer.} — la négation de \emph{can} est \emph{cannot} (en un seul mot à l'écrit) ou \emph{can't}.
\end{enumerate}
\textbf{Barème indicatif (type Bac)} : compréhension 5 × 2 pts = /10 ; language work 3 questions = /6 ; qualité de la langue = /4. Total /20.
\end{corrige}

\begin{corrige}[Exercice 10 — Traduction]
\textbf{Traduction de référence} :

\eng{In my school, many students use their mobile phones to learn English. They watch videos, listen to songs and download free applications. Our teacher encourages us to use the Internet to do research, but he asks us not to spend too much time on social media. Thanks to ICTs, I revise my lessons every evening, and my marks have improved.}

\textbf{Points de vigilance} :
\begin{itemize}
\item « encourager quelqu'un à faire » = \eng{to encourage somebody to do} ; « demander à quelqu'un de ne pas faire » = \eng{to ask somebody not to do} (la négation porte sur \emph{not to spend}).
\item « passer du temps sur » = \eng{to spend time on} ; « les réseaux sociaux » = \eng{social media} (pluriel, sans article ici).
\item « mes notes se sont améliorées » : résultat présent d'une évolution passée $\Rightarrow$ present perfect \eng{have improved}. \eng{Marks} (britannique) ou \eng{grades} (américain) = « notes ».
\item « grâce à » = \eng{thanks to} + nom ; « tous les soirs » = \eng{every evening} (sans préposition).
\item \eng{to do research} et non \emph{to make research} ; \eng{free applications} = « applications gratuites ».
\end{itemize}
\textbf{Barème indicatif} : respect du sens /8 ; temps verbaux /4 ; prépositions et constructions /4 ; orthographe et vocabulaire /4. Total /20.
\end{corrige}

\begin{corrige}[Exercice 11 — Jeu de rôle et composition — type Bac]
\textbf{Part A — Modèle de dialogue (12 à 15 répliques)} :

\eng{— Hi Marlyse! I heard you want to improve your English with ICTs. Could I ask you a few questions?}\\
\eng{— Of course, go ahead!}\\
\eng{— What do you use your phone for?}\\
\eng{— Mostly for WhatsApp and music, to be honest.}\\
\eng{— I see. How often do you watch videos in English?}\\
\eng{— Almost never, actually. I find them difficult.}\\
\eng{— Really? You should start with short videos with subtitles. Have you ever tried a learning application?}\\
\eng{— No, I haven't. Which one do you recommend?}\\
\eng{— You could download a free dictionary app with games. Why don't you also join our online study group?}\\
\eng{— That sounds interesting! How often does the group meet?}\\
\eng{— Twice a week, in the evening. It really helped me improve my pronunciation.}\\
\eng{— Great! I will try both. Thank you very much for your advice.}\\
\eng{— You are welcome. Good luck with your English!}

\textbf{Barème indicatif (Part A, /20)} : interaction (questions, relances, clôture) /4 ; vocabulaire des TIC et formules fonctionnelles /4 ; prononciation, accentuation, intonation /4 ; correction grammaticale /4 ; aisance et équilibre des prises de parole /4.

\textbf{Part B — Proposition de rédaction modèle (environ 190 mots)} :

\eng{\textbf{Computers in our classrooms: a chance for every student}}\\
\eng{Our school is planning to buy computers and connect every classroom to the Internet. This is excellent news, because ICTs can deeply change the way we learn.}\\
\eng{First of all, connected classrooms will give every student access to an enormous amount of information. Instead of depending only on one textbook, we shall be able to do research, watch science experiments on video and download past examination papers to prepare for our exams. Moreover, learning with videos, songs and educational games is more enjoyable, so students are usually more motivated.}\\
\eng{However, there is one possible danger: distraction. If students spend their time on social media instead of studying, computers will not help them at all. Teachers will therefore have to guide us and teach us to use these tools responsibly.}\\
\eng{In my opinion, the advantages are far greater than the risks. Thanks to ICTs, learning becomes easier, faster and more attractive. To conclude, I warmly welcome this project, and I hope every classroom will soon be connected.}

\textbf{Points de vigilance (Part B)} : titre obligatoire ; connecteurs variés (\eng{first of all}, \eng{moreover}, \eng{however}, \eng{in my opinion}, \eng{to conclude}) ; présent de vérité générale ; prépositions correctes (\eng{access to}, \eng{to depend on}, \eng{to spend time on}) ; \eng{information} indénombrable ; respecter la fourchette de 180 à 200 mots.

\textbf{Barème indicatif (Part B, /20)} : respect du sujet et du format (article avec titre) /4 ; organisation en paragraphes et connecteurs /4 ; richesse du vocabulaire des TIC /4 ; correction grammaticale /4 ; pertinence des arguments et opinion personnelle /4.
\end{corrige}

\newpage

\coursec{Fiche bilan}

\begin{popfigure}
\begin{tikzpicture}[mindmap, grow cyclic, every node/.style={concept, text=white, font=\small},
  root concept/.append style={concept color=popblue, font=\bfseries},
  level 1/.append style={level distance=3.4cm, sibling angle=60},
  level 1 concept/.append style={font=\small}]
\node[root concept]{ICTs for work and learning}
  child[concept color=poporange]{ node[align=center]{Key vocabulary\\ devices, apps, tools} }
  child[concept color=popgreen]{ node[align=center]{Asking\\ questions} }
  child[concept color=poppurple]{ node[align=center]{Giving\\ information} }
  child[concept color=poppink]{ node[align=center]{Follow-up\\ and closing} }
  child[concept color=popgold]{ node[align=center]{Present tense\\ and modals} }
  child[concept color=popmuted]{ node[align=center]{Stress and\\ intonation} };
\end{tikzpicture}
\legende{Carte mentale de la leçon : échanger des informations sur les TIC}
\end{popfigure}

\begin{bilan}
\textbf{1. Lexique clé (à connaître par cœur)}
\begin{itemize}
  \item \eng{ICT (Information and Communication Technology)} : TIC, technologies de l'information
  \item \eng{a device} : un appareil ; \eng{a smartphone / a laptop / a tablet} : un smartphone / un ordinateur portable / une tablette
  \item \eng{an app (application)} : une application ; \eng{software} : un logiciel ; \eng{a search engine} : un moteur de recherche
  \item \eng{to browse the Internet} : naviguer sur Internet ; \eng{to download / to upload} : télécharger (vers l'aval / vers l'amont)
  \item \eng{to enhance / to improve} : améliorer, renforcer ; \eng{online learning} : l'apprentissage en ligne
  \item \eng{a tutorial} : un tutoriel ; \eng{to save time} : gagner du temps ; \eng{reliable} : fiable
\end{itemize}

\textbf{2. Formules fonctionnelles}
\begin{center}
\begin{tabularx}{\linewidth}{|l|X|}\hline
\rowcolor{popblueL} Fonction & Formules anglaises \\ \hline
Demander & \eng{How do you use your phone for studies?} ; \eng{Could you tell me which app you use?} \\ \hline
Donner une information & \eng{I use WhatsApp to share notes.} ; \eng{Actually, Google helps me do research.} \\ \hline
Relancer & \eng{Really? How often?} ; \eng{What do you mean exactly?} ; \eng{Can you give an example?} \\ \hline
Conclure & \eng{So, ICTs really save time.} ; \eng{Thanks for sharing. That's all for now.} \\ \hline
\end{tabularx}
\end{center}

\textbf{3. Grammaire à connaître par cœur}
\begin{itemize}
  \item \cle{Present simple} pour les habitudes : \eng{I use my laptop every day.} (« J'utilise mon ordinateur tous les jours. »)
  \item \cle{Present continuous} pour une action en cours : \eng{I am downloading a tutorial now.}
  \item Questions avec \eng{do/does} et mots interrogatifs : \eng{What do you use it for?} ; \eng{How often does she connect?}
  \item \cle{Modaux} : \eng{can} (capacité), \eng{could you} (politesse), \eng{should} (conseil) : \eng{You should try this app.}
  \item But avec \eng{to + base verbale} : \eng{I use the Internet to do research.}
\end{itemize}

\textbf{4. Prononciation et intonation}
\begin{itemize}
  \item \eng{technology} se prononce « tèk-no-lo-dji » (accent sur la 2\up{e} syllabe).
  \item Questions en \eng{wh-} : intonation descendante ; questions fermées (\eng{Do you...?}) : intonation montante.
  \item Accentuer les mots porteurs de sens : \eng{I USE my PHONE to LEARN English.}
\end{itemize}

\textbf{5. Méthode express pour un dialogue}
\begin{enumerate}
  \item Saluer et introduire le sujet (\eng{Hi! Can I ask you something about ICTs?}).
  \item Poser des questions ouvertes et écouter la réponse.
  \item Relancer avec \eng{Really? Why? How?}
  \item Donner sa propre expérience.
  \item Conclure poliment.
\end{enumerate}
\end{bilan}

\begin{aretenir}[Checklist avant le Bac]
\begin{itemize}
  \item $\square$ Je connais au moins 15 mots de vocabulaire sur les TIC (devices, apps, actions).
  \item $\square$ Je sais saluer et introduire le sujet du dialogue en anglais.
  \item $\square$ Je sais poser des questions avec \eng{do/does} et les mots interrogatifs (\eng{what, how, how often, why}).
  \item $\square$ Je sais donner une information au present simple avec \eng{use ... to + verbe}.
  \item $\square$ Je sais relancer mon partenaire (\eng{Really? What do you mean? Can you give an example?}).
  \item $\square$ Je sais exprimer le conseil avec \eng{should} et la capacité avec \eng{can}.
  \item $\square$ Je sais conclure un dialogue poliment (\eng{Thanks for sharing.}).
  \item $\square$ Je fais attention à l'intonation : descendante pour les questions en \eng{wh-}, montante pour les questions fermées.
  \item $\square$ Je prononce correctement \eng{technology} (« tèk-no-lo-dji ») et \eng{Internet} (« in-tèr-nèt »).
  \item $\square$ Je peux tenir un dialogue de 6 à 8 répliques sans lire mes notes.
  \item $\square$ J'évite les faux amis : \eng{actually} = « en fait » (pas « actuellement »).
  \item $\square$ Je n'oublie pas le \eng{-s} à la 3\up{e} personne : \eng{She uses}, pas \eng{she use}.
\end{itemize}
\end{aretenir}

\findecours

\end{document}
